\documentclass[10pt,a4paper]{article}

% ----------------- Paquetes -------------------%
\usepackage[T1]{fontenc}
\usepackage[utf8]{inputenc}
\usepackage[a4paper,margin=2.5cm]{geometry}
\usepackage{mathptmx}             
\usepackage{changepage} 
\usepackage{setspace}
\usepackage[spanish]{babel}
\usepackage[labelfont=bf,labelsep=space]{caption}
\usepackage{amssymb}
\usepackage{amsmath}
\usepackage{ebgaramond}
\usepackage{placeins}
\usepackage{etoolbox}
\usepackage{setspace}
\usepackage{parskip} 
\usepackage{tcolorbox}
\usepackage{needspace}
\usepackage{xparse}
\usepackage{ocgx2}
\usepackage{hyperref}
\usepackage{hypcap}
\usepackage{soul}\sethlcolor{yellow}% resaltado amarillo: \hl{texto} (Panel TCEE, ⌃H)


%%%%%%%%%%%%%%%%%%%%%%%%%%%%%%%%%%%%%%%%%%%%%%%%%%%%%%%%%%%%%%%%
% --- Fija primero tus valores globales "buenos" ---
\setstretch{1.2}                 % Interlineado global
\setlength{\parskip}{0.7\baselineskip}
\setlength{\parindent}{0pt}
\newlength{\GlobalParskip}
\setlength{\GlobalParskip}{\parskip}

\newcounter{eqnum}[subsection]
\renewcommand{\theeqnum}{\arabic{eqnum}}
\newcounter{alphaitem}
\newcounter{numitem}

%% ------------------- Recuadros----------------------%%
\newcommand{\puntosclave}[1]{%
  \begin{tcolorbox}[
    colframe=black,
    title={Puntos clave},
    before upper={\setlength{\parindent}{0.5cm}\setlength{\parskip}{0.5\baselineskip}}
  ]
  #1
  \end{tcolorbox}
}

\newcommand{\modificaciones}[1]{
  \begin{tcolorbox}[
    colback=white,
    colframe=red,
    title={Modificaciones a realizar},
    before upper={\setlength{\parindent}{0.5cm}\setlength{\parskip}{0.5\baselineskip}}
  ]
  #1
  \end{tcolorbox}
}

%% ------------------- CITA ----------------------%%
% \begin{cita}[Autor][Año][Obra] texto \end{cita}: cita sangrada y, debajo a la derecha, «AUTOR (Año), Obra». Los tres datos son opcionales.
% En el Panel TCEE: escribir \cita e Intro (el cursor va al texto; Tab pasa a Autor, Año y Obra).
\NewDocumentEnvironment{cita}{ O{} O{} O{} +b }{%
  \begin{adjustwidth}{2cm}{0pt}
  \raggedright
  #4\par
  \raggedleft
  \if\relax\detokenize{#1#2#3}\relax
    % sin firma
  \else
    #1%
    \if\relax\detokenize{#2}\relax\else\if\relax\detokenize{#1}\relax\else\space\fi(#2)\fi
    \if\relax\detokenize{#3}\relax\else\if\relax\detokenize{#1#2}\relax\else, \fi\textit{#3}\fi
  \fi
  \end{adjustwidth}
}{}



%% ------------------- Listas----------------------%%
    % --- Lista alfabética ---
    \newenvironment{la}
      {\begin{list}{\alph{alphaitem})}{%
          \usecounter{alphaitem}%
          \setlength{\leftmargin}{1cm}%
          \setlength{\parskip}{3pt}% 
          \setlength{\itemsep}{0pt}% ← espacio entre ítems
          \setlength{\parsep}{0pt}%  ← párrafos dentro de un ítem
      }}
      {\end{list}}
      
    % --- Lista numérica ---
    \newenvironment{lnum}
      {\begin{list}{\arabic{numitem}.}{%
          \usecounter{numitem}%
          \setlength{\leftmargin}{1cm}%
          \setlength{\parskip}{3pt}% 
          \setlength{\itemsep}{0pt}% ← espacio entre ítems
          \setlength{\parsep}{0pt}%  ← párrafos dentro de un ítem
      }}
      {\end{list}}
      
% ------------------- Imágenes----------------------%
    \graphicspath{{Img/}}% Rutas por defecto 
    % --- Imagen adaptativa ---
    % Registros de dimensión definidos una sola vez
    \newdimen\imgcap
    \newdimen\imgmin
    \newdimen\imgmax
    \newdimen\imgremain
    \newdimen\imgh
    \newdimen\imgneed
    
    \newcommand{\imagenfit}[3]{%
      \addvspace{10pt}% separación antes
      \par\begingroup
        % Parámetros de composición (asignaciones, no \newdimen)
        \imgcap=1\baselineskip            % alto estimado de título+fuente
        \imgmin=0.15\textheight
        \imgmax=0.15\textheight
        \imgremain=\pagegoal
        \ifdim\pagegoal=\maxdimen \imgremain=0pt\fi
        \advance\imgremain by -\pagetotal
        \advance\imgremain by -\imgcap
        % Decisión de altura destino
        \ifdim\imgremain<\imgmin
          \newpage
          \imgh=0.20\textheight
        \else
          \imgh=\imgremain
          \ifdim\imgh>\imgmax \imgh=\imgmax \fi
        \fi
        % Reservar espacio para evitar cortes del bloque completo
        \imgneed=\imgh \advance\imgneed by \imgcap
        \Needspace{\imgneed}
        % Cuerpo
        \noindent\begin{minipage}{\textwidth}
          \centering
          \captionof{figure}{\textemdash\ #2}\par\vspace{0.3em}% título arriba
          \includegraphics[width=\linewidth,height=\imgh,keepaspectratio]{\detokenize{#1}}\par
          \vspace{0.3em}\small\textit{Fuente: }#3% fuente abajo
        \end{minipage}%
      \endgroup
      \par\addvspace{10pt}%
    }

%------------ Bloque de RECUADRO ESPECIAL -------------%
    \newenvironment{specialblock}[1]{%
      \begin{quote} % crea sangría a izquierda y derecha
      \small % texto más pequeño
      \noindent\textbf{#1}\\[-0.3em] % título en negrita
      \rule{\linewidth}{0.4pt}\\[0.6em]}{
      \end{quote}}
      

%-------------------------- Author Font -----------------------%
    \newcommand{\authorfont}[1]{{\fontfamily{EBGaramond-TLF}\selectfont #1}}

%-------------------------- Anexos -----------------------%
    \makeatletter
    \newcounter{anexo}
    \renewcommand{\theanexo}{\Roman{anexo}}
    \newcommand{\anexo}[1]{%
      \clearpage
      \phantomsection
      \refstepcounter{anexo}%
      \noindent\textbf{Anexo \theanexo.- }#1\par\medskip
      \addcontentsline{stc}{section}{Anexo \theanexo.- #1}%
    }
    \makeatother

    %-------------------------- Lista de figuras (personal) -----------------------%
\makeatletter
\newcommand{\MyListOfFigures}{%
  \clearpage
  \phantomsection
  {\centering\bfseries\Large Índice de figuras\par}%
  \medskip
  % Entrada en nuestro índice de contenido (stc)
  \addcontentsline{stc}{section}{Índice de figuras}%
  % Recuperación del listado (sin cabecera propia)
  \begingroup
    \parindent=0pt\parskip=2pt
    \@starttoc{lof}%
  \endgroup
}
\makeatother

%-------------------------- Bibliografía -----------------------%
\makeatletter
% Contador para las referencias
\newcounter{myref}
% Cabecera de la bibliografía, entra en el índice 'stc'
\newcommand{\MyBibliography}{%
  \clearpage
  \phantomsection
  {\centering\bfseries\Large Bibliografía y referencias\par}%
  \medskip
  \addcontentsline{stc}{section}{Bibliografía y referencias}%
  \setcounter{myref}{0}%
}
\newcommand{\MyBibItem}[4]{%
  \refstepcounter{myref}%
  \noindent[\themyref]\ #1.\ \emph{#2}. #3. #4\par
}
\makeatother

%--------- Establecer cuatro niveles de índice --------%
    \setcounter{tocdepth}{4}   
    \setcounter{secnumdepth}{4}% numera hasta \paragraph

%%%%%%%%%%%%%%%%%%%%%%%%%%%%%%%%

\makeatletter

    %--------- Interlineado Global --------%
    \edef\GlobalBaseLineStretch{\baselinestretch}
    
    %--------- Espaciado de seccinones --------%
    \renewcommand\section{\@startsection{section}
      {1}{\z@}
      {2.5ex \@plus 1ex \@minus .2ex} % espacio antes
      {1.5ex \@plus .2ex}             % espacio después
      {\normalfont\Large\bfseries}    % formato del título
    }
    \renewcommand\subsection{\@startsection{subsection}{2}{\z@}
      {3ex \@plus 0.8ex \@minus .2ex}
      {1.5ex \@plus .2ex}
      {\normalfont\large\bfseries}
    }
    \renewcommand\subsubsection{\@startsection{subsubsection}{3}{\z@}
      {3ex \@plus 0.5ex \@minus .2ex}
      {1ex \@plus .2ex}
      {\normalfont\normalsize\bfseries}
    }
    \renewcommand\paragraph{\@startsection{paragraph}{4}{\z@}%
    {-2ex \@plus -0.5ex \@minus -.2ex}% espacio antes
    {0.8ex \@plus .2ex}% espacio después
    {\normalfont\normalsize\itshape}}% ← cursiva en lugar de negrita

    %--------- Ecuaciones --------%   
    % Variables globales para almacenar títulos y notas
    \gdef\leq@current@section{}%
    \gdef\leq@current@subsection{}%
    \gdef\leq@last@printed@section{}%
    \gdef\leq@last@printed@subsection{}%
    
    % Comando para añadir nota (invisible en el cuerpo)
    \newcommand{\eqnote}[1]{%
      \addcontentsline{leq}{leqnote}{#1}%
    }
    
    % Comando para ecuaciones
    \newcommand{\eqblock}[2]{%
      % Verificar si necesitamos imprimir el título de sección
      \ifx\leq@current@section\leq@last@printed@section\else
        \addcontentsline{leq}{leqsection}{\leq@current@section}%
        \global\let\leq@last@printed@section\leq@current@section
        \gdef\leq@last@printed@subsection{}% Reset subsección al cambiar sección
      \fi
      % Verificar si necesitamos imprimir el título de subsección
      \ifx\leq@current@subsection\leq@last@printed@subsection\else
        \ifx\leq@current@subsection\@empty\else
          \addcontentsline{leq}{leqsubsection}{\leq@current@subsection}%
          \global\let\leq@last@printed@subsection\leq@current@subsection
        \fi
      \fi
      % Ecuación
      \refstepcounter{eqnum}%
      \begin{equation}
        #1\tag{E.\theeqnum}
      \end{equation}%
      \addcontentsline{leq}{leqt}{[E.\theeqnum] -- #2}%
      \addcontentsline{leq}{leqm}{#1}%
    }
    
    %%% ----- Índice de ecuaciones ----------%%%
    % Tamaño para []
    \newcommand{\leqIndexFont}{\normalsize}
    \newcommand{\leqTitleFont}{\tiny}
    \newcommand{\leqNoteFont}{\tiny}
    % Cabecera de sección 
    \newcommand\l@leqsection[2]{%
      \addvspace{25pt}% Mayor separación antes de nueva sección
      \noindent\leqIndexFont
      \makebox[\linewidth]{\rule{.38\linewidth}{0.4pt}\ \ \textbf{  \MakeUppercase{#1}}\ \ \rule{.38\linewidth}{0.4pt}}%
      \par\vspace{-10pt}
    }
    % Cabecera de subsección 
    \newcommand\l@leqsubsection[2]{%
      \par\addvspace{15pt}%
      \noindent\leqIndexFont #1
      \par\vspace{-7pt}
    }
    % Título de ecuación 
    \newcommand\l@leqt[2]{%
    \par\addvspace{10pt}%
      \par\noindent\leqTitleFont #1\par
    }
    % Miniatura de ecuación
    \newcommand\l@leqm[2]{%
      \vspace{0.3\baselineskip}%
      \par\scriptsize\noindent\hspace{1.5em}\(#1\)\par%
    }
    % Nota de ecuación 
    \newcommand\l@leqnote[2]{%
      \vspace{0.2\baselineskip}%
      \par\noindent\leqNoteFont\hspace{1.5em}\textbf{Nota.--} #1\par\vspace{0.5\baselineskip}%
    }
    % Generación del índice
    \newcommand{\mylistofequations}{%
      \clearpage
      \phantomsection
      % Título visual de la lista (ajústalo a tu gusto):
      {\centering\bfseries\Large Índice de ecuaciones\par}
      % Entrada en nuestro índice 'stc':
      \addcontentsline{stc}{section}{Índice de ecuaciones}%
      % Llamada a tu lista real (usa el nombre exacto de tu macro):
      \@starttoc{leq}%
    }
    
    %--------- Captura de títulos de sección --------%
    \let\leq@original@section\section
    \let\leq@original@subsection\subsection
    % Flag para saber si estamos en una sección asterisco
    \newif\ifLeqStarSection
    \LeqStarSectionfalse
    
    % Redefinir \section CON CONTROL DE ASTERISCO
    \renewcommand{\section}{%
      \@ifstar{\leq@section@star}{\@dblarg\leq@section@nostar}%
    }
    \def\leq@section@star#1{%
      \global\LeqStarSectiontrue
      \gdef\leq@current@section{#1}%
      \gdef\leq@current@subsection{}%
      \leq@original@section{#1}%
      \global\LeqStarSectionfalse
    }
    \def\leq@section@nostar[#1]#2{%
      \global\LeqStarSectionfalse
      \gdef\leq@current@section{#2}%
      \gdef\leq@current@subsection{}%
      \leq@original@section[#1]{#2}%
    }
    % Redefinir \subsection
    \renewcommand{\subsection}{%
      \@ifstar{\leq@subsection@star}{\@dblarg\leq@subsection@nostar}%
    }
    \def\leq@subsection@star#1{%
      \global\LeqStarSectiontrue
      \gdef\leq@current@subsection{#1}%
      \leq@original@subsection{#1}%
      \global\LeqStarSectionfalse
    }
    \def\leq@subsection@nostar[#1]#2{%
      \global\LeqStarSectionfalse
      \gdef\leq@current@subsection{#2}%
      \leq@original@subsection[#1]{#2}%
    }

    % ----- Restauración de valores Globales de estilo ----------%
    % Flags
    \newif\ifInFootnote
    \InFootnotefalse
    \pretocmd{\@footnotetext}{\InFootnotetrue}{}{}
    \apptocmd{\@footnotetext}{\InFootnotefalse}{}{}
    % Profundidad de listas (soporta anidadas)
    \newcount\MyListDepth \MyListDepth=0
    \AtBeginEnvironment{la}{\advance\MyListDepth by 1}
    \AfterEndEnvironment{la}{\advance\MyListDepth by -1}
    \AtBeginEnvironment{lnum}{\advance\MyListDepth by 1}
    \AfterEndEnvironment{lnum}{\advance\MyListDepth by -1}
    \AtBeginEnvironment{itemize}{\advance\MyListDepth by 1}
    \AfterEndEnvironment{itemize}{\advance\MyListDepth by -1}
    \AtBeginEnvironment{enumerate}{\advance\MyListDepth by 1}
    \AfterEndEnvironment{enumerate}{\advance\MyListDepth by -1}
    % Restauración diferida: hazlo al INICIO del próximo párrafo real
    \newif\if@pendingRestore
    \@pendingRestorefalse
    \newcommand{\RequestRestore}{\global\@pendingRestoretrue}
    \everypar{%
      \if@pendingRestore
        \global\@pendingRestorefalse
        \ifInFootnote\else
          \ifnum\MyListDepth=0
            \setlength{\parskip}{\GlobalParskip}%
            \setstretch{\GlobalBaseLineStretch}%
          \fi
        \fi
      \fi
    }
    % Al final de bloques:
    \AfterEndEnvironment{equation}{\RequestRestore}
    \AfterEndEnvironment{equation}{\RequestRestore}
    \AfterEndEnvironment{la}{\RequestRestore}
    \AfterEndEnvironment{lnum}{\RequestRestore}
    % Al final de macros:
    \apptocmd{\eqblock}{\RequestRestore}{}{}
    \apptocmd{\imagenfit}{\RequestRestore}{}{}
    
\makeatother

% ===================== ToC propio con 4 niveles (único bloque) =====================
\makeatletter

% --- Escritores: añadimos una línea al .stc en cada cabecera numerada ---
% Guardamos las versiones actuales para llamar al comportamiento normal
\let\MyTOC@orig@section\section
\let\MyTOC@orig@subsection\subsection
\let\MyTOC@orig@subsubsection\subsubsection
\let\MyTOC@orig@paragraph\paragraph

% SECTION
\renewcommand{\section}{%
  \@ifstar{\MyTOC@section@star}{\@dblarg\MyTOC@section@nostar}%
}
\def\MyTOC@section@star#1{%
  \MyTOC@orig@section{#1}% sin entrada al .stc
}
\def\MyTOC@section@nostar[#1]#2{%
  \MyTOC@orig@section[{#1}]{#2}% comportamiento normal (escribe al .toc)
  % Escribimos también nuestra línea al .stc (texto con \numberline)
  \addcontentsline{stc}{section}{\protect\numberline{\thesection}#2}%
}

% SUBSECTION
\renewcommand{\subsection}{%
  \@ifstar{\MyTOC@subsection@star}{\@dblarg\MyTOC@subsection@nostar}%
}
\def\MyTOC@subsection@star#1{%
  \MyTOC@orig@subsection{#1}%
}
\def\MyTOC@subsection@nostar[#1]#2{%
  \MyTOC@orig@subsection[{#1}]{#2}%
  \addcontentsline{stc}{subsection}{\protect\numberline{\thesubsection}#2}%
}

% SUBSUBSECTION
\renewcommand{\subsubsection}{%
  \@ifstar{\MyTOC@subsubsection@star}{\@dblarg\MyTOC@subsubsection@nostar}%
}
\def\MyTOC@subsubsection@star#1{%
  \MyTOC@orig@subsubsection{#1}%
}
\def\MyTOC@subsubsection@nostar[#1]#2{%
  \MyTOC@orig@subsubsection[{#1}]{#2}%
  \addcontentsline{stc}{subsubsection}{\protect\numberline{\thesubsubsection}#2}%
}

% PARAGRAPH
\renewcommand{\paragraph}{%
  \@ifstar{\MyTOC@paragraph@star}{\@dblarg\MyTOC@paragraph@nostar}%
}
\def\MyTOC@paragraph@star#1{%
  \MyTOC@orig@paragraph{#1}%
}
\def\MyTOC@paragraph@nostar[#1]#2{%
  \MyTOC@orig@paragraph[{#1}]{#2}%
  % Si no numeras los \paragraph, cambia \theparagraph por texto plano sin \numberline
  \addcontentsline{stc}{paragraph}{\protect\numberline{\theparagraph}#2}%
}

% --- Impresor del índice propio (lee el .stc) ---
\newcommand{\MyTOC}{%
  \par\bigskip
  {\centering\bfseries\Large Índice de contenido\par}%
  \medskip
  \begingroup
    \parindent=0pt\parskip=2pt
    \@starttoc{stc}%
  \endgroup
}

\makeatother
% ===================== fin ToC propio =====================


%%%%%%%%%%%%%%


% --- Datos del tema ---
\title{Sistema Fiscal en España (III): Fiscalidad Internacional\\
\large Impuesto sobre la Renta de los No Residentes. Convenios de doble imposición, con especial referencia al modelo de la OCDE. El problema de la erosión de bases imponibles y propuestas de solución}
\author{Marcelo Ortega}
\date{Septiembre 2026}

\begin{document}
\maketitle
\newpage

% --- Página de resumen y modificaciones ---
\puntosclave{ %% Escribir puntos clave Apuntes CECO
     Éste es un tema de Derecho tributario, lo que implica dos cosas:
     \begin{la}
         \item Por un lado, se requiere rigor jurídico en el empico de los conceptos técnicos;
         \item Por otro, el opositor debe priorizar el trasfondo económico que existe detrás de la regulación. 
     \end{la}
    
    \textcolor{red}{\textbf{OJO ->} No es necesaria la memorización de los artículos, directivas ni jurisprudencia}. Sin perder de vista el contenido jurídico, se debe priorizar y demostrar que se comprenden los conceptos económicos fundamentales del Tema: arbitraje, interacción estratégica, etc.
    
    Se ha adoptado la sucesión de epígrafes adoptada en el BOE (en el temario de la oposición). Por tanto, existe mucha flexibilidad para escoger el lugar concreto en el que tratar cada materia, ampliar contenido, etc.\footnote{%
        Por ejemplo, jurídico. Lo más natural sería ampliar con uso exclusivo de fuentes primarias (por ser un tema de tanta actualidad). Por ejemplo, profundizar en temas concretos a través de la web BEPS, texto consolidado del IRNR, etc. Complementariamente se podría profundizar en la bibliografía del tema original.
    } Los bloques responden a cada una de las vertientes donde se valoran y tratan las consecuencias de una economía abierta: doméstico, bilateral y multilateral. La distribución del tiempo debería ser equilibrada, incluso con un mayor peso al enfoque doméstico.
    }
\vspace{1cm}

\modificaciones{
    \textcolor{red}{OJO ->} Contenido tiene que estar mucho más enfocado a lo económico.

    Eliminar el Segundo bloque. La distribución de tiempos tiene que ser: 7' (IRNR, Doble Imposición) y luego 12' (Multiletaral), y enfocarlo mucho más económico presentando los problemas, país pequeño país grande. 

    \textcolor{red}{\textbf{NOTA.-} Muy interesante la excepción por \textbf{deuda pública}: para no distorsinoar las decisiones en un contexto de economía abierta, sin influencia sobre el interés mundial!!! Poner de ejemplo en el tema correspondiente y aquí.}
    }

\newpage
\MyTOC
\newpage

\mylistofequations
\clearpage

\MyListOfFigures
\clearpage


\pagenumbering{arabic}
\setcounter{page}{1}


\section{Introducción}



\subsection{Contextualización}


\subsubsection{Relevancia}




\subsubsection{Problemática}



\subsection{Estructura de la exposición}




La globalización y la digitalización está poniendo de manifiesto las grietas de la vigente arquitectura de la fiscalidad internacional, en especial en el ámbito del impuesto sobre sociedades.

Estos dos fenómenos acentúan el trade-off en materia fiscal al que se enfrentan los Estados: elegir unos tipos impositivos que permitan el sostenimiento del sistema, pero que les permitan mantenerse competitivos para evitar la elusión/huida de contribuyentes y capital. Es decir, ha provocado una competencia fiscal feroz entre Estados que buscan atraer capital mediante el uso activo de su autonomía tributaria. Además de las evidentes oportunidades a la elusión fiscal que ofrecen la integración vertical de empesas multinacionales, con sus procesos de fabricación/distribución/actividad fragmentados en un número muy variado de países, erosionando las bases imponibles y permitiendo el traslado de beneficios a jurisdicciones más favorables. 

El marco de coordinación internacional estaba construido sobre un conjunto de principios reglas basados: la presencia física de las empresas para desarrollar sus negocios, la coordinación y adopción de compromisos bilaterales, etc. Sin embargo, la digitalización se ha sumado al fenómeno anterior, transformando una parte importante de la actividad económica y facilitando a las empresas el desarrollar su actividad en países en los que no tienn presencia física. 

En definitiva, este cóctel ha provocado y, en cierto modo, agravado tras el transcurso de los años la competencia fiscal entre estados, la planificación fiscal agresiva de las empresas multinacionales (EMN), erosionado las bases imponibles mediante la fragmentación global de la producción (operaciones intraempresa) y el traslado de beneficios a las jurisdicciones más favorables, incluso con el uso de precios de transferencia. Las implicaciones también son inequívocas. Se ha experimenta una reducción notable de la recaudación proveniente de determinadas fuentes, una proliferación de nuevas distorsiones fiscales o agravamiento de las existentes para suplir este vacío recaudatorio, una asignación de recursos ineficientes provocado por la competencia fiscal, que oscurece las virtudes de la ventaja comparativa pura, y un agravamiento de las distorsiones en la estructura de los mercados, pues estas dinámica tiende a reforzar la posición dominante de empresas multinacionales y grandes empresas digitales.

¿Qué podemos hacer al respecto? Existen diversas alternativas. Tradicionalmente estos retos se han intentado resolver/parchear en diferentes niveles. Empezando por el nivel unilateral, con disposiciones especiales para la tributación de no residentes, siguiendo por las soluciones bilaterales, siendo los Convenidos de Doble Imposición (CDI) el ejemplo paradigmático; y, más recientemente, en la esfera multilateral, liderada por organismos internacionales, en buene medida obligados por sus Estados Miembros, mediante acciones de coordinación o iniciativas conjuntas que tratan de resolver los problemas de fondo.

Con el fin de presentar, de manera clara y didáctica, el objeto principal de este tema, estructuraremos la exposición en tres partes, que de manera no exhaustiva nos permitirán comprender los problemas y soluciones que han ido surgiendo desde la óptica español. Por tanto, trataremos:
\begin{lnum}
    \item En primer lugar, las soluciones unilaterales a la erosión de bases, cuyo ejemplo claro es el Impuesto sobre la Renta de los No Residentes;
    \item En segundo lugar, nos centraremos en los problemas que emanan entre países y las soluciones bilaterales que se aplican como remedio; y,
    \item Por último, trataremos los problemas de ámbito local, presentando las principales iniciativas de ámbito multilateral llamadas a solucionarlos. 
\end{lnum}




\clearpage
\section{Fiscalidad internacional}
\puntosclave{
    La fiscalidad internacional se enfrenta a dos problemas adicionales: la doble imposición internacional y la no imposición internacional. 
    
    La D11 poder ser jurídica y económica. Existen múltiples vías para eliminarla. tanto a nivel unilateral, bilateral y multilateral. 
    
    La elusión fiscal no implica incumplimiento legal (a diferencia de la evasión) pero supone un uso interesado de la misma y puede estar sujeta a puede estar sujeta a escrutinio y regulación por parte de las autoridades fiscales para evitar abusos. Un ejemplo claro son las iniciativas desarrolladas en BEPS por la OCDE y el Marco Inclusivo.
}

La fiscalidad internacional es el conjunto de normas que regula el funcionamiento de los sistemas fiscales en su vertiente exterior, en el marco de las relaciones internacionales.

La fiscalidad internacional se rige por cuatro fuentes del Derecho Tributario Internacional.

El primero, y más importante, es el Derecho Tributario autónomo (como por ejemplo el IRNR en España). Dentro de esta categoría deberíamos también encuadrar el Derecho Tributario Europeo, pues al emanar de las Instituciones de la Unión, constituye parte integral de nuestro ordenamiento jurídico y no debe ser tratado ni como una extensión ni como un ordenamiento diferente. El segundo, y más relevante históricamente, es el Derecho Tributario Internacional. Este es el cuerpo normativo internacional propiamente dicho, formado por tratados internacionales, con eficacia inmediata a nivel nacional (porque se incorporan inmediatamente al ordenamiento jurídico tras su publicación en el BOE) preeminencia sobre el ordenamiento jurídico nacional (esta se deriva del art. 96 CE: ·'Sus disposiciones solo podr,1n ser derogadas. modificadas o suspendidas en la fomrn prevista en los propios Tratados o de acuerdo con las normas generales de Derecho internacional").
4. Soft Law, que son normas. directrices y estándares no vinculantes que se establecen a nivel internacional. No son obligatorias ni jurídicamente vinculantes. Ejemplos: modelo de Convenio de la OCDE para evitar la doble imposición. Directrices de la OCDE sobre precios de transferencia. o los estándares de transparencia e intercambio de información fiscal de la OCDE,

Dentro de los tratados internacionales. puede haber tratados bilaterales (CDIs) o multilaterales, por ejempol, el Instrumento Multilateral para la modificación de CDI en el marco de la acción 15 del BEPS o el Convenio Multilateral del Pilar I de BEPS 2.0 que veremos a lo largo del tema).

\textbf{Problemas de la fiscalidad internacional}
La fiscalidad internacional, por sus particularidades, ha tenido que hacer frente a dos tipos de problemas que operan en sentido opuesto: la doble imposición y la doble no imposición.

La \textbf{doble imposición internacional} se produce cuando el mismo hecho imponible resulta gravado en dos o más jurisdicciones fiscales. Una de las causas habituales es la aplicación de distintos principios de tributación. Por ejemplo, si en un país existe una figura de imposición directa que se construye de manera personal (obligación personal), mientras que en otro se construye como una obligación real (principio de territorialidad), el mismo hecho imponible será gravado dos veces. 

El caso más claro de este problema se da actualmente en la mayoría de países desarrollados, que establecen la obligación personal (IRPF e IS en España) sobre sus residentes, y relegan el criterio real a sus no residentes (IRNR en España). Gravando, de facto, dos veces el mismo hecho imponible, no en una misma jursidicción pero sí en dos. 

Evidentemente esto merma la rentabilidad de las inversiones, distorsiona la asignación de recursos y viola los principios de neutralidad internacional. Además, puede tener efectos sobre la movilidad internacional de factores productivos, tanto trabajadores como capitales, y provocar la deslocalización de inversiones y riqueza/renta a territorios de baja o nula tributación. 

Por otro lado, la \textbf{doble no imposición internaional}, o erosión de bases, se produce por varias causas. La más evidente sucede cuando no se contempla la obligación tributaria, per se, bien se de carácter real o personal, fenómeno que conocemos como no imposición pura. También puede darse como consecuencia de la planificación fiscal internacional, como el uso abusivo de los Convenios de Doble Imposición o de precios de transferencia, fenómeno que conocemos como erosión de bases en sentido estricto (merma la base, pero no la anula).

\begin{specialblock}{Doble Imposición Internacional: Dos tipos diferentes}
    Se pueden distinguir dos tipos de doble imposición: jurídica y económica.

    La doble imposición jurídica ocurre cuando dos o más países gravan el mismo hecho imponible, sobre el mismo sujeto pasivo y con de la misma forma o muy similar. Las causas pueden ser muy diversas.
    
    Por ejemplo, el sujeto puede considerarse residente en ambos dos estados: un país de fine la residencia en función del número de días y el otro como centro de intereses económicos. También puede darse porque un estado establezca una obligación real sobre las rentas de sus no residentes y otro como obligación personal sobre sus residentes: el sujeto pasivo tributará dos veces por una fracción de su hecho imponible, una vez como obligación real y otra como obligación personal. Por último, puede también suceder que dos países graven simultáneamente la renta de un sujeto no residente. Por ejemplo, imaginemos un sujeto que reside en Alemania, pero genera rentas en Chile, a través de una empresa que opera en toda Sudamérica. Tanto Chile como el resto de países suramericanos podrían gravar sus rentas aplicando el principio de fuente (obligación real).\footnote{%
        Este caso particular es muy peliagudo puesto que, en la mayoría de casos, no se resuelve mediante Convenio de Doble Imposición. Como es natural, los Convenios de Doble Imposición solo se extienden a los sujetos residentes en alguno de los países con los que hay dicho compromiso. Como nuestro sujeto no reside ni en Chile ni, pongamos, Bolivia, el Convenio que hubiere entre estos no sería de aplicación. En todo caso si aplicaría el que hubiere entre Chile y Alemania, si lo hubiere.
    } Como podemos suponer, el segundo caso es el más habitual, cuando un país utiliza el principio personalista y otro de territorialidad.
   
    Por otro lado, la doble imposición económica ocurre cuando dos o más países gravan la misma renta más de una vez, pero a personas distintas. Los ejemplos más claros se dan sobre dividendos y sobre los ajustes en los precios de transferencia.

    El primer caso es evidente cuando el beneficio de una sociedad se grava en sede (residencia) y en el país perceptor del dividendo. El modelo de convenio propuesto por la OCDE prefiere dejar a la elección de los Estados la solución al problema, sea adoptando medidas unilaterales o por vía convencional. El segundo caso es más particular. Cuando una empresa realiza operaciones entre empresas del grupo o vinculadas están obligadas a aplicar precios de transferencia que respeten el principio de plena competencia que, básicamente, implica utilizar precios de mercado en las operaciones de disposición intraempresa. La doble imposición se podría producir si un país obliga a la realización de ajustes en los precios de transferencia en base a un CDI y el otro país no obliga o impide la realización de estos ajustes en sentido contrario.

    Como vemos, la diferencia entre ambas es que la jurídica se produce en un mismo contribuyente, y la económica en dos o más contribuyentes. Tradicionalmente, la doble imposición económica ha recibido menos atención.
    
    Sin embargo su importancia es creciente debido al creciente papel, en términos de volumen relativo, que tienen los grupos multinacionales en las operaciones transfronterizas. Estos grupos aplican políticas de precios de transferencia que permiten que los ingresos impositivos de los distintos países se vean afectados, ya que los utilizan para desviar los beneficios a países de baja tributación.

    
\end{specialblock}

Las soluciones para eliminar la DII
Put::den ser unilaterales, bilaterales o multilaterales, y pueden usarse de forma complementaria, no son excluyentes entre sí:
1 .Unilaterales: las que adopta un país independientemente del resto de países. Puede tomar básicamente dos formas:
1 . Exención y n o sujeción: e l país renuncia a la imposición de rentas susceptibles de imposición en otro Estado. Por ejemplo las exenciones del IS español para determinadas rentas de fuente extranjera cuando se cumplen ciertos requisitos, o del IRNR que estudiaremos en el apartado 111.
11. Deducciones de la carga fiscal sufrida por una renta en atención a los impuestos pagados en otro país. Por ejemplo deducciones por O l l de algunas figuras impositivas en España como el IRPF, el I S o el JRNR.

2-Bilaterales: mediante l a forna de CDI que veremos en el apartado IV.

3-M ultilaterales, mediante la finna de:
1 . Acuerdos de integración amplios que incluyan nonnativa e n materia de fiscalidad: Mercosur, TLCAN (previamente denominado N A FTA), UE. En la UE hay numerosas disposiciones, por ejemplo la Directiva de interés y cánones (Directiva 2003/49/CE)
n. Tratados multilaterales específicos de ámbito fiscal, como el I nstrumento Mul tilateral de BEPS (MLI) o el Convenio Multilateral (MLC), si bien éste último está más relacionado con evitar la DNII que la DII como veremos en el apartado V.


\textcolor{red}{\textbf{OJO ->} Sería interesante decir que es perfectamente legal, pero explicar las razones económicas por las que no es óptimo}

Con esta final idad. las jurisdicciones incluyen en su Derecho interno las denominadas normas anti-abuso, que pueden ser de dos tipos: GAAR (General Anti-Avoidance Rules) y SAAR (Specitic Anti-Avoidance Rules). Por un lado. las primeras establecen principios que determinan cuándo las normas tributaria!> se están aplicando de manera contraria a su espíritu. con el mero fin de obtener un tratamiento fiscal más ventajoso del que tiene previsto el legislador. Por otro, las SAAR son disposiciones anti-abuso específicas para una norma tributaria en concreto (por ejemplo, la tributación de los dividendos).





\clearpage
\section{Dimensión doméstica: Impuesto sobre la Renta de los No Residentes}
\puntosclave{
    El IRNR. es un i mpuesto directo, real. no progresivo n i cedido, que grava a los no residentes, por las rentas obtenidas en España. 
    
    Es m uy i mpm1ante distinguir si las rentas se obtienen mediante establecimiento permanente o no, ya que afecta a l a normativa aplicable. 
    
    Cuando existe establecimiento permanente el JRNR se guía por el IS, y se liquida al final de afio de forma similar al IS 
    
    Cuando no hay establecimiento permanente el IRNR se guía por el IRPF. y se l iquida un impuesto por cada devengo de renta, lo que junto con la obligación general de retener y la figura del responsable. perm ite reducir la evasión.
}

Correcciones que conviene hacer en el resto del bloque 3

Al revisar el bloque he detectado varios errores o desactualizaciones fuera del 3.2. Te los señalo para que decidas si los corriges:

La directiva sobre rendimientos del ahorro. [Seguro] El documento cita una "Directiva 2014/48/UE sobre fiscalidad de los rendimientos del ahorro". La directiva sobre fiscalidad del ahorro era la 2003/48/CE, y fue derogada en 2015 por la Directiva (UE) 2015/2060, al quedar absorbida por el intercambio automático de información de la DAC2 que vimos en el 5.4.2. La exención de intereses a residentes en la UE sigue vigente, pero su justificación actual es la libre circulación de capitales y el intercambio automático de información, no esa directiva.
La figura del responsable. [Probable] El documento dice que el sujeto pasivo "deberá nombrar un responsable". En realidad, el responsable no se nombra: la propia ley declara responsables solidarios a quienes pagan las rentas o custodian los bienes del no residente en España (art. 9 del texto refundido).
La opción de tributar por el IRPF de los residentes en la UE. [Probable] El documento presenta como requisitos acumulativos lo que en realidad son dos vías alternativas: que las rentas del trabajo y de actividades económicas obtenidas en España supongan al menos el 75% de su renta total, o bien que su renta en España sea inferior al 90% del mínimo personal y familiar y su renta fuera de España también sea inferior a ese mínimo. En ambos casos esas rentas deben haber tributado efectivamente por el IRNR. Conviene verificarlo con el art. 46 del texto refundido.
El tipo de los establecimientos permanentes. [Probable] El documento menciona un tipo general "24,75 → 24%" para contribuyentes con establecimiento permanente. Hoy el establecimiento permanente tributa al tipo del Impuesto sobre Sociedades, con carácter general el 25%; el 24% es el tipo general de las rentas sin establecimiento permanente.
Referencias a paraísos fiscales. La nota 6 ya recoge el cambio a "jurisdicciones no cooperativas"; puede remitirse ahora al 4.3.4, donde está desarrollado.


Tradicionalmente en España la tributación de residentes y no residentes se regulaba en la misma normativa que el Impuesto de la Renta de las Personas Físicas y el Impuesto sobre Sociedades, que aplicaba el criterio de obligación personal de contribuir para los residentes en territorio español y la obligación real de contribuir para los no residentes. Esta situación se modificó en 1998 con la entrada en vigor de la Ley sobre la Renta de No Residentes.





\subsection{Marco jurídico: TR de la Ley del Impuesto sobre la Renta de No Residentes (TRLIRNR, 2004)}
Comencemos estudiando el Impuesto sobre la Renta de los No Residente, la figura impositiva española para estos casos. Su regulación se encuentra recogida en el Texto Refundido de la Ley del Impuesto sobre la Renta de No Residentes (TRLIRNR) aprobado por Real Decreto Legislativo 5/2004, de 5 de marzo. Es una norma muy breve ya que en muchas cuestiones remite a la normativa del IRPF y el IS.

\textcolor{red}{[\textbf{Reformular y explicar mejor esto. No se entiende muy bien}] Mediante la Ley 26/2014, el IRPF e IRNR fueron ligeramente modificados. Respecto al IRNR, introdujo medidas para favorecer las libertades de circulación en el ámbito de la Unión, como el calculo de los gastos deducibles para contribuyentes sin establecimiento permanente - las personas físicas según la norma del IRPF, y las personas jurídicas según la norma del IS: incluye un nuevo supuesto para favorecer la libre circulación de ciudadanos de la Unión por el que se permite optar por tributar por el IRPF a contribuyentes residentes en otros Estados de la Unión de bajos ingresos; y admite la exención de ganancias patrimoniales por transmisión de la vivienda habitual en España si se aplica reinversión, al igual que ocurre en el IRPF. También se ajustaron elementos del IRNR como consecuencia de cambios en el IRPF, como los tipos de gravamen. El tipo general para contribuyentes con establecimiento permanente pasó del $24,75 \to 24\%$ actual, mientras que para residentes en otros Estados de la Unión el tipo se estableció en el $19\%$, coincidente con el tipo marginal más bajo de la tarifa del IRPF. Adicionalmente, se equiparó el tipo aplicable a los establecimiento permanente al que corresponda con arreglo a la normativa del IS.}

Se trata de un impuesto directo, pues grava una manifestación de capacidad de pago directa, concretamente las rentas obtenida por no residentes, personas físicas y jurídicas en España y, como veníamos diciendo, establece una obligación de carácter real. Es un impuesto administrado y gestionado íntegramente por el Estado, ya que los diferentes niveles de Gobierno, subsectores, ni tienen capacidad normativa ni comparten la recaudación.\footnote{%
    Evidentemente, la exceción común aplica, puesto que el ámbito de aplicación territorial excluye tanto a los Territorios Históricos, en los que prevé una abstención completa, y particularidades propias para Canarias, Ceuta y Melilla, aplicables en virtud de su normativa específica y de lo dispuesto en el TRLIRN.
}

Estudiaremos en primer lugar las características y elementos definitorios y continuaremos con la liquidación del impuesto distinguiendo entre las situaciones en las que existe o no establecimiento permanente.


\subsubsection{Hecho Imponible}
El hecho imponible que grava el impuesto es el aumento de la capacidad de pago de un sujeto por fuentes de renta que se encuentren en la jurisdicción española. El periodo impositivo es anual, aunque su gestión, como veremos, difiere atendiendo a la circunstancia concreta del sujeto. Por último, el devengo se produce el último día del periodo impositivo.

Aunque este aumento de la capacidad de pago puede originarse por innumerables fuentes, el art. $13$ de la ley enumera algunas, como: (i) rentas de actividades económicas obtenidas tanto con establecimiento permanente como sin establecimiento permanente;\footnote{%
    Cuando no media establecimiento permanente el control es mucho más difícil. Estas actividades pueden incluir prestaciones de servicios en España y actuación personal de artistas o deportistas en España, etc.
} (ii) rendimientos del trabajo, incluido pensiones; (iii) rendimientos del capital mobiliario (dividendos, intereses y cánones); (iv) rendimientos del capital inmobiliario, inmuebles situados en territorio español; (v) rentas imputadas por inmuebles urbanos en territorio español; y, (vi) ganancias patrimoniales. 

Por el contrario, no se consideran rentas obtenidas en España: (i) las derivadas de compraventas internacionales de mercancías; y, (ii) \textcolor{red}{las rentas pagadas a no residentes por establecimiento permanente situados en el extranjero cuando dichas rentas son por trabajos relacionados con la actividad del establecimiento permanente situado en el extranjero. [\textbf{explicar mejor que es esto}]}

Se consideran exentas las rentas que en el IRPF lo estén y sean percibidas por personas físicas. Por ejemplo las pensiones por invalidez. Existen además un conjunto de rentas adicionales exentas que persiguen incentivar la inversión en España o evitar la doble imposición. Por ejemplo, los intereses satisfechos a residentes en otros Estados Miembros, los rendimientos de las cuentas de no residentes, los rendimientos por tenencia de títulos de deuda pública (obtenidos sin establecimiento permanente), así como las exenciones para evitar la doble imposición por aplicación de un CDI.

\begin{specialblock}{\textbf{Exenciones:} Listado completo}
    Están exentas del IRNR, las siguiente rentas:
    \begin{lnum}
        \item Las rentas previstas en el art. $7$ de la LIRPF, siempre que sean percibidas por persona física. Las tres más importantes son las prestaciones por invalidez, indemnizaciones por despido y las prestaciones familiares (por hijo a cargo, orfandad, maternidad, y paternidad). Suponen dos tercios del total de rentas exentas;
        \item Las becas percibidas por personas físicas satisfechas por Administraciones Públicas;
        \item Intereses y ganancias patrimoniales mobiliarias comunitarias. Es decir, los rendimientos del capital intangible así como las ganancias patrimoniales por enajenación de bienes muebles sin mediación ele establecimiento permanente. Siempre y cuando el sujeto sea residentes en otro Estado de la Unión o del Espacio Económico Europeo o por establecimientos permenantes de dichos residentes situados en otro Estado de la Unión o del Espacio Económico Europeo.\footnote{%
            Se aplica también al EEE siempre que exista acuerdo de intercambio de información. Esto es aplicable para todos los casos de este capitulo: cuando una norma es aplicable a un Estado de la Unión, lo es también a los países del EEE que no pertenecen a la Unión siempre que exista un acuerdo de intercambio de información.
        } Esta exención es la traspoción de la Directiva 2014/48/UE sobre fiscalidad de los rendimientos del ahorro que determina que los rendimientos percibidos en un Estado por personas físicas de otro Estado deben tributar efectivamente en el Estado Miembro de residencia (es decir en ese segundo país) y por tanto se aplica exención en origen de los rendimientos. El objetivo de esta Directiva es garantizar la libre circulación de capitales. No se aplica la exención a las ganancias patrimoniales derivadas de la transmisión de acciones en sociedades o cuando el activo sea fundamentalmente de carácter inmobiliario, cuando se trate de personas físicas y hubiera tenido en los doce meses anteriores una participación superior al $25\%$ del capital de la entidad o, por último, cuando no se cumplan los requisitos del art. $21$ LIS .
        \item Los rendimientos generados por los títulos de Deuda Pública, sin mediación de establecimiento permanente;
        \item Las rentas derivadas de valores emitidos en España por personas físicas o entidades no residentes sin mediación de establecimiento permanente (bonos matador);
        \item Los rendimientos de cuentas de no residentes satisfechos a contribuyentes del IRNR, a menos que el pago sea realizado a un establecimiento permanente en España por cualquier institución registrada como inermediario de crédito y depositante;
        \item Las rentas derivadas de la transmisión de valores o el reembolso de participaciones en fondos de inversión realizadas en mercados secundarios oficiales de valores españoles, obtenidas por personas físicas o entidades no residentes sin mediación de establecimiento permanente. que sean residentes en un Estado que tengan suscrito con España un convenio con cláusula de intercambio de información.\footnote{%
            Los CDI basados en el Modelo OCOE u ONU, contienen siempre un articulo de intercambio de información. Además, varios países han firmado recientemente Acuerdos de Intercambio de Información Fiscal con países con los que no tienen un convenio, como con paraísos fiscales. El intercambio de información también está cubierto por acuerdos multilaterales entre los países nórdicos, la Unión, y un Convenio conjunto sobre Asistencia Administrativa Mutua en Materia Fiscal del Consejo de Europa y la OCOE, que entró en vigor en 1995 y, que después de un comienzo lento, más de noventa países ya lo han firmado.
        }
        \item Los beneficios distribuidos por sociedades filiales residentes en España a sus sociedades matrices residentes en otros Estados Miembros de la Unión cuando se cumplan los siguientes requisitos: (i) que ambas sociedades revistan una forma jurídica de las previstas en el anexo de la Directiva 201 1/96/UE; (ii) que ambas sociedades estén sujetas, y no exentas, a tributación por sus beneficios societarios en España y por un impuesto equivalente en el otro Estado de la Unión; (iii) la distribución del beneficio no sea consecuencia de la liquidación de la filial; (iv) tienen la consideración de sociedad matriz aquella que posea una participación directa de al menos el $5\%$ del capital. 
        \item Los dividendos y participaciones en beneficios obtenidos por fondos de pensiones equivalentes e instituciones de inversión colectiva de la Directiva 2009/65/CE sin establecimiento permanente.
        \item Los cánones o regalías satisfechos por una sociedad residente o un establecimiento permanente en España a una sociedad residente en otro Estado o a un establecimiento permanente en otro Estado, cuando concurran determinadas condiciones. Por ejemplo, tratarse de sociedades asociadas, que la sociedad que reciba los pagos lo haga en su propio beneficio y que los pagos estén relacionados con la actividad.
    \end{lnum}

    \textcolor{red}{\textbf{NOTA.-} Se podría dar una vuelta adicional a todas la lista expuesta. Con información y ejemplos para que se entendiese bien cada cosa y circunstancia, así como su justificación}
\end{specialblock}

\subsubsection{Sujeto Pasivo}
Es sujeto pasivo del IRNR cualquier persona, física o jurídica, no residente y que obtengan rentas cuya fuente está en España. Para determinar la condición de sujeción, esto es, la no residencia, la Ley se remite explícitamente a los criterios del IRPF e IS ([\authorfont{Ver Tema 4.B.17}]).\footnote{%
    En el caso de las personas físicas: $183$ días en territorio español o tener el centro de interés económico en España. Como presunción \textit{iuris tantum}, también se considera residente aquel cuyo cónyuge y/o hijos menores residan habitualmente en España. En el caso de las personas jurídicas: constituidas bajo la LSC, domicilio social en España o sede de dirección efectiva en España.
}

Por tanto, el ámbito de aplicación subjetivo del IRNR es \textbf{\textit{sensu contrario}}: cualquier persona, física o jurídica, que no sea residente a efectos del IRPF o IS. La única excepción son los residentes de cualquier otro Estado de la Unión, que pueden optar por tributar por el IRPF (a pesar de ser no residentes) si concurren dos circunstancias (aparte de residencia):\footnote{%
    \textbf{NOTA.-} Siempre y cuando dicho territorio no sea considerado un \textit{paraíso fiscal}. Aunque la denominación de de paraíso fiscal está en desuso y, desde 2021, se viene sustitutyendo por el de \textit{jurisdicción no cooperativa}, que incluye aquellas que no comparten información fiscal de forma efectiva o los territorios de baja o nula tributación.
} (i) que las rentas del trabajo o de actividades económicas obtenidas en España sean al menos un $75\%$ de su renta global; y, (ii) \textcolor{red}{[\textbf{Habría que explicar mejor las dos circunstancias, no se entiende bien.}] que dichas rentas hayan tributado efectivamente por el IRNR}; (iii) que la renta obtenida durante el ejercicio en España haya sido inferior al $90\%$ del mínimo personal y familiar que le hubiese correspondido de acuerdo con el IPRF de haber sido residente en España. Requiere además que la renta obtenida fuera de España haya sido también inferior a dicho mínimo. Esta segunda condición se introduce en la reforma de 2006, su intuición económica es favorecer la libre circulación de ciudadanos de la UE, evitando discriminar a ciudadanos de bajos ingresos de otros Estados frente a los españoles, que podrían aplicar la exención del mínimo personal y familiar.

Además, como \textbf{excepción} contraria, se considera que los miembros de misiones diplomáticas, oficinas consulares o aquellos que desempeñen un cargo público de la misma índole en España (trabajando para un empleador de Derecho Público extranjero), son no residentes a efectos fiscales, independientemente de su residencia formal/material, y por tanto quedan sujetos al IRNR. \textcolor{red}{\textbf{Hay que explicar mejor esto, ¿qué es este régimen? Incluir/Mejorar explicación y poner un ejemplo}Por último, las \textbf{entidades en régimen de atribución de rentas} constituidas en el extranjero y con actividades económicas en territorio español (art. 38 TRLIRNR) serán también consideradas sujetos pasivos a efectos del IRNR.}

Como se puede intuir, una de las características diferenciales del IRNR es que al gravar a sujetos no residentes, se tiene un escaso contacto con el y existe mayor riesgo de impago. Con esto en mente, el legislador ha establecido un sistema dual para asegurar el cumplimiento de la obligación tributaria: retenciones y figuras personales. Veamos en detalle a las personales.

\paragraph{Representante}
La primera de estas figuras es la del representante. Los sujetos pasivos están obligados a nombrar a un representante con residencia en España, antes del vencimiento del plazo de declaración.

No obstante, esta obligación no tiene la misma fuerza sobre todos los sujetos. Solo es obligatoria si: (i) opera en España mediante establecimiento permanente; (ii) obtienen rendimientos sin establecimiento permanente por prestaciones de servicios, asistencia técnica, obras de instalación o montaje derivados de contratos de ingeniería o, en general, de explotación económicas; (iii) se trata de residentes en territorios con los que no existe intercambio efectivo de información y sean titulares de bienes, excluidos los valores negociados en mercados secundarios oficiales; (iv) se trata de entidades en régimen de atribución de rentas con presencia en territorio español (las del art. $38$); y, (v) si así lo requiere la Administración Tributaria.

El contribuyente o su representante están obligados a poner en conocimiento de la Administración Tributaria el nombramiento del representante en un plazo de dos meses a partir de la fecha del nombramiento.

\paragraph{Responsable: }
La segunda garantía personal es la figura del responsable. Cuando no existe establecimiento permanente, la Ley establece una obligación subsidiaria sobre el pagador de las rentas obtenidas por el sujeto pasivo del impuesto, así como sobre el depositario o gestor de sus bienes en el territorio. 

Es decir, al no mediar establecimiento permanente, el pagador tiene la obligación de retener sobre los pagos efectuadas el porcentaje que corresponde a la obligación tributaria. Por ello, en caso de que el suejto pasivo efectivamente incumpliese la obligación, la AEAT puede ir directamente contra el pagador de las rentas puesto que es el responsable de efectuar las retenciones, razón por la cual se establece la obligación solidaria. 

Si hubiese establecimiento permanente, también podría nombrar como responsable al mismo representante que hubiese nombrado. En definitiva, la Ley busca crear un incentivo sobre las personas residentes a ingresar la deuda tributaria. No obstante, esto no llega a ser necesario en muchos casos porque el importe de la retención, que el pagador está obligado a hacer, suele ser exactamente igual a la cuantía de la obligación.

\subsection{Esquema de Liquidación}
Una vez estudiadas las características y elementos definitorios, veamos su esquema de liquidación (muy básico). 

Lo primero que debemos tener en cuenta es que, como nos podemos imaginar por lo presentado hasta aquí, para liquidar el IRNR es decisivo saber si el no residente opera en España mediante un establecimiento permanente o no. 

Se entiende que existe establecimiento permanente cuando disponga de instalaciones o lugares de trabajo en los que realice toda o parte de su actividad, de forma continuada, o actúe en España por medio de un agente autorizado. Por ejemplo se considera establecimiento permanente una sede de dirección, sucursal, oficina, fábrica, taller, almacén, entre otros. Una filial no se considera establecimiento permanente porgue tiene personas jurídicas en España y por tanto es residente. y tributaría por el IS.


Si lo hay, el impuesto se parece al Impuesto sobre Sociedades; si no, al IRPF. Ahora bien, ¿por qué el legislador distingue según exista o no establecimiento permanente?

\subsubsection{¿Por qué el legislador distingue estas dos situaciones?}
La respuesta tiene cinco razones, y todas conducen a la misma idea de fondo: la distinción resuelve un conflicto entre medir bien la capacidad económica y poder cobrar el impuesto.

En primer lugar, el impuesto sobre la renta pretende gravar la renta neta: lo que se ingresa menos lo que cuesta obtenerlo. Un establecimiento permanente es una presencia estable, con contabilidad, instalaciones y empleados en España; sus ingresos y gastos se producen aquí y pueden verificarse, igual que los de una sociedad española. Por eso puede gravarse su beneficio neto con las reglas del Impuesto sobre Sociedades. Sin establecimiento permanente, en cambio, España solo ve un pago aislado: un dividendo, un alquiler, una factura por un servicio. Los costes de obtener esa renta se producen normalmente fuera de España y la Administración española no puede comprobarlos. La solución práctica es gravar el importe íntegro, sin gastos, a un tipo relativamente bajo que actúa como aproximación a lo que pagaría la renta neta. Por eso los tipos del régimen sin establecimiento son relativamente bajos; al no admitirse deducciones, aunque sean necesarias para obtener la renta, se compensa con los tipos. Por tanto, esta distinción persigue \textbf{ajustar la medición de la renta neta} a las posibilidades reales de la AEAT.

En segundo lugar, cuando hay establecimiento permanente, el control es más fácil y existe la figura del representante. Cuando no lo hay, el control es muy difícil, y la forma de garantizar el cobro es que el impuesto se devengue en cada operación y que el pagador tenga la obligación de retener. Es exactamente la misma lógica de gravar donde el Estado puede observar que vimos en la historia de la imposición indirecta, desde los peajes romanos hasta el crédito de factura del IVA ([\authorfont{Ver Tema 4.B.12}]). Un no residente sin presencia en España no tiene bienes aquí que embargar ni contabilidad que inspeccionar; el único punto observable y controlable es el pagador residente, que sí está en España. Por eso el sistema traslada al pagador la obligación de retener, y fija esa retención de modo que coincida con la cuota final, el impuesto queda cobrado en el mismo momento del pago. Por tanto, persigue también un objetivo pragmático: \textbf{poder cobrar el impuesto}. 

Por último, existen dos razones que emanan directamente del articulado de los convenios suscritos por España. En primer lugar, el art. $7$ del Modelo OCDE, como veremos, solo permite al país de la fuente gravar los beneficios empresariales de un no residente si este tiene allí un establecimiento permanente. La ley interna española se alinea con ese mismo umbral: el concepto de establecimiento permanente es, a la vez, la puerta de entrada del IRNR al régimen de beneficio neto y el requisito que los convenios exigen para que España pueda gravar esos beneficios. La intuición económica del régimen del establecimiento permanente es no discriminar a los no residentes frente a las empresas residentes: no tendría sentido que una sucursal (sin personalidad propia) sujeta al IRNR tuviera una obligación distinta de la de una filial (con personalidad propia) sujeta al Impuesto sobre Sociedades. Este principio no es solo una preferencia del legislador español, sino que lo exige el art. $24.3$ del Modelo OCDE, recogido en los convenios españoles, que prohíbe gravar a un establecimiento permanente de forma menos favorable que a las empresas residentes que realicen la misma actividad. Y lo exige el Derecho de la Unión a través de la libertad de establecimiento. La consecuencia económica es la neutralidad en la forma jurídica, la multinacional debe elegir entre sucursal y filial por razones de negocio, no fiscales. Por tanto, la distinción busca \textbf{no discriminar al inversor extranjero}.

Por último, una filial española paga el Impuesto sobre Sociedades y, cuando reparte dividendos a su matriz extranjera, esos dividendos soportan además una retención en España. Una sucursal paga el Impuesto sobre Sociedades, pero cuando transfiere su beneficio a la casa central no hay ningún dividendo que gravar, porque jurídicamente es la misma persona. Para igualar ambos casos, el art. $19.2$ del texto refundido establece un gravamen complementario del $19\%$ sobre las rentas que el establecimiento transfiere al extranjero. No se aplica cuando la casa central reside en otro Estado de la Unión, salvo que sea una jurisdicción no cooperativa, porque dentro de la Unión la Directiva matriz-filial ya exime los dividendos; ni cuando un convenio lo impide.

\paragraph{¿Qué pasaría si el legislador no distinguiera?}
Podemos preguntarnos, para aclarar definitivamente esta distinción, ¿qué sucederia si el legislador no distinguiera? Podrían darse los siguientes casos polares. 

Por un lado, si todo se gravara como sin establecimiento, en bruto, una sucursal con altos costes pagaría impuestos sobre sus ingresos aunque tuviera pérdidas. De esta forma, la inversión productiva extranjera quedaría gravemente penalizada frente a la nacional. Sin entrar, evidentemente, en la vulneración que esto supondría de los convenios suscritos y el Derecho de la Unión.

Por otro lado, si todo se gravara como con establecimiento, en neto y por declaración anual, España debería esperar a que cada no residente sin presencia aquí presentara una declaración anual de sus rentas españolas y sus gastos, sin ningún medio para comprobarla ni para cobrarla si no la presenta. Por lo tanto, cabría esperar como resultado un fraude masivo, pues no habría incentivos a declarar algunas fuentes de renta.

La distinción es, por tanto, un compromiso entre precisión y eficacia. Precisión en la medida de la renta cuando hay presencia que permite medirla, eficacia en el cobro cuando no la hay.

\paragraph{Implicaciones: ¿tiene algún coste establecer esta frontera?}
Por último, una vez hemos visto las razones que llevan al legislador a trazar esta linea, y hemos visto brevemente que sucedería si no se establecier. ¿Es efectiva esta linea respecto a las implicaciones materiales que tiene? Dicho de otra forma, ¿genera costes trazar esta linea divisorio? La respuesta es que sí. 

\textcolor{red}{[\textbf{No queda claro si siempre existe incentivo, ¿siempre lo hayo solo cuando concurre alguna circunstancia particular? Al fin y al cabo, el IRNR sin establecimiento \textit{penaliza más} pues presenta una asimetría hacia ganancias: no permite la deducción de costes, ¿siempre hay este incentivo, parece contraintuitivo?}] En primer lugar, existe un incentivo a no tener establecimiento permanente. Si el régimen sin establecimiento resulta más favorable, el contribuyente tiene incentivo a organizar su actividad para no alcanzar ese umbral. Por ejemplo, fragmentando actividades o vendiendo a través de comisionistas. Llevado al extremo, es precisamente el problema que presenta la economía digital: una empresa puede obtener grandes beneficios en España sin establecimiento permanente, y por tanto sin que España pueda gravar su beneficio, lo que, como veremos, dio origen a un buen número de iniciativas en el marco multilateral y el Derecho de la Unión.}

En segundo lugar, existe un riesgo de discriminación en el régimen sin establecimiento. Gravar en bruto puede hacer que el no residente pague más que un residente por la misma renta. El TJUE lo corrigió en el caso Gerritse (C-234/01, 2003). Un batería neerlandés que actuaba en Berlín tributaba en Alemania sobre sus ingresos íntegros, sin poder deducir los gastos de la actuación, mientras los residentes sí podían hacerlo; por lo que el Tribunal concluyó que eso vulneraba la libre prestación de servicios. Esta jurisprudencia explica por qué la ley española permite a los residentes en la Unión y el EEE deducir los gastos directamente vinculados a sus rentas españolas. Además, la sentencia Schumacker (C-279/93, 1995) explica, en el mismo sentido, la opción de tributar por el IRPF que tienen los residentes en la Unión que obtienen en España la mayor parte de su renta. Si casi toda su renta procede del Estado de la fuente, su Estado de residencia no puede tener en cuenta sus circunstancias personales, y debe hacerlo el de la fuente. \textcolor{red}{\textbf{NOTA.- Este último caso no está bien explicado, hay que explicar un poco mejor qué es lo que se quiere decir, qué estaba sucediendo y por qué se alcanzó esa conclusión. Además, teniendo en cuenta el punto anterior, parece contradictorio: por un lado se dice que hay incentivos permanantes, por otro que se penaliza siempre, ¿no es contradictorio? No queda bien explicado...}}

\subsubsection{Establecimiento permanente}
Si existe establecimiento permanente, el sujeto pasivo tributa por la totalidad de la renta imputable al establecimiento, cualquiera que sea el lugar donde se haya obtenido o producido. Como es normal, puede existir más de un establecimiento permanente para el mismo sujeto, cuando realicen actividades claramente diferenciables y la gestión se lleve de forma separada. En tal caso, cada establecimiento permanente tributa por separado, sin posibilidad de compensar entre ellos rendimientos ni bases imponibles, ni trasladar deducciones o bonificaciones de uno a otro.

\paragraph{Base Imponible: Impuesto sobre Sociedades}
Como dijimos al principio, si existe establecimiento permanente el esquema de liquidación es similar al del IS. De esta forma, la base imponible del impuesto se determinará según la normativa del Impuesto sobre Sociedades, salvando las siguientes peculiaridades.

En primer lugar, en lo relativo a la \textbf{valoración de operaciones vinculadas}. \textcolor{red}{En caso de sujeción al IRNR, las operaciones del establecimiento permanente con su matriz o con otros establecimiento permanente de la misma, situados o no en España, o con otras sociedades vinculadas a la matriz, se valoran conforme a las reglas aplicables a las operaciones vinculadas.}\footnote{%
    Aunque en las disposiciones del IRNR se habla de casa central, hemos sustituido el término por matriz para facilitar la comprensión.
}

En segundo lugar, los gastos deducibles. El establecimiento permanente puede deducir la parte razonable de los gastos de dirección y gastos generales de administración en los que haya incurrido la matriz y se deban imputar al establecimiento permanente. No obstante, no serán deducibles: (i) los pagos del establecimiento permanente a la matriz o a otros establecimiento permanente en concepto de cánones, intereses, comisiones, etc. en contraprestación de servicios de asistencia técnica o por el uso o cesión de bienes o derechos; y/o, (ii) el coste de los capitales propios de la matriz afectos, directa o indirectamente, al establecimiento permanente en España. La razón económica de este trato diferenciado es, precisamente, evitar la erosión artifical de bases imponibles trasladando mayor beneficio a la matriz o a otros establecimiento permanente. En tercer lugar, la compensación por pérdidas tendrá un trato diferenciado. Por norma general, la compensación de bases imponibles negativas se realiza de acuerdo con la LIS, pero el establecimiento permanente no podrá compensar sus bases imponibles negativas con las bases imponibles positivas de otro establecimiento permanente. Por último, la integración de rentas por cese. En caso de cese de actividad o traslado al extranjero de activos se integra en la bases imponibles la diferencia entre el valor de mercado y el valor contable de los elementos.

Una vez calculada la base imponible del impuesto, la cuota íntegra se determina como el resultado de aplicar a la bases imponibles el tipo impositivo o tipo de gravamen señalado en la LIS. Además, se aplicarán sobre la cuota íntegra las bonificaciones y deducciones contempladas en el régimen general del IS, y las retenciones, ingresos a cuenta y pagos fraccionados ya efectuados conforme a la LIS. Tras estas operación, el resultante será la obligación tributaria del sujeto pasivo. 


De forma esquemática, podemos resumir las operaciones mediante el siguiente esquema:

\begin{table}[h]
    \centering
    \caption{Placeholder Caption}
    \label{tab:placeholder_label}
    \begin{tabular}{p{0.92\linewidth}}
        Ingresos del establecimiento permanente \\
        $-$ Gastos deducibles según la LIS \\
        \hspace{1.5em}$\cdot$ incluidos los gastos de dirección y generales de la casa central imputables en proporción razonable \\
        \hspace{1.5em}$\cdot$ excluidos: cánones, intereses y comisiones pagados a la casa central o a otros EP del grupo; coste de los capitales propios de la matriz \\
        $\pm$ Valoración de operaciones con la matriz y vinculadas a valor de mercado \\
        $\pm$ Integración de rentas por cese o traslado de activos al extranjero \\
        \hline
        $=$ Base imponible (sin compensar con otros EP del mismo contribuyente) \\
        $\times$ Tipo de gravamen del Impuesto sobre Sociedades (general: 25\%) \\
        \hline
        $=$ Cuota íntegra \\
        $-$ Bonificaciones y deducciones de la LIS \\
        \hline
        $=$ Cuota líquida \\
        $-$ Retenciones, ingresos a cuenta y pagos fraccionados \\
        \hline
        $=$ Cuota diferencial (a ingresar o devolver) \\
        \\
        Y, al transferir beneficios a la casa central: \\
        $+$ Gravamen complementario del 19\% sobre las rentas transferidas (salvo casa central en la UE/EEE no cooperativa, o convenio que lo impida) \\
    \end{tabular}
\end{table}

De la misma forma, imaginemos una sociedad con sede en un país no perteneciente a la Unión, sin convenio que impida el gravamen complementario, y que tiene una sucursal en Madrid. En el ejercicio en curso ha ingresado $5.000.000$ euros, ha soportado unos gastos operativos propios (personal, alquileres, suministros) de $3.500.000$ euros, ha pagado un canon a la casa central por uso de la marca de $400.000$ euros, ha imputado $200.000$ euros en concepto de dirección de la casa central a la sucursal, ha llevado a cabo una inversión de $25.000$ en I+D, íntegramente deducibles, y ha efectuado pagos fraccionados por valor de $150.000$ euros.

El primer baso para calcular la base imponible será tener en cuenta que el canon pagado a la casa central no es deducible, mientras que los gastos de dirección sí, por tanto: $5.000.000 - 3.500.000 - 200.000 = 1.300.000$ euros. Sin la regla que excluye el canon, la base habría sido de $900.000$ euros. De esta forma, la sucursal habría trasladado $400.000$ euros de beneficio a la casa central, justo el tipo de erosión que pretendemos evitar. 

¿Cuál será la cuota? En primer lugar, calculamos la cuota íntegra: $1.300.000 × 25\% = 325.000$ euros. Después, la cuota líquida: $325.000 - 25.000 = 300.000$ euros; y, por último, la cuota diferencial: $300.000 - 150.000 = 150.000$ euros a ingresar.

Por último, debemos ver si resulta aplicable el gravamen complementario. Para ello debemos preguntarnos, ¿ha transferido la sucursal todos sus beneficios a la filial? Si la sucursal lo ha hecho, entonces aplicamos el gravamen complementario como $1.300.000 - 300.000 = 1.000.000$ euros, soportando un 19\% adicional: $190.000$ euros. En caso de que los haya transferido, la obligación total será de: $340.000$ euros; en caso contrario, de $150.000$ euros.

Pensemos ahora, para recapitular todo lo expuesto antes, ¿qué pasaría si en vez de una sucursal estuviésemos ante una filial de la matriz? Si en lugar de una sucursal fuera una filial española con idénticas cifras, pagaría el mismo Impuesto sobre Sociedades, y al repartir un dividendo de $1.000.000$ de euros a su matriz soportaría una retención del $19\%$, también de $190.000$ euros. 

Por tanto, el gravamen complementario iguala así la carga de la sucursal y la de la filial. Se trata de la neutralidad en su forma jurídica. Si la casa central residiera en otro Estado de la Unión, ni la sucursal pagaría el gravamen complementario ni la filial la retención sobre dividendos, por efecto de la Directiva matriz-filial; y, por tanto, la igualdad se mantendría.

\subsubsection{Sin establecimiento permanente}
Ahora que conocemos el esquema básico de liquidación en caso de que disponga de una establecimiento permanente, ¿qué ocurre cuando esto no es así, cuando no dispone de uno?

En primer lugar, debemos tener claro que si las rentas se obtienen sin mediación de establecimiento permanente, el sujeto tributa de forma separada por cada devengo total o parcial de la renta sometida a gravamen, con total separación unas de otras, y sin que sea posible compensación alguna entre ellas.

Como es evidente, si no media un establecimiento permanente el control es mucho más complejo. Por ello, para evitar la posible ocultación de rentas se exige que estas se declaren devengo a devengo. Es decir se obliga al sujeto pasivo a contabilizar de forma separada cada renta, en el momento de cada devengo, y sin la posibilidad de compensar entre éstas. Como vemos, la diferencia con el caso anterior es muy sustancial, pues permite el cómputo conjunto, la compensación por establecimiento y el devengo en el último día del ejercicio fiscal. 

No obstante, la diferencia crucial, que en realidad es consecuencia de las diferencias ya enunciadas, es más de carácter práctico que formal: \textbf{el pagador está obligado a retener una cuantía igual al impuesto}. 

\paragraph{Base Imponible: Impuesto sobre la Renta de las Personas Físicas}
La base imponible, en este segundo supuesto, se aproxima mucho más a la forma general prevista en la LIRPF. Concretamente, en la determinación de la renta gravada se tienen en cuenta las siguientes reglas:
\begin{la}
    \item Con carácter general, la bases imponibles correspondiente a rendimientos está constituida por su importe íntegro, determinado conforme a las normas del IRPF;
    \item Como excepción, en el caso de prestaciones de servicios, asistencia técnica, obras de instalación o montaje derivadas de contratos de ingeniería y, en general, de actividades o explotaciones económicas, la base imponible es igual a la diferencia entre los ingresos íntegros y los gastos de personal, aprovisionamiento de materiales incorporados a las obras, trabajos y suministros;
    \item Tratándose de ganancias patrimoniales, la base imponible se calcula aplicando a cada alteración patrimonial las normas previstas en la ley del IRPF;
    \item En el caso de personas físicas no residentes, la renta imputada de los bienes inmuebles urbanos situados en territorio español se determina con arreglo a lo dispuesto en la LIRPF;
    \item En el caso de contribuyentes residentes en otro Estado de la Unión, para la determinación de la base imponible se podrán deducir los gastos previstos en la LIRPF siempre que se acredite que están relacionados directamente con los rendimientos obtenidos en España. Para las ganancias patrimoniales obtenidas también se aplicarán las normas del IRPF. En el caso de entidades serán deducibles los gastos indicados en la LIS.
\end{la}

Por tanto, vemos como, en general, los rendimientos se consideran por su importe íntegro y no hay gastos deducibles ni reducciones.\footnote{%
    Salvo algunas excepciones. Por ejemplo los rendimientos de actividades económicas permiten deducir determinados gastos.
}

Una vez calculada la base imponible, la cuota íntegra será el resultado de aplicar a la bases imponibles: (i) el tipo general, del $24\%$; (ii) el tipo reducido del $19\%$, para residentes en otros Estados de la Unión; o, (iii) otros tipos específicos en función de la fuente de renta, por ejemplo del $19\%$ para dividendos y ganancias patrimoniales, gravando la misma fuente a un tipo similar al del IRPF.

Como vemos, todos los tipos de gravamen son relativamente bajos, precisamente porque hay que tener presente que, de forma general, no se admiten deducciones/reducciones, aunque sean necesarios para obtener las rentas. Una vez tengamos la cuota íntegra, debemos sustraer las deducciones por donativos del IRPF, sin bonificaciones, y las retenciones o ingresos a cuenta que se hayan materializado previamente, no se tienen que realizar pagos fraccionados, ya que los rendimientos se tributan devengo a devengo. El resultado final será la cuantía de la obligación tributaria. En muchos casos, debido a la obligación de retener y el esquema proporcional del impuesto, el sujeto pasivo tendrá una obligación de cuantía nula, y solamente tendrá que presentar declaración.

De forma esquemática, podemos resumir las operaciones mediante el siguiente esquema. Para CADA devengo, por separado y sin compensar con otros:

  \begin{table}[h]
    \caption{Placeholder Caption}
    \label{tab:placeholder_label}
    \centering
    \begin{tabular}{p{0.97\linewidth}}
        \hline
        \textbf{Cálculo (esquema)} \\
        \hline
        Importe íntegro de la renta \\
        $-$ Gastos deducibles, solo en los casos previstos: \\
        \hspace{1.5em}$\cdot$ prestaciones de servicios, asistencia técnica, montaje y actividades económicas: gastos de personal, materiales y suministros \\
        \hspace{1.5em}$\cdot$ residentes en la UE/EEE: gastos del IRPF (personas físicas) o de la LIS (entidades) directamente vinculados a la renta \\
        \hspace{1.5em}$\cdot$ ganancias patrimoniales: valor de transmisión $-$ valor de adquisición, según las normas del IRPF \\
        \hline
        $=$ Base imponible de esa renta \\
        $\times$ Tipo de gravamen: \\
        \hspace{1.5em}$\cdot$ general: 24\% \\
        \hspace{1.5em}$\cdot$ residentes en la UE/EEE: 19\% \\
        \hspace{1.5em}$\cdot$ dividendos, intereses y ganancias patrimoniales: 19\% \\
        \hspace{1.5em}$\cdot$ o el tipo reducido que fije el convenio aplicable \\
        \hline
        $=$ Cuota íntegra \\
        $-$ Deducciones por donativos (según el IRPF) \\
        $-$ Retenciones practicadas por el pagador \\
        \hline
        $=$ Cuota diferencial (normalmente cero: la retención equivale a la cuota) \\
        \hline
    \end{tabular}
\end{table}

Supongamos, ahora con un ejemplo númerico, que una persona residente en Estados Unidos, obtiene rentas en España sin establecimiento permanente. El total de sus rentas española en este periodo impositivo ha sido de cuatro.

Por ejemplo, una de dividendos de una sociedad española por $10.000$ euros. El tipo del IRNR es del $19\%$, es decir, $1.900$ euros. Además, el convenio hispano-estadounidense limita la tributación en la fuente de estos dividendos al $15\%$: $1.500$ euros.
Si la sociedad pagadora retuvo el $19\%$ porque el perceptor no acreditó su residencia, este puede pedir la devolución de $400$ euros. Si la acreditó con un certificado de residencia fiscal, la sociedad retiene directamente el $15\%$.

La segunda de sus rentas proviene de un alquiler de un piso en Madrid por $12.000$ euros al año, con $3.000$ euros de gastos (IBI, comunidad, reparaciones). Como no reside en la Unión, no puede deducir gastos. Por tanto, tributa sobre $12.000$ euros al $24\%$, es decir, $2.880$ euros. En contraste, si en vez de residir en Estados Unidos residiera en Francia, podría deducir los gastos y tributaría al $19\%$ sobre $9.000$ euros, $1.710$ euros. La diferencia es la consecuencia práctica de la jurisprudencia Gerritse y de la libertad de circulación, que protege al residente en la Unión pero no al estadounidense.

Supongamos que su tercera fuente es por un servicio de consultoría prestado en España por $20.000$ euros, con $8.000$ euros de gastos de personal y materiales. Al tratarse de una prestación de servicios, sí puede deducir esos gastos. Por tanto, su base imponible sería de $12.000$ euros al 24\%, es decir, una obligación de $2.880$ euros.

Por últmio, imaginemos que el sujeto sufre unas pérdidas patrimoniales, por la venta de acciones de una sociedad española, por valor de $5.000$ euros. Por mucho que haya sufrido una pérdida patrimonial, como cada renta tributa por separado, no puede compensarse con los dividendos, el alquiler ni la consultoría. Por tanto, esos $5.000$ euros de pérdida no reducen ninguna otra cuota.

En total, la carga total que experimenta en España es de $1.500 + 2.880 + 2.880 = 7.260$ euros, sobre unas rentas netas reales de $10.000 + 9.000 + 12.000 - 5.000 = 26.000$ euros. El tipo efectivo al que tributa es del $27,9\%$. Como sabemos ([\authorfont{Ver Tema 4.B.16}]), un residente en España con esas mismas rentas podría deducir los gastos del alquiler y compensar la pérdida con otras ganancias patrimoniales. La diferencia, por tanto, es el coste de la regla de cada renta por separado, que protege la recaudación a cambio de una medida menos precisa de la capacidad económica.









\clearpage
\section{Dimensión bilateral: Doble imposición}
\puntosclave{
    Mediante los CDIs se busca solución a los problemas de doble imposición y a los problemas de soberanía tributaria.
    
    El objetivo esencial de un CDI es la distribución de los rendimientos gravables de los ciudadanos de los Estados contratantes entre el Estado de la fuente (aquel en el que se origina la renta) y el Estado de residencia (Estado en el que reside el perceptor de la renta distinto del de generación).
    
    Los CDI suscritos por España están basados en el modelo de la OCDE. Éste utiliza de forma general el criterio de residencia y permite dos métodos para corregir la doble imposición: exención e imputación. España ha optado de forma general por la imputación.
}

Hasta aquí hemos visto la solución española para gravar las rentas que obtienen en España los no residentes: el IRNR. En contraposición, el IRPF y el Impuesto sobre Sociedades gravan a los residentes por la totalidad de sus rentas mundiales. Cuando dos países aplican a la vez estos criterios, una misma renta puede quedar gravada en el país donde se genera, por obligación real, y en el país de residencia del perceptor, por obligación personal.

Imaginemos una sociedad residente en España que tiene una filial en Alemania y repatria, anualmente, un dividendo de $1.000$ euros. Alemania gravará ese dividendo en la fuente, porque se genera allí. Por ejemplo, supongamos que aplica una retención del $25\%$ (250 euros). Por otro lado, España gravará a la sociedad española por su renta mundial, que incluye ese dividendo. Por ejemplo, supongamos que tributa al $25\%$, otros 250 euros. En total, el resultado es que sobre $1.000$ euros de renta, el contribuyente paga $500$ en impuestos. Sin embargo, una inversión idéntica en una filial española solo habría pagado $250$, produciéndose así una distorsión por motivos fiscales que podría afectar enormemente al comercio internacional de bienes y servicios y a los flujos internacionales de capital.

Eso es la doble imposición: el mismo hecho imponible gravado por figuras tributarias de naturaleza análoga en dos jurisdicciones distintas. Su efecto económico está claro: la inversión transfronteriza soporta una carga que no soporta la doméstica, de modo que el sistema fiscal penaliza la internacionalización con independencia de su rentabilidad real. 

Pasamos ahora a estudiar como se trata este fenómeno, que requiere como podemos intuir la colaboración entre diferentes jursidicciones fiscales, por lo que su respuesta clásica son los convenios bilaterales.

\subsection{Convenios de doble imposición (CDI): ¿qué son, de dónde vienen y cómo se organizan?}
Los convenios para evitar la doble imposición (CDI) son tratados internacionales suscritos entre dos Estados, y excepcionalmente entre varios, que determinan en qué país puede gravarse cada tipo de renta y, cuando pueden gravarla ambos, con qué límites. Dicho de otro modo, distribuyen la potestad de gravar entre el Estado en cuya jurisdicción está la fuente de la renta y el Estado donde reside quien la percibe.\footnote{%
    Los más habituales son bilaterales y se refieren sobre todo a la imposición directa sobre la renta y el patrimonio. Existen también algunos multilaterales, como el Convenio Nórdico entre Dinamarca, Finlandia, Islandia, Noruega y Suecia, de 1983.
}

Conviene precisar una idea que a menudo se entiende al revés. Un convenio no crea impuestos, solo limita los que ya existen en la legislación interna de cada Estado. Si la ley española no grava una renta, el convenio no puede hacer que se grave; si la grava, el convenio puede impedir que se aplique o limitar su cuantía. Por eso, para saber \textbf{cuánto paga un no residente en España hay que leer siempre dos normas a la vez}: el IRNR, que dice qué grava España, y el convenio aplicable, que dice hasta dónde puede hacerlo.

\subsubsection{Genealogía: ¿cómo se han desarrollado a lo largo de la historia?}
Antes de que los Estados se pusieran de acuerdo para eliminar la doble imposición, firmaron acuerdos de cooperación administrativa en materia tributaria. El primero conocido lo firmaron Francia y Bélgica en 1843: sus funcionarios intercambiaban documentos e información útiles para recaudar los impuestos de ambos países. Es un dato revelador, pues la cooperación para recaudar precede a la cooperación para no gravar dos veces, en contra de la idea extendida de que los convenios nacieron sobre todo para proteger al contribuyente.

El primer convenio de esta naturaleza concreta es del Reino Unido, de 1872, y se refería a los impuestos sobre sucesiones de los ciudadanos británicos residentes en el cantón suizo de Vaud, ya que muchos ingleses se instalaban en Lausana o Montreux, y al morir pagaban el impuesto sucesorio suizo y el británico. Sin embargo, el primer convenio general dedicado expresamente a evitar la doble imposición se firmó entre el Imperio austrohúngaro y Prusia el 21 de junio de 1899. Su contexto es el de la aparición de los impuestos modernos sobre la renta en Europa central a finales del siglo XIX, pues en cuanto dos Estados vecinos empiezan a gravar la renta de forma general, la doble imposición se convierte en un problema real para quien vive en uno y trabaja o invierte en el otro.

Tras la Primera Guerra Mundial, el aumento de los tipos para financiar la reconstrucción hizo la doble imposición mucho más gravosa. La Conferencia Financiera Internacional de Bruselas de 1920 pidió actuar, y en 1921 el Comité Financiero de la Sociedad de Naciones encargó a cuatro economistas de Italia, Países Bajos, Reino Unido y Estados Unidos un estudio sobre los aspectos económicos de la doble imposición internacional. EINAUDI, BRUINS, SELIGMAN y STAMP publicaron su informe en 1923, que desarrolló por primera vez la idea de \textbf{pertenencia económica} (\textit{economic allegiance}). 

Una renta pertenece en parte al país donde se genera y en parte al país donde reside quien la percibe, y la cuestión es cómo repartirla. De ahí nace el llamado compromiso de los años veinte, que sigue vertebrando el sistema actual. Según ese compromiso, \textbf{los beneficios empresariales y profesionales deben tributar principalmente en la fuente, donde se genera la actividad, mientras que las rentas de inversión, como dividendos, intereses y cánones, deben tributar principalmente en el país de residencia}. No obstante, conviene no idealizar ese origen. La historiografía reciente sostiene que el informe reflejaba desacuerdos de fondo, en la teoría y en la práctica, entre los dos grandes exportadores de capital de la época, Estados Unidos y Reino Unido y que los trabajos posteriores de la Sociedad de Naciones cambiaron sustancialmente el enfoque. Sin embargo, la lección es útil: el reparto clásico entre fuente y residencia no es una verdad técnica neutral, sino el resultado de una negociación entre los países que entonces exportaban capital. Como ya podremos haber intuido el reparto guarda en realidad una estrecha relación con la política, más que con la economía. ¿A quién pertenece la riqueza de un país? ¿Por qué gravar las fuentes materiales de renta en origen, y las fuentes nominales de renta en destino? Contestar a esta pregunta nos ayudará a entender mejor el tema.

Entre 1923 y 1927, un grupo de altos funcionarios tributarios elaboró los primeros proyectos de convenio bilateral para evitar la doble imposición en los impuestos directos, en los de sucesiones y en materia de asistencia administrativa. Esos trabajos culminaron en los modelos de 1928. Durante la Segunda Guerra Mundial se aprobaron dos modelos más con orientaciones opuestas: el de México (1943), más favorable a la tributación en la fuente, elaborado con fuerte participación latinoamericana, y el de Londres (1946), más favorable a la residencia.\footnote{%
    El mismo conflicto de fondo que hoy enfrenta al Modelo OCDE con el de Naciones Unidas.
}

Los modelos de México y Londres dieron cierta uniformidad a los convenios de posguerra, pero seguía habiendo muchas diferencias entre ellos, y la incertidumbre resultante para las empresas llevó en 1956 a crear el Comité Fiscal de la Organización Europea de Cooperación Económica (OECE), con el encargo de redactar un modelo de convenio bilateral. Trabajo que dio lugar al primer Modelo de la OCDE en 1963, revisado en 1977 y transformado en 1992 con un enfoque dinámico, y la intención de actualizarlo de forma periódica.

Desde entonces, la arquitectura fiscal internacional ha descansado en convenios bilaterales. Hoy en día existen más de $3.000$, y la mayoría sigue el Modelo de la OCDE o el de Naciones Unidas. Precisamente, esta densa red descentralizada es la mayor virtud del sistema, porque permite adaptar cada convenio a la relación económica concreta entre dos países, y a la vez su mayor debilidad, porque es imposible cambiarla rápidamente. Por eso, las reformas recientes, como veremos, han necesitado instrumentos multilaterales para modificar de golpe miles de convenios.

\paragraph{Naturaleza: ¿cuál es su naturaleza?}
Como cualquier tratado, se interpretan conforme a la Convención de Viena sobre el Derecho de los Tratados de 1969. Su regla general es interpretar los términos de buena fe, en su contexto y a la luz del objeto y fin del tratado. Esto explica la importancia del cambio de título y preámbulo del Modelo en 2017 que veremos más adelante.

Por otro lado, debemos tener claro que los convenios cumplen cuatro funciones: (i) eliminar la doble imposición y, desde la última década, también prevenir la doble no imposición; (ii) resolver el problema de la soberanía tributaria, pues firmar un convenio implica repartir de mutuo acuerdo el poder de gravar entre dos Estados que, sin él, lo reclamarían ambos; (iii) dar seguridad jurídica a los contribuyentes, que saben de antemano cómo tributará su inversión; y, (iv) estimular la inversión extranjera y las relaciones económicas entre los dos países, al reducir la carga fiscal y la incertidumbre. Se podría incluso añadir una quinta: la \textbf{cooperación administrativa}. El intercambio de información y la asistencia en la recaudación entre Administraciones, que se ha convertido en una de las piezas centrales del sistema.

Podríamos preguntarnos si cumplen sus funciones declaradas. Por ejemplo, ¿estimulan realmente la inversión? Esta cuarta función suele darse por supuesta, pero no está tan clara en la evidencia. Los estudios empíricos sobre el efecto de los convenios en la inversión directa extranjera ofrecen resultados contradictorios. Algunos encuentran un efecto positivo, otros no encuentran ninguno significativo. Para los países en desarrollo, en particular, varios autores sostienen que firmar convenios con países ricos, que casi siempre limitan la tributación en la fuente, les hace perder recaudación sin que la inversión aumente de forma demostrable. 

\subsubsection{Modelo OCDE: estructura, lógica de reparto y el contrapunto del Modelo ONU}
La mayoría de convenios existentes tienen un contenido muy parecido porque, como dijimos, se elaboran a partir de un modelo. 

Eso tiene varias ventajas. Sirve como punto de partida y agiliza la negociación, da seguridad jurídica al usar un lenguaje consensuado, y facilita la aplicación e interpretación de sus disposiciones. Además, el modelo suele ir acompañado de unos Comentarios, elaborados por los países miembros de la OCDE, que son muy útiles para aplicar los convenios y resolver controversias; y, a pesar de que éstos no son jurídicamente vinculantes, los tribunales de muchos países, incluidos los españoles, los utilizan como criterio interpretativo de referencia. Por último, cada país puede además formular reservas (no aceptar un artículo) y observaciones (discrepar de una interpretación de los Comentarios), lo que permite conocer de antemano su posición negociadora.

El modelo nació como Draft Double Taxation Convention on Income and Capital en 1963 y adoptó en 1977 su denominación actual, Model Tax Convention on Income and on Capital. En 1992 se publicó en un formato de hojas intercambiables que reconocía que su revisión era un proceso continuo, no la culminación de un trabajo cerrado, y desde entonces se ha actualizado de forma periódica en diez ocasiones: 1994, 1995, 1997, 2000, 2003, 2005, 2008, 2010, 2014 y 2017.

Dos actualizaciones merecen atención especial. En primer lugar, en 2017 se dio lo que se conoce como \textbf{giro antiabuso}. Se modificaron el título y el preámbulo para dejar claro que los convenios no solo buscan eliminar la doble imposición, sino también prevenir la elusión y la evasión fiscal, sin crear oportunidades de no imposición o de imposición reducida, incluido el abuso de los convenios por residentes de terceros Estados. Además, se añadió el art. $29$, que regula quién tiene derecho a los beneficios del convenio, para impedir lo que se conoce como \textit{treaty shopping}. En segundo lugar, en 2025 se llevó a cabo la primera revisión completa desde 2017. Se publicó el 19 de noviembre de 2025, aprobada por el Consejo de la OCDE ese mismo mes, y añadió un nuevo apartado al art. $25$ sobre el procedimiento amistoso, revisando también los Comentarios de varios artículos. Los cambios más detallados se refieren al art. $5$. Se introduce un marco para decidir cuándo el domicilio de un trabajador en remoto puede constituir un establecimiento permanente de su empresa, con una referencia orientativa del $50\%$ del tiempo total de trabajo, centrada en un uso regular y sustancial y no intermitente. Incluye también una disposición opcional para las actividades extractivas, que permite pactar bilateralmente un umbral temporal más bajo para que exista establecimiento permanente. Como podemos intuir, es la respuesta del Modelo a la extensión del teletrabajo tras la pandemia, un fenómeno que el concepto clásico de lugar fijo de negocios no contemplaba.

\paragraph{Estructura: ¿cómo se organiza?}
El articulado se organiza en siete capítulos:

\begin{table}[h]
    \caption{Placeholder Caption}
    \label{tab:placeholder_label}
    \centering
    \begin{tabular}{lll}
        \hline
        Capítulo & Artículos & Contenido \\
        \hline
        I. Ámbito & 1-2 & Personas e impuestos a los que se aplica \\
        II. Definiciones & 3-5 & Definiciones generales, residencia (con reglas de desempate) y establecimiento permanente \\
        III. Imposición de las rentas & 6-21 & Reglas de reparto para cada tipo de renta \\
        IV. Imposición del patrimonio & 22 & Reglas de reparto del patrimonio \\
        V. Métodos & 23 A y 23 B & Exención e imputación \\
        VI. Disposiciones especiales & 24-29 & No discriminación, procedimiento amistoso, intercambio de información, asistencia en la recaudación, diplomáticos y derecho a los beneficios \\
        VII. Disposiciones finales & 30-32 & Extensión territorial, entrada en vigor y denuncia \\
        \hline
    \end{tabular}
\end{table}

La lógica del reparto sigue tres tipos de reglas. El modelo utiliza como regla general el criterio de residencia y como excepción el de la fuente. Para entender cómo lo hace, conviene clasificar sus reglas de reparto en tres grupos:
\begin{la}
    \item \textbf{Tributación exclusiva en la residencia}. Solo puede gravar el país de residencia. Es el caso general de los beneficios empresariales obtenidos sin establecimiento permanente (art. $7$), las ganancias por venta de acciones en general (art. $13$) o las pensiones privadas (art. $18$). Aquí no hay doble imposición posible, porque el país de la fuente renuncia por completo a gravar;
    \item \textbf{Tributación compartida con límite en la fuente}. Pueden gravar ambos, pero el país de la fuente solo hasta un tipo máximo. Es el caso de los dividendos (art. $10$), con un límite del $5\%$ para participaciones significativas entre sociedades y del $15\%$ en los demás casos, y de los intereses (art. $11$), con un límite del $10\%$. Los cánones (art. $12$) tributan en el Modelo OCDE exclusivamente en la residencia, aunque muchos convenios reales, incluidos varios españoles, permiten una retención limitada en la fuente; y,
    \item \textbf{Tributación preferente en la fuente}. Puede gravar sin límite el país de la fuente, y el de residencia debe eliminar la doble imposición. Es el caso de las rentas de bienes inmuebles (art. $6$), los beneficios atribuibles a un establecimiento permanente (art. $7$) o las rentas del trabajo por actividad desarrollada en el otro país (art. $15$).
\end{la}

En los grupos (b) y (c) es donde puede producirse doble imposición, y donde entran en juego los métodos de exención e imputación del art. $23$ que estudiaremos más adelante. 

\paragraph{Principios de aplicación: ¿qué exige cualquier convenio?}
El modelo se apoya en cuatro principios:
\begin{la}
    \item \textbf{No discriminación} (art. $24$). Este principio prohíbe tratar peor a los nacionales o empresas del otro Estado que a los propios en la misma situación;
    \item \textbf{Procedimiento amistoso} (art. $25$). Si un contribuyente considera que se le aplica una medida contraria al convenio, puede plantearlo ante su Estado de residencia, y ambos Estados deben intentar resolverlo de mutuo acuerdo; si no lo consiguen, puede acudirse al arbitraje cuando el convenio lo prevé;
    \item \textbf{Intercambio de información} (art. $26$). Como condición previa al buen desarrollo de la relación prevista en el convenio, está el intercambio de información que genere confianza mutua. Por tanto, en virtud de este principio, los Estados deben intercambiar información para aplicar correctamente el convenio y prevenir el fraude; y,
    \item \textbf{Primacía}. Los tratados internacionales forman parte del ordenamiento interno y prevalecen sobre la ley nacional. Por ejemplo, en España, el art. $96$ de la Constitución dispone que sus disposiciones solo pueden derogarse, modificarse o suspenderse en la forma prevista en los propios tratados o conforme a las normas generales del Derecho internacional. El art. $4$ del texto refundido del IRNR lo recoge de forma expresa para este impuesto. Por eso, si un no residente obtiene rentas en España y existe convenio con su país de residencia, se aplica el convenio antes que el IRNR.\footnote{%
        Esta primacía tiene una consecuencia importante. En otros ordenamientos, como el estadounidense, una ley posterior puede prevalecer sobre un tratado anterior, práctica conocida como \textit{treaty override}. En España eso no es posible, ya que el legislador no puede desactivar un convenio aprobando una ley contraria, y si quiere cambiar su contenido tiene que renegociarlo o denunciarlo.
    }
\end{la}

Algunos convenios, sobre todo los suscritos con países menos desarrollados, incorporan aspectos del Modelo de Naciones Unidas en lugar del de la OCDE. Sin ser exhaustivos en su análisis, sí debemos tener en cuenta algunas diferencias. En particular, el Modelo ONU (1980) tiene la misma estructura que el de la OCDE, pero reparte los derechos de gravamen de forma más favorable al país de la fuente, que suele ser el país en desarrollo. ¿Cómo hace esto? Las diferencias esenciales que hay que conocer es que:
\begin{lnum}
    \item \textbf{Establecimiento permanente}. Define el establecimiento permanente de forma más amplia. Por ejemplo, las obras de construcción crean establecimiento permanente a partir de seis meses en lugar de doce;
    \item \textbf{Cánones}. Permite al país de la fuente gravar los cánones y, desde 2017, las remuneraciones por servicios técnicos;
    \item \textbf{Servicios digitales}. Incorpora desde 2021 el art. 12B sobre servicios digitales automatizados.
\end{lnum}

La coexistencia de los dos modelos es la manifestación, en el terreno técnico, del conflicto de intereses entre exportadores e importadores de capital que atraviesa todo la historia de la fiscalidad internacional, desde la oposición entre los modelos de México y Londres en los años cuarenta hasta el pulso actual entre la OCDE y Naciones Unidas.

\subsection{Metodología: cómo coordinar las jurisdicciones fiscales}
Como hemos visto, el Modelo OCDE clasifica cada renta en tres grupos: tributación exclusiva en la residencia, compartida con límite en la fuente, y preferente en la fuente. 

En el primer grupo no hay problema, porque solo grava un país. En los otros dos, ambos Estados pueden gravar, y hace falta una regla que diga cómo evita el país de residencia que la renta pague dos veces. El convenio deja a los Estados la elección entre dos métodos que pasamos ahora a analiza: la exención o la imputación. Para explicarlos, recurriremos fundamentalmente a un ejemplo numérico, el mismo aplicando ambos métodos, con el fin de asegurar la comparabilidad y la valoración conjunta. Luego veremos el problema inverso: cómo los propios convenios pueden generar doble no imposición.

Imaginemos un contribuyente residente en España, que obtiene $50.000$ euros de renta en España y $20.000$ euros en el país $X$. Para simplificar, supondremos que la tarifa española tiene dos tramos: $20\%$ hasta $40.000$ euros y $40\%$ a partir de ahí.

Sin ningún mecanismo corrector, España gravará la renta mundial de $70.000$ euros (obligación personal), por tanto, la obligación tributaria será de $40.000 × 20\% + 30.000 × 40\% = 8.000 + 12.000 = 20.000$ euros. Como vemos, el tipo medio efectivo español es del $28,57\%$. No obstante, el país $X$, donde se genera la renta, también gravará en virtud de una obligación real. Supongamos, que grava los $20.000$ euros al $30\%$, $6.000$ euros.

\subsubsection{Doble imposición: ¿cómo eliminarla?}
En total, la carga que soporta el sujeto, sin ninguna corrección, es de $26.000$ euros. ¿Cómo podemos solucionar esta inequidad? (Por ejemplo, el mismo contribuyente con rentas solo en España, cuya obligación es de $20.000$ euros)


\paragraph{Método de exención}
Bbajo el método de exención, el \textbf{Estado de residencia renuncia a gravar las rentas que pueden gravarse en el Estado de la fuente}. Por tanto, la renta extranjera queda exenta en el país de residencia. 

Este tipo de exención puede adoptar \textbf{dos formas}. La primera es la forma de exención íntegra, la renta extranjera simplemente no se incluye en la base imponible.

De esta forma, España grava solo los $50.000$ euros españoles, resultando en una obligación de $40.000 × 20\% + 10.000 × 40\% = 12.000$ euros. Mientras que el país de la fuente gravará la renta al $30\%$, resultando en una carga total de $12.000 + 6.000 (en X) = 18.000$ euros.

Fíjate en que el contribuyente paga menos que si toda su renta fuera española ($20.000$ euros), aunque la renta de $X$ haya pagado un $30\%$. La razón es que la renta extranjera no solo deja de tributar en España: tampoco empuja a la renta española hacia el tramo alto de la tarifa. Por tanto, la exención íntegra configurada de esta forma rompe la progresividad del impuesto personal, como vimos ([\authorfont{Ver Tema 4.B.16}]).

Bajo la segunda modalidad, de exención con progresividad, la renta extranjera se integra en la base solo para determinar el tipo de gravamen, y ese tipo se aplica después únicamente a la renta española. Por tanto, se busca remediar la quiebra de progresividad anterior, intentando que la renta extranjera refleje la capacidad de pago, aunque no se grave para evitar la doble imposición.

De esta forma, se calcula el tipo medio sobre la renta mundial primero, $20.000 / 70.000 = 28,57\%$. España aplica ese tipo solo a la renta española, $50.000 × 28,57\% = 14.285,71$ euros, mientras que el país extranjero grava su renta al tipo proporcional. La carga total del sujeto será entonces de $14.285,71 + 6.000 = 20.285,71$ euros, ligeramente superior al primer caso.

Por tanto, la renta española tributa al mismo tipo al que tributaría si toda la renta fuera española, preservandp la progresividad y eliminando la doble imposición.

\paragraph{Método de imputación}
La alternativa es el método de imputación. Bajo esta modalidad, el Estado de residencia sí grava la renta extranjera, pero deduce de su impuesto el impuesto pagado en la fuente. Como antes, hay dos modalidades.

La primera es la imputación íntegra, donde se deduce todo el impuesto pagado en la fuente. Por ejemplo, España calcula su impuesto sobre la renta mundial, resultando en $20.000$ euros. Entonces deduce los $6.000$ pagados en $X$: $20.000 - 6.000 = 14.000$ euros de obligación tributaria con España y dejando una carga fiscal total de $14.000 + 6.000 = 20.000$ euros, exactamente lo mismo que si toda la renta fuera española.

El problema de esta modalidad aparece cuando el impuesto extranjero es más alto que el español. En este caso, España acabaría devolviendo parte del impuesto pagado en $X$ con cargo a la recaudación de sus rentas españolas. Como es normal, ningún país acepta subvencionar así a otros Estados, y por eso esta modalidad es casi inexistente.

La segunda modalidad es la imputación ordinaria. Con esta, el contribuyente puede deducir, como máximo, lo que habría pagado en España por esa misma renta. En este caso, el impuesto español que corresponde a la renta de $X$, sería de $20.000 × 28,57\% = 5.714,29$ euros, por lo que este constituye el límite de la deducción. Una vez calculada, se deduce el menor de los dos importes: el impuesto pagado en $X$ ($6.000$) o el límite ($5.714,29$) y, por tanto, se permite la deducción de $5.714,29$. El resultado es una obligación con el fisco español
Impuesto español de $20.000 - 5.714,29 = 14.285,71$ euros; y una carga total soportada de $14.285,71 + 6.000 = 20.285,71$ euros. Por tanto, los $285,71$ euros de exceso pagados en X no se recuperan.

Esta segunda modalidad, dentro del segundo método es el enfoque general de los convenios españoles. En resumen:

\begin{table}[h]
    \centering
    \caption{Placeholder Caption}
    \label{tab:placeholder_label}
    \begin{tabular}{lrrr}
        \hline
        Método & Impuesto en España & Impuesto en X & Carga total \\
        \hline
        Sin corrección & 20.000 & 6.000 & 26.000 \\
        Exención íntegra & 12.000 & 6.000 & 18.000 \\
        Exención con progresividad & 14.285,71 & 6.000 & 20.285,71 \\
        Imputación íntegra & 14.000 & 6.000 & 20.000 \\
        Imputación ordinaria & 14.285,71 & 6.000 & 20.285,71 \\
        \hline
    \end{tabular}
\end{table}

Cada uno de los enfoques da unos resultados distintos, excepto cuando el impuesto extranjero es superior al español. En estos casos, la exención con progresividad y la imputación ordinaria dan el mismo resultado, y el contribuyente acaba soportando el tipo más alto de los dos países.

Si, por el contrario, el país de la fuente grava menos, entonces el caso que decide la preferencia por cada método es fácil de ver. Suponiendo que el país X aplica un tipo del $10\%$ ($2.000$ euros).

Adoptando el enfoque de exención con progresividad, España cobrará $14.285,71$ y la carga total serán $16.285,71$. Por tanto el contribuyente se beneficia plenamente del tipo bajo de $X$. Sin embargo, adoptando el enfoque de imputación ordinaria, el límite serán igualmente $5.714,29$, pero solo se pagaron $2.000$ en $X$, así que la deducción límite es de $2.000$. En este caso, España cobra $18.000$ y la carga total será de $20.000$. El contribuyente pagará lo mismo que si hubiera invertido en España.

En definitiva, con la imputación la ventaja fiscal que $X$ quería ofrecer a los inversores desaparece. Lo que $X$ deja de cobrar lo cobra España. Por el contrario, con la exención, esa ventaja se mantiene íntegra. 

\paragraph{Valoración: ¿por qué no hay convergencia en el método?}
Esta diferencia explica realmente por qué no hay convergencia en la aplicación de los métodos y todo el debate en torno a la fiscalidad internacional.

Como podemos intuir, estos dos métodos son la traducción práctica de los dos criterios de neutralidad que vimos en el tema de imposición directa, al estudiar el sujeto del impuesto ([\authorfont{Ver tema 4.B.12}]).

La imputación busca la neutralidad en la exportación de capital. Esto es, que todos los residentes de un país soporten la misma carga con independencia de dónde inviertan, de modo que la fiscalidad no influye en la decisión de invertir dentro o fuera. En el ejemplo, el inversor español paga $20.000$ euros tanto si invierte en España como en $X$ con tipo bajo. Por otro lado, la exención busca la neutralidad en la importación de capital. Es decir, que todas las inversiones realizadas en un mismo país soporten la misma carga, con independencia de la residencia del inversor. En el ejemplo, el inversor español en $X$ compite en igualdad de condiciones fiscales con los inversores locales de $X$.

Ambas neutralidades solo pueden alcanzarse a la vez si todos los países aplican el mismo tipo. Por eso la elección entre métodos no es técnica, sino una elección entre dos objetivos incompatibles.

La exención suele preferirse desde la óptica de los países importadores de capital, porque respeta los incentivos que estos crean para atraer inversión, que la imputación anularía. La imputación suele preferirse desde la óptica de los países exportadores de capital, pues protege su recaudación, porque cuando el impuesto extranjero es bajo, recuperan la diferencia, y eliminan el incentivo fiscal a invertir fuera. No obstante, su gran inconveniente es precisamente ese, que anula los incentivos fiscales del país de la fuente.

Para resolver este inconveniente, muchos convenios entre un país desarrollado y uno en desarrollo incluyeron una cláusula especial, cláusula de crédito por impuesto no pagado (tax sparing). El país de residencia permite deducir no el impuesto efectivamente pagado en la fuente, sino el que se habría pagado sin el incentivo fiscal. En el ejemplo con $X$ al $10\%$, si el tipo normal de $X$ fuera del $30\%$, España permitiría deducir $6.000$ euros, aunque $X$ solo hubiera cobrado $2.000$, preservando así el incentivo de $X$. España incluyó este tipo de cláusulas en varios convenios con países iberoamericanos. No obstante, la propia OCDE se ha mostrado crítica con ellas desde finales de los años noventa, porque se prestan a abusos y tienden a mantener incentivos que, con el tiempo, ya no son necesarios.\footnote{%
    Con el Pilar 2 de la iniciativa BEPS, pierden buena parte de su sentido para los grandes grupos: si el incentivo reduce la tributación por debajo del $15\%$, la diferencia se cobra igualmente por otra vía.
}

Para terminar, conviene señalar que, en la práctica, la mayoría de los países, incluida España, no elige un método puro sino una combinación de ambos. Por un lado, suelen aplicar la exención a las rentas empresariales activas, como los beneficios de un establecimiento permanente o los dividendos de filiales con actividad real, para que sus empresas compitan en igualdad en el exterior. Y, por otro lado, suelen aplicr la imputación a las rentas pasivas y de las personas físicas, para evitar que las inversiones financieras se desplacen al exterior solo por razones fiscales.

\subsubsection{Elusión fiscal: ¿cómo solucionar el problema inverso?}
Lo primero que debemos señalar es que un convenio está diseñado para que una renta no tribute dos veces. Por ello, sus reglas pueden producir un resultado inesperado: que una renta no tribute en ninguno de los dos países. 

Esto no ocurre por un fallo de diseño, sino porque el convenio reparte la potestad de gravar dando por supuesto que el país al que se la asigna la ejercerá, y esa suposición no siempre se cumple.\footnote{%
    Conviene delimitar este apartado respecto del bloque 5: aquí estudiamos la doble no imposición que nace de la mecánica de los propios convenios; en el bloque 5 estudiamos el fenómeno global de traslado de beneficios, que va mucho más allá de los convenios.
}

Las cuatro causas comparten una misma raíz. El convenio reparte la potestad de gravar entre dos Estados, pero no puede obligar a ninguno a ejercerla. Toda la evolución reciente de los convenios, desde las cláusulas de sujeción hasta el art. $29$, consiste en condicionar la renuncia de un Estado a que el otro grave efectivamente. Es la misma lógica que, a escala global, veremos en la dimensión multilateral.

\paragraph{Exención: ¿qué pasa cuando la fuente no grava?} 
Este primer caso es el más claro y conviene verlo con números. Imaginemos una sociedad española que tiene un establecimiento permanente en el país $X$. El convenio asigna a $X$ la potestad de gravar sus beneficios, y España aplica el método de exención. El establecimiento obtiene $1.000.000$ de euros de beneficio, que España declara exentos, porque el convenio le dice que los grave $X$. Pero la ley interna de $X$ concede una exención temporal a las actividades nuevas, y $X$ no cobra nada. EL resultado, es que ese $1.000.000$ de euros \textbf{no tributan en ninguna parte}.

         Renta: 1.000.000 € (EP en X)
                    |
     ---------------------------------
     |                               |
 País X (fuente)             España (residencia)
 Convenio: puede gravar      Convenio: debe eximir
 Ley interna: exención       -> no grava
 -> no grava
     |                               |
     ---------- Tributación total: 0 € ----------

La respuesta técnica es la llamada cláusula de cambio de método (switch-over) o cláusula de sujeción. El propio articulado del Modelo OCDE permite al país de residencia dejar de aplicar la exención cuando el país de la fuente interpreta el convenio de forma que la renta queda exenta o sujeta a un límite en su territorio. Otras cláusulas van más allá. Por ejemplo, cuando el país de residencia solo exime si la renta ha estado efectivamente sujeta a tributación en la fuente; y, si no lo ha estado, cambia al método de imputación y la grava. En el ejemplo, España pasaría a gravar el $1.000.000$ de euros y, como $X$ no cobró nada, no habría nada que deducir.

\paragraph{Conflictos de calificación: ¿qué pasa cuando hay definiciones divergentes?}
Dentro de los convenios, el mismo problema aparece cuando los dos Estados califican de forma distinta una renta o una entidad, y cada uno aplica en consecuencia un artículo distinto del convenio. 

Si el país de la fuente considera que un pago es un interés, y lo trata como deducible y sujeto a un límite bajo, mientras el país de residencia considera que es un dividendo exento, la renta puede no tributar en ninguno de los dos. Los Comentarios del Modelo OCDE resuelven parte de estos conflictos obligando al país de residencia a aceptar, en ciertos casos, la calificación que hace el país de la fuente.

\paragraph{Doble no residencia: ¿y si no reside en ninguno de los dos paises?} 
Una persona o entidad puede no ser residente en ninguno de los dos países, porque cada uno aplica criterios de residencia distintos. Si además el convenio le atribuye la renta a su residencia, esta puede quedar sin gravar. 

\textcolor{red}{Es el mismo fenómeno que explotaba la estructura Double Irish que vimos en el 5.1.2: sociedades irlandesas que, por diferencias entre las leyes irlandesa y estadounidense, no eran residentes fiscales en ningún país.}

\paragraph{Abuso de convenios} 
El treaty shopping consiste en crear una sociedad residente en un país con el único fin de acceder a la red de convenios que este ha firmado: residentes de un tercer Estado se interponen en uno de los dos Estados contratantes para obtener ventajas que no obtendrían actuando directamente. La respuesta por parte de los tratados es doble:
\begin{la}
    \item Limitación de beneficios (LOB, por sus siglas en inglés). La cláusula de limitación de beneficios implica que solo se concedan los beneficios del convenio a quienes cumplen ciertos requisitos objetivos de vinculación con su país de residencia, como cotizar en bolsa allí o estar controlados por residentes;
    \item \textbf{Cláusula de propósito principal}. La cláusula del propósito marginal implica que se nieguen los beneficios del convenio cuando obtenerlos fue uno de los objetivos principales de una operación. \textcolor{red}{Es la que la mayoría de países ha incorporado a sus convenios a través del Instrumento Multilateral, como vimos en el 5.2.3.}
\end{la}

Ambas se recogen hoy en el art. $29$ del Modelo OCDE.



\subsection{España: la red de convenios, su gestión y su evaluación}
Actualmente, a junio de 2026, la red de convenios de esta índole en España se sitúan en los $99$ activos, con países como Alemania, Francia, Estados Unidos, China, México o Australia; y su evolución muestra un crecimiento sostenido.

En enero de 2013, Hacienda contaba $85$ convenios en vigor y otros $10$ en distintas fases de tramitación, un dato que presentaba como muestra de la internacionalización de la economía española. Pero, para diciembre de 2021, al ratificar el Instrumento Multilateral, ya había $94$ convenios en vigor, con contenido y redacción dispares, de los que $88$ podrían modificarse a través de ese instrumento a medida que los demás países lo ratificaran.

Muchos convenios en vigor se renegocian para actualizarlos. Por ejemplo, el convenio firmado en Bonn el 5 de diciembre de 1966 con Alemania, se sustituyó por uno nuevo, firmado el 3 de febrero de 2011, porque en algunos aspectos había quedado obsoleto y no reflejaba ni las relaciones económicas actuales ni los sucesivos cambios del Modelo OCDE. Otro ejemplo serían los convenios con el Reino Unido. El nuevo convenio sustituyó al firmado en Londres el 21 de octubre de 1975. Resolvió problemas que el anterior no podía abordar, como el tratamiento de los residentes británicos no domiciliados y de los trusts; estableció la tributación exclusiva en la residencia para los dividendos de participaciones mayoritarias, los intereses y los cánones; e introdujo una cláusula de arbitraje. Por otro lado, el caso de Estados Unidos muestra un ejemplo concreto del coste de no renegociar a tiempo. El convenio hispano-estadounidense de 1990 quedó desfasado cuando EEUU renegoció a comienzos de este siglo sus convenios con la mayoría de países de la Unión para reducir la tributación de dividendos, intereses y cánones: la inversión cruzada entre España y EEUU quedó penalizada frente a la de otros países europeos. El protocolo que lo actualizó se firmó en enero de 2013, pero no entró en vigor hasta 2019, por los retrasos en su ratificación por el Senado estadounidense.

La red española cubre a todos los socios de la Unión, a las grandes economías del mundo y, con \textbf{especial densidad, a Iberoamérica}. Esta última orientación no es casual. Como es normal, responde a la posición de las empresas españolas como grandes inversoras en la región desde los años noventa. Se combina además con un régimen interno que ha convertido a España en plataforma de inversión hacia Iberoamérica: el de las entidades de tenencia de valores extranjeros (ETVE), que permite a sociedades españolas cuya actividad es gestionar participaciones en filiales extranjeras recibir sus dividendos y plusvalías exentos. Sus defensores lo presentan como un instrumento legítimo para atraer sedes de grupos internacionales. No obstante, algunas organizaciones de la sociedad civil lo han criticado por facilitar que grupos extranjeros canalicen sus inversiones en Iberoamérica a través de España sin actividad real significativa aquí, aprovechando la red de convenios, el mismo tipo de uso que las cláusulas antiabuso intentan limitar.

Resulta llamativo, por otro lado, que los convenios entre Estados de la Unión sigan siendo bilaterales, y que no exista una norma común para evitar la doble imposición ni un modelo europeo propio impulsado por las instituciones de la Unión. Existen, no obstante, directivas que evitan la doble imposición en rentas concretas. Las principales son tres: (i) la Directiva matriz-filial (hoy Directiva 2011/96/UE), que elimina la doble imposición de los dividendos entre sociedades de distintos Estados miembros vinculadas por una participación significativa; (ii) la Directiva de intereses y cánones (Directiva 2003/49/CE), que suprime la retención en la fuente sobre estos pagos entre sociedades asociadas; y, (iii) los mecanismos de resolución de controversias, el Convenio de Arbitraje de 1990, para los ajustes de precios de transferencia, y la Directiva (UE) 2017/1852, que obliga a resolver en un plazo determinado cualquier conflicto de doble imposición entre Estados miembros. La razón por la que no existe un convenio europeo único es, como en muchas otras materias, que exigiría unanimidad, y ningún Estado quiere renunciar a los equilibrios que ha negociado bilateralmente con cada socio.

En definitiva, España dispone de una red amplia, alineada con el Modelo OCDE, especialmente densa con sus principales socios comerciales y con Iberoamérica, y modernizada de forma sistemática en las últimas dos décadas. Ha incorporado con rapidez los estándares antiabuso de BEPS a través del Instrumento Multilateral, y su ley interna contiene mecanismos de exención con cláusula de sujeción que evitan tanto la doble imposición como la doble no imposición. No obstante, también presenta algunas debilidades. En particular, destaca la lentitud de los procedimientos, que hace que la entrada en vigor de un convenio se prolongue muchos años, y la renegociación de convenios antiguos se acumule. También la heterogeneidad de la red. Como reconocía el propio Ministerio al ratificar el Instrumento Multilateral, los convenios españoles son dispares en contenido y redacción, fruto de negociaciones realizadas a lo largo de medio siglo con modelos distintos, lo que dificulta su aplicación y genera diferencias de trato difíciles de justificar según el país de origen de la inversión. Por último, los procedimientos amistosos siguen siendo lentos en la práctica, pese a los compromisos en el ámbito multilateral; y el arbitraje solo está disponible en los convenios que lo prevén expresamente y, dentro de la Unión, a través de la Directiva de 2017. 

\subsubsection{Competencia y procedimiento}
La competencia para celebrar convenios corresponde en exclusiva al Estado. Se apoya en el art. $149$ de la Constitución, que atribuye al Estado las relaciones internacionales, y sobre Hacienda en general. De esta forma, las Comunidades Autónomas no negocian convenios, aunque sus efectos alcanzan a los tributos cedidos que gestionan, y los territorios forales los aplican dentro de su propio sistema tributario.

Durante la negociación juega un papel esencial la Subdirección General de Fiscalidad Internacional, dentro de la Dirección General de Tributos de la Secretaría de Estado de Hacienda, responsable de negociar los convenios y de su aplicación. Todos los convenios españoles se basan en el Modelo OCDE.

Un convenio sigue un procedimiento en seis fases: (i) \textbf{negociación técnica} entre las dos Administraciones, sobre la base del Modelo OCDE y de las posiciones de cada país; (ii) \textbf{rúbrica del texto} por los negociadores, que fija el contenido acordado; (iii) \textbf{autorización de la firma} por el Consejo de Ministros y firma del convenio; (iv) \textbf{dictamen del Consejo de Estado}, sobre la necesidad o no de autorización parlamentaria; (v) \textbf{autorización de las Cortes Generales}, porque los convenios fiscales afectan a la Hacienda Pública y modifican la aplicación de leyes tributarias, supuestos para los que el art. $94$ de la Constitución exige esa autorización; y, (vi) \textbf{ratificación}, intercambio de instrumentos con el otro Estado y publicación en el BOE, momento en que el convenio se incorpora al ordenamiento interno.

El procedimiento explica por qué entre la firma de un convenio y su entrada en vigor pueden pasar años: ambos Estados deben completar sus propios trámites internos, y basta con que uno se retrase, como ocurrió con EEUU, para que el convenio no se aplique.

\subsubsection{Política española para eliminar la doble imposición}
En los convenios españoles, como ya hemos visto, se conviene la imputación como regla general. La elección es coherente con su posición histórica de país que recibía más inversión de la que emitía, y que por tanto quería proteger su recaudación sobre las rentas que sus residentes obtenían fuera.

No obstante, en la ley interna se establece una exención para las rentas empresariales. Por tanto, como vemos, con independencia de lo que digan los convenios, la ley española ofrece a las sociedades mecanismos unilaterales que en la práctica son más relevantes:
\begin{lnum}
    \item La exención de dividendos y plusvalías de participaciones (art. $21$ de la LIS). Se aplica cuando la participación en la filial es de al menos el $5\%$ y la filial ha estado sujeta a un impuesto análogo con un tipo nominal de al menos el $10\%$, requisito que se entiende cumplido si reside en un país con convenio que incluya intercambio de información. Desde 2021, la exención cubre el $95\%$ de la renta y no el $100\%$: el $5\%$ restante tributa como compensación de los gastos de gestión de la participación, que el grupo ya ha deducido. Este recorte lo introdujo la Ley de Presupuestos para 2021;
    \item La exención de las rentas de establecimientos permanentes en el extranjero (art. $22$), con requisitos análogos;
    \item Las deducciones por doble imposición internacional (arts. $31$ y $32$), que aplican el método de imputación ordinaria cuando no procede la exención;
    \item En el IRPF, la deducción por doble imposición internacional del art. $80$, también de imputación ordinaria.
\end{lnum}

El requisito del tipo nominal mínimo del $10\%$ del art. $21$ es un buen ejemplo práctico de lo que vimos antes. Se trata de una cláusula de sujeción incorporada a la ley interna, donde España solo renuncia a gravar los dividendos de una filial extranjera si esta ha tributado razonablemente en su país; mientras que si ha tributado a tipo cero, España grava el dividendo. Anticipándose, a menor escala y veinte años antes, la lógica que veremos en el marco multilateral (BEPS $2.0$).

\paragraph{Países sin convenio: de los paraísos fiscales a las jurisdicciones no cooperativas}
Durante tres décadas, España identificó los paraísos fiscales mediante una lista aprobada por el Real Decreto 1080/1991, que inicialmente incluía $48$ territorios. La lista se fue reduciendo a medida que España firmaba con ellos acuerdos de intercambio de información o convenios con cláusula de intercambio.

La Ley 11/2021 sustituyó el término paraíso fiscal por el de jurisdicción no cooperativa, siguiendo los trabajos de la Unión y la OCDE, y amplió los criterios. Ya no se atiende solo a la transparencia, sino también a la equidad fiscal: se identifican los territorios que facilitan sociedades extraterritoriales dirigidas a atraer beneficios sin actividad económica real, que aplican una tributación baja o nula, o que carecen de transparencia. Un cambio coherente con lo que veremos después: lo relevante ya no es solo si un territorio comparte información, sino si permite trasladar beneficios sin actividad.

La lista actual se aprobó por la Orden HFP/115/2023, de 9 de febrero, y comprende, entre otros, Anguila, Bahréin, Barbados, Bermudas, Gibraltar, Guernsey, Jersey, la Isla de Man, las Islas Caimán o las Islas Vírgenes Británicas. Como podemos suponer, la inclusión de Gibraltar tiene un evidente componente político para España. En mayo de 2026, el Ministerio de Hacienda sometió a información pública un proyecto de orden que modifica la de 2023, excluyendo o incluyendo países, territorios y regímenes según los criterios de transparencia y equidad fiscal de la ley. \textcolor{blue}{Conviene comprobar si esa modificación se ha aprobado y con qué contenido antes del examen.}

Pero, ¿qué implicaciones tiene estar en estas listas? Como sabemos, estar en la lista tiene efectos a lo largo de todo el sistema tributario español. Por ejemplo, presunciones de residencia en España para las personas que se trasladan a estos territorios ([\authorfont{Ver Tema 4.B.16}]), exclusión de las exenciones de dividendos del art. $21$, aplicación reforzada de la transparencia fiscal internacional, retenciones más elevadas, y obligaciones reforzadas de documentación en las operaciones vinculadas ([\authorfont{Ver Tema 4.B.17}]).


\clearpage
\section{Dimensión multilateral: Marco BEPS}
\puntosclave{
    La globalización y digitalización han facilitado la erosión de bases imponibles y traslado de beneficios a territorios de nula o baja tributación.
    
    Para resolver esta externalidad global, la OCDE y el G-20 ha propuesto una solución global a través del paquete BE PS, al que posteriormente se han sumado un gran número de países que forman el Marco Inclusivo
    
    Actualmente se está trabajando para aplicar las 15 acciones de BEPS de 2015. así como los Pilares 1 y 2 de BEPS 2.0 acordadas en ocLUbre de 202 1.
    
    Los dos pilares pretenden resolver los retos que genera la digitalización, mediante un instrumento que permita gravar de fo1ma más justa a las grandes multinacionales y repartir los beneficios según donde se hayan generado (pilar 1 ) y mediante la introducción de un tipo mínimo efectivo para sociedades del 15% (pilar 2)
}   

Los dos bloques anteriores han mostrado cómo un Estado grava por su cuenta a quien no reside en él (el IRNR) y cómo dos Estados se reparten entre sí el derecho a gravar una misma renta (los CDI). Cada uno resuelve bien sus respectivos objetivos. En particular, los convenios bilaterales resuelven con eficacia el problema para el que nacieron, la doble imposición entre dos países que se conocen y se reconocen.

Pero, de la misma forma, ambos instrumentos comparten un límite que ninguno puede superar por sí solo: solo obligan a quien los firma y no a actuar de manera óptima. Imaginemos dos países idénticos, $A$ y $B$, que pueden fijar un tipo alto ($25\%$) o bajo ($12,5\%$). Las cifras son ilustrativas, en recaudación:

\begin{table}[h]
    \centering
    \caption{Placeholder Caption}
    \label{tab:placeholder_label}
    \begin{tabular}{lcc}
        \hline
        & B mantiene 25\% & B baja a 12,5\% \\
        \hline
        A mantiene 25\% & A: 25 / B: 25 & A: 10 / B: 30 \\
        A baja a 12,5\% & A: 30 / B: 10 & A: 12,5 / B: 12,5 \\
        \hline
    \end{tabular}
\end{table}

Si ambos mantienen el tipo alto, cada uno recauda $25$. Si uno baja y el otro no, el que baja atrae base ajena y recauda $30$ (un tipo menor sobre una base mucho mayor), mientras el otro cae a $10$. Para cada país, bajar el tipo es la mejor opción con independencia de lo que haga el otro. Ambos bajan, y terminan recaudando $12,5$ cada uno, \textbf{la mitad de lo que tendrían cooperando}. Ningún Estado puede salir de este equilibrio por su cuenta, y un convenio bilateral entre $A$ y $B$ tampoco basta si existe un país $C$ dispuesto a ofrecer un $5\%$. Esta es la razón de fondo por la que el problema exige una solución multilateral: hace falta un acuerdo que ate a la vez a todos los jugadores relevantes.

El problema que debemos abordar, por tanto, no es un defecto técnico de redacción de los convenios, sino un problema de acción colectiva, y los problemas de acción colectiva solo se resuelven coordinando las acciones de todos los jugadores. Por eso la respuesta ha tenido que ser multilateral.

\subsection{Problema: ¿cuál es el problema multilateral?}
La expresión inglesa \textit{Base Erosion and Profit Shifting} (BEPS) une dos fenómenos distintos que conviene separar. La erosión de la base consiste en reducir artificialmente la base imponible en la jurisdicción donde se genera la renta. El ejemplo típico es la deducción de grandes pagos de intereses a otra sociedad del grupo. El traslado de beneficios, por su parte, consiste en desplazar el beneficio gravable de un país de tributación alta a otro de tributación baja. Siendo el ejemplo típico el traslado de la propiedad intelectual y los ingresos que genera desde Estados Unidos a Bermudas. 

En la práctica, ambos fenómenos casi siempre aparecen juntos: \textbf{el pago que erosiona la base de un país es, a la vez, el ingreso que aparece en otro con tributación nula o reducida}.

Como es evidente, nos encontramos ante un fallo de mercado, concretamente una \textbf{externalidad fiscal}. La génesis de esta rama teórica (competencia fiscal) nace en dos artículos casi simultáneos del Journal of Urban Economics de 1986, uno de ZODROW y MIESZKOWSKI y otro de WILSON. Su argumento central es que, cuando el capital puede moverse entre jurisdicciones, cada gobierno que sube su impuesto sobre el capital pierde base imponible en favor de sus vecinos. Al decidir su tipo, cada gobierno tiene en cuenta esa pérdida propia, pero no tiene en cuenta la ganancia que su subida genera en los demás. por tanto, es una externalidad positiva ignorada: cada uno fija un tipo más bajo del que fijaría si internalizara el efecto sobre el conjunto. El resultado es una tributación del capital inferior a la socialmente deseable y, con ella, una provisión insuficiente de bienes públicos, o bien un desplazamiento de la carga hacia los factores que no pueden moverse, sobre todo el trabajo. Además, KANBUR y KEEN (1993), mostraron de manera célebre cómo y por qué la competencia fiscal no es simétrica. El país pequeño tiene un incentivo mucho mayor a rebajar su tipo. La razón es aritmética. Para un país pequeño, lo que pierde bajando el tipo sobre su escasa base propia es poco, y lo que gana atrayendo base de sus grandes vecinos es mucho. Para un país grande ocurre al revés. Esto explica por qué Irlanda, Luxemburgo, Países Bajos, Singapur o las jurisdicciones caribeñas han sido protagonistas de este fenómeno, y por qué sus Gobiernos se resistieron después con más fuerza a cualquier suelo mínimo.

La competencia fiscal también tiene defensores. El primero procede de TIEBOUT (1956) y fue tratado en otros temas ([\authorfont{Ver Tema 4.B.14}]). La competencia entre jurisdicciones permite a los ciudadanos elegir la combinación de impuestos y servicios que prefieren. El segundo procede de BRENNAN y BUCHANAN (1980) y su Estado leviatán: si el Estado tiende a comportarse como un Leviatán que maximiza recaudación, la competencia entre Estados es un freno útil a esa tendencia. Precisamente, quienes se opusieron a los acuerdos de la OCDE se apoyaron en la idea de que un suelo mínimo global equivale a un cártel de gobiernos que elimina la disciplina que la competencia impone.

Por qué la elusión, siendo legal, no es socialmente óptima. Respondo aquí a tu comentario del bloque 2, porque es donde encaja con más naturalidad. [Probable] 
La elusión es legal por definición, pero genera al menos cuatro costes que no aparecen en ninguna liquidación:

En primer lugar, distorsiona la asignación de recursos. Los beneficios se localizan según el tipo, no según la productividad. Una filial con dos empleados en Bermudas puede declarar más beneficio que una fábrica con miles de trabajadores en Alemania. En segundo lugar, consume recursos reales improductivos. Los asesores, las estructuras societarias y los litigios cuestan dinero y talento que no producen nada; solo redistribuyen recaudación entre Estados y entre la empresa y el fisco. Es la misma idea de coste social de la evasión y la elusión que SLEMROD ha desarrollado en su obra sobre cumplimiento tributario. En tercer lugar, la elusóni genera competencia desleal. La multinacional que traslada beneficios compite con una ventaja fiscal frente a la empresa local que no puede hacerlo, aunque esta sea más eficiente: la ventaja no nace de producir mejor, sino de planificar mejor. Por ultimo, la elusón desplaza la carga y erosiona la legitimidad: lo que no pagan las bases móviles lo pagan las inmóviles, sobre todo el trabajo y el consumo. Y la percepción de que las grandes empresas no pagan lo que les corresponde deteriora la moral fiscal del resto de contribuyentes.

\subsubsection{¿Cómo se erosiona la base imponible?}
Conviene conocer los canales concretos, a través de los cuales se erosiona la base. 

\paragraph{Canal 1: los intangibles}
Es el canal más poderoso, y la razón es sencilla: una patente, una marca o un algoritmo no tienen ubicación física, y su valor es muy difícil de medir porque no existe un mercado de referencia con el que compararlos. La técnica consiste en ubicar la titularidad del intangible en una sociedad del grupo situada en un territorio de baja tributación, y hacer que las sociedades operativas le paguen cánones por usarlo. Por ejemplo, supongamos que un grupo vende en España por $100$ millones de euros, con $70$ millones de costes operativos; su beneficio real en España es de $30$ millones. Si ese beneficio tributa en España al $25\%$, paga $7,5$ millones. Pero supongamos que la filial española paga a otra sociedad del grupo, situada en una jurisdicción con tipo cero y titular de la marca, un canon de $25$ millones por usarla. El beneficio declarado en España baja a $5$ millones, que pagan $1,25$ millones de impuesto. Los $25$ millones restantes aparecen en la jurisdicción de tipo cero y no pagan nada. El resultado final es que el grupo paga $1,25$ millones sobre un beneficio real de $30$, un tipo efectivo del $4,2\%$ en lugar del $25\%$. La base española se ha erosionado en $25$ millones y el beneficio se ha trasladado.

Los precios de transferencia siguen una lógica muy similar. Las operaciones entre sociedades del mismo grupo no están sometidas a la disciplina del mercado, de modo que el grupo puede fijar precios distintos de los de mercado para decidir dónde se queda el beneficio. Ya vimos que la respuesta clásica es el principio de plena competencia (arm's length): valorar esas operaciones como si se hubieran pactado entre empresas independientes. Por tanto, el canal de los intangibles es, en el fondo, el caso extremo de este mismo problema: el principio de plena competencia funciona razonablemente bien cuando existe un precio de mercado comparable, y muy mal cuando no existe, que es precisamente lo que ocurre con los intangibles únicos.

La estructura Double Irish y el caso Apple. La estructura combinaba sociedades irlandesas que, por diferencias entre la ley irlandesa y la estadounidense sobre la residencia de sociedades, no eran residentes fiscales en ningún sitio. En el caso Apple, la Comisión Europea concluyó en 2016 que dos resoluciones fiscales irlandesas (de 1991 y 2007) habían permitido a dos filiales del grupo tributar sobre un beneficio muy inferior al normal. El núcleo del problema era que Irlanda aceptó que las licencias de propiedad intelectual se atribuyeran a unas sedes centrales de esas filiales que no tributaban en ningún país. La Comisión ordenó recuperar $13.000$ millones de euros más intereses, unos $14.300$ millones en total. El Tribunal General anuló la decisión en 2020, pero el TJUE la restableció definitivamente el 10 de septiembre de 2024 (asunto C-465/20 P). Irlanda cerró la estructura Double Irish a nuevas incorporaciones en 2015, con un periodo transitorio que terminó a finales de 2020.


\paragraph{Canal 2: la deuda intragrupo} 
Una sociedad del grupo en un país de tipo alto se endeuda con otra del mismo grupo en un país de tipo bajo. Los intereses se deducen en el primero y tributan poco o nada en el segundo. Se trata del mismo mecanismo, con \textbf{intereses en lugar de cánones}. Aprovecha, además, el sesgo a favor de la deuda de casi todos los impuestos sobre sociedades: los intereses son deducibles y los dividendos no.

Otra situación en la misma linea aprovecha que dos países califican de forma distinta un mismo instrumento o una misma entidad. Un pago puede ser deducible en el país del pagador (que lo considera interés) y no tributar en el del receptor (que lo considera dividendo exento). O una entidad puede ser transparente para un país y opaca para otro, de modo que su renta no tributa en ninguno.

\paragraph{Canal 6: la economía digital} 
La digitalización permite obtener grandes beneficios en un país sin tener allí presencia física. Como el reparto clásico del derecho a gravar exige presencia física (el establecimiento permanente), un país puede tener millones de usuarios y clientes de una plataforma y no poder gravar ni un euro de su beneficio. Este canal es distinto de los anteriores, pues no depende de ninguna maniobra artificial, sino de que las reglas del siglo XX no reconocen el valor que se crea en el mercado de destino. Por eso requerirá una respuesta propia: el Pilar 1.


\subsubsection{Magnitud del problema: ¿qué dicen los datos y por qué se discuten?}
La OCDE, en el informe final de la Acción $11$ del plan BEPS (2015), estimó que la erosión de bases costaba a los Estados entre el $4\%$ y el $10\%$ de la recaudación mundial del impuesto sobre sociedades, entre $100.000$ y $240.000$ millones de dólares al año. El propio intervalo, muy amplio, ya muestra lo difícil que es medir algo que, por definición, se intenta ocultar.

La estimación académica más influyente, ha sido el trabajo de TØRSLØV, WIER y ZUCMAN (2023), que cambió la escala del debate. Utilizando nuevas estadísticas macroeconómicas sobre filiales extranjeras, muestran que las filiales de multinacionales extranjeras son un orden de magnitud más rentables que las empresas locales en varios países de baja tributación, y a partir de esa diferencia estiman que el $36\%$ de los beneficios de las multinacionales se traslada a paraísos fiscales en el mundo; las multinacionales estadounidenses trasladan el doble que las demás en relación con sus beneficios en el exterior, muy lejos de las estimaciones anteriores que situaban la cifra cerca del $40\%$. 

La lógica de su método es sencilla y conviene entenderla. La rentabilidad media de las empresas locales es parecida en paraísos y en países que no lo son, lo que indica que la actividad económica no es en general más rentable en los paraísos; lo desproporcionado es la rentabilidad de las filiales de multinacionales en ellos, y eso se explica por el traslado de beneficios. Dicho de otro modo: si una filial en Irlanda gana muchísimo más por euro de salario que cualquier empresa irlandesa, ese exceso no puede proceder de la productividad irlandesa.\footnote{%
    Evidentemente, estas cifras no son pacíficas. BLOUIN y ROBINSON sostienen que buena parte de las estimaciones basadas en datos estadounidenses están infladas por un error contable. Los datos de la Oficina de Análisis Económico de EEUU atribuyen toda la renta de las filiales de niveles inferiores al país de la filial de nivel superior. Es habitual que una multinacional estadounidense controle sus filiales extranjeras a través de una sociedad holding situada en un país de baja tributación. Por ejemplo, una filial alemana propiedad de una holding irlandesa. La renta de esa holding irlandesa incluye entonces lo que gana la filial alemana, que ya ha tributado en Alemania. El resultado es que beneficio realmente obtenido en Alemania parece obtenido en Irlanda. 

Evidentemente, el debate sigue abierto. CLAUSING había estimado que el traslado de beneficios costaba a EEUU entre $77.000$ y $111.000$ millones de dólares en 2012; BLOUIN y ROBINSON, corrigiendo el doble cómputo, lo rebajaron a entre $10.000$ y $32.000$ millones. CLAUSING respondió con datos de 2017 y volvió a obtener pérdidas superiores a $100.000$ millones, argumentando que el método de sus críticos también tiene problemas y puede excluir algunos canales. Conviene precisar que la corrección de BLOUIN y ROBINSON afecta sobre todo a los estudios que usan esos datos estadounidenses. Sin embargo, el trabajo de TØRSLØV, WIER y ZUCMAN se apoya en otras fuentes, aunque tampoco libres de críticas. Una revisión de la literatura encuentra estimaciones que van desde menos del $5\%$ de la renta trasladada hasta más del $30\%$, y concluye que todavía es pronto para sacar conclusiones definitivas sobre la importancia cuantitativa de la elusión internacional.
}

Entonces, ¿quién gana y quién pierde? Si los beneficios trasladados volvieran a sus países de origen, los beneficios domésticos aumentarían en torno a un $20\%$ en los países de la Unión con tributación alta, un $10\%$ en EEUU y un $5\%$ en los países en desarrollo, mientras caerían un $55\%$ en los paraísos fiscales. Los principales ganadores son Irlanda, Luxemburgo o Singapur, que aplican tipos bajos pero sobre una base muy grande, y los principales perdedores, los países de la Unión que no son paraísos.


\subsection{BEPS 1.0 (2013-2015): Reforzar el sistema}
Tras la crisis financiera de 2008, los países del G20 afrontaban a la vez déficits públicos elevados, programas de austeridad y una opinión pública indignada por los casos antes citados. Resultaba políticamente insostenible pedir sacrificios fiscales a los ciudadanos mientras las grandes multinacionales pagaban tipos efectivos mínimos. 

En 2013, la OCDE y el G20 unieron fuerzas y desarrollaron un Plan de Acción con $15$ acciones, destinadas a dar coherencia a las normas sobre actividades transfronterizas, mejorar la transparencia y la certidumbre, y asegurar que los beneficios tributen allí donde se desarrolla la actividad económica y se crea el valor. Con esto en mente, existían dos grandes caminos de solución:
\begin{la}
    \item Mantener el sistema de entidades separadas y el principio de plena competencia, tapando sus agujeros. Cada filial sigue tributando como si fuera independiente, y se endurecen las reglas para que los precios intragrupo reflejen la realidad;
    \item Sustituirlo por un reparto por fórmula (formulary apportionment). Se calcula el beneficio consolidado del grupo mundial y se reparte entre países según indicadores objetivos, como ventas, empleo o activos. Es el sistema que ya aplican los estados de EEUU entre sí y la lógica de la propuesta europea de base consolidada que veremos en el 5.4.
\end{la}

Evidententemente, los grandes países europeos (Francia, Alemania, Reino Unido) fueron los impulsores más activos ya que eran los principales perdedores del traslado de beneficios. Estados Unidos adoptó una posición ambivalente. Por un lado, compartía el interés en combatir la elusión. Por otro, las principales multinacionales afectadas eran estadounidenses, y buena parte del beneficio trasladado por ellas lo era a costa de otros países, no del propio EEUU.\footnote{%
    De hecho, Estados Unidos nunca firmó el Instrumento Multilateral que veremos más adelante.
} Por otro lado, en el extremo opuesto, estaban las jurisdicciones de baja tributación (Irlanda, Luxemburgo, Países Bajos), que aceptaron el marco general, pero defendieron con firmeza su soberanía para fijar tipos lo que condujo a que BEPS $1.0$, respetando esa línea roja, no fijara ningún tipo mínimo. Por último, los países en desarrollo criticaron desde el principio que las reglas se diseñaran en una mesa donde apenas estaban representados. 

Finalmente, el Plan se publicó en julio de 2013 y lo respaldaron los líderes del G20 en la cumbre de San Petersburgo de septiembre de ese año. Los informes finales de las $15$ acciones se publicaron en octubre de 2015, y los líderes del G20 aprobaron el paquete en la cumbre de Antalya de noviembre de 2015. Y, con el fin de salvar las distancias, optó por \textbf{reforzar el principio de plena competencia} sin no sustituirlo, eligiendo por tanto el primer camino de los dos.\footnote{%
    Las razones fueron tres. Primero, sustituir el principio de plena competencia exigía renegociar miles de convenios y un consenso político inexistente. Segundo, cualquier fórmula de reparto genera ganadores y perdedores claros. Los países con mucho mercado ganaban, los que concentran producción o intangibles pierden, y ningún perdedor aceptaba el cambio. Tercero, la propia OCDE había defendido durante décadas el principio de plena competencia como estándar internacional.
} Este es el punto conceptual más importante que debemos retener. 

Sin embargo, esta decisión es crucial para entender lo que vino después: al no cambiar la lógica del reparto, BEPS $1.0$ no pudo resolver el problema de la economía digital, donde el valor se crea en el mercado sin presencia física. Esa es la razón por la que hará falta BEPS $2.0$.

\subsubsection{Plan de Acción: 15 Acciones, 3 Bloques}
¿Cuáles son, entonces, las acciones previstas por el Plan? La OCDE agrupa las acciones en torno a tres objetivos: coherencia, sustancia y transparencia. A ellos se suman una acción sobre la economía digital y otra sobre la forma de aplicarlo todo. 

\paragraph{Bloque 1: Coherencia}
El primer bloque, relativo a la \textbf{coherencia}, buscaba que las diferencias entre sistemas nacionales \textbf{no generasen rentas que no tributen} en ningún sitio. Dentro de este marco se adoptaron las siguientes acciones:
\begin{lnum}
    \item \textbf{Mecanismos híbridos} (Acción 2). Neutraliza los desajustes entre normas. Si un pago es deducible en un país y no tributa en el otro, se niega la deducción en el primero o se obliga a integrar la renta en el segundo. Es una regla "de enlace": la respuesta de un país depende de lo que haga el otro;
    \item \textbf{Transparencia fiscal internacional} (Acción 3) (CFC). Permite al país de la matriz gravar directamente la renta que sus filiales acumulan en territorios de baja tributación, sin esperar a que se reparta como dividendo. España ya la tenía desde los años noventa; la acción la refuerza y armoniza;
    \item \textbf{Intereses} (Acción 4). Limita la deducción de gastos financieros netos a un porcentaje del beneficio operativo (EBITDA), recomendando una horquilla entre el $10\%$ y el $30\%$. Por ejemplo, si una filial española tiene un EBITDA de $10$ millones y paga $6$ millones de intereses a otra sociedad del grupo. Con un límite del $30\%$, solo puede deducir $3$ millones; los otros $3$ no reducen la base española;
    \item \textbf{Prácticas fiscales perniciosas} (Acción 5). Esta acción tiene dos piezas. En primer lugar, exige que los regímenes preferenciales, en particular los patent box que gravan a tipo reducido las rentas de patentes, estén ligados a actividad real de investigación en el país, el llamado \textbf{enfoque del nexo}. En segundo lugar, obliga a intercambiar espontáneamente información sobre las resoluciones fiscales (tax rulings) que las Administraciones conceden a empresas concretas.
\end{lnum}

\paragraph{Bloque 2: Sustancia}
El segundo bloque, relativo a la \textbf{sustancia}, busca que la tributación siga a la actividad económica real, no a la forma jurídica. Se compone de tres acciones:
\begin{lnum}
    \item \textbf{Abuso de convenios} (Acción 6). La cláusula de limitación de beneficios y la cláusula de propósito principal, que ya trabajamos antes, que niega los beneficios del convenio cuando obtenerlos fue uno de los objetivos principales de una operación;
    \item \textbf{Establecimiento permanente} (Acción 7). Amplía la definición para impedir la elusión artificial: fragmentar actividades para que cada una parezca auxiliar, o vender a través de comisionistas que formalmente no contratan en nombre de la empresa;
    \item \textbf{Precios de transferencia} (Acciones 8 a 10). Refuerzan el principio de plena competencia en los tres ámbitos donde más fallaba: intangibles, riesgos y capital, y operaciones de alto riesgo. La idea central es que el beneficio de un intangible debe asignarse a quien realmente desarrolla, mejora, mantiene, protege y explota ese intangible, y no a la sociedad que simplemente figura como titular jurídico. En nuestro ejemplo, una sociedad sin empleados en un territorio de tipo cero ya no podría quedarse con el canon de $25$ millones solo por ser titular de la marca.
\end{lnum}

\paragraph{Bloque 3: Transparencia}
El tercer bloque, relativo a la \textbf{transparencia}, parte de una idea sencilla: las Administraciones no pueden combatir lo que no ven. Esta compuesto por cuatro acciones:
\begin{lnum}
    \item \textbf{Medición} (Acción 11). Encomienda a la OCDE estimar la magnitud del problema; es de donde procede la estimación de $100.000$ a $240.000$ millones de dólares que vimos;
    \item \textbf{Declaración obligatoria} (Acción 12). Obliga a los asesores y contribuyentes a comunicar a la Administración los esquemas de planificación fiscal agresiva;
    \item \textbf{Informe país por país} (Country-by-Country Reporting) (Acción 13). Obliga a los grupos multinacionales con una facturación consolidada igual o superior a $750$ millones de euros a informar, jurisdicción por jurisdicción, de sus ingresos, beneficios, impuestos pagados, empleados y activos. Las Administraciones se intercambian estos informes. Es, probablemente, la \textbf{acción de mayor impacto práctico de todo el paquete}. Por primera vez, cada Administración ve dónde declara beneficios el grupo y dónde tiene actividad real, y puede detectar de un vistazo si una filial con dos empleados declara cientos de millones de beneficio;
    \item \textbf{Resolución de controversias} (Acción 14). Mejora los procedimientos amistosos entre Administraciones cuando surgen conflictos de doble imposición.
\end{lnum}

Hay que tener en cuenta que no todas las acciones obligan igual. Existe obligación de aplicar como estándar mínimo cuatro de ellas: la lucha contra las prácticas fiscales perniciosas y el intercambio de resoluciones fiscales (acción 5), las cláusulas antiabuso en los convenios (acción 6), el informe país por país (acción 13) y los mecanismos de resolución de controversias (acción 14). Su aplicación se revisa regularmente por el Marco Inclusivo mediante evaluaciones entre pares. Sin embargo, el resto de acciones son recomendaciones, enfoques comunes o buenas prácticas que cada país decide si adopta. Distinción que explica los resultados desiguales: donde hubo obligación y vigilancia, el avance fue claro; donde no, dependió de la voluntad de cada país.

Además, en respuesta a la crítica de los países en desarrollo, la OCDE creó en 2016 el Marco Inclusivo sobre BEPS, abierto a cualquier jurisdicción. Sus miembros participan en igualdad de condiciones con los países de la OCDE y del G20, pero a cambio deben comprometerse a aplicar los cuatro estándares mínimos. Ha pasado de alrededor de $100$ miembros en 2017 a más de $140$ en la actualidad. Los resultados de las evaluaciones entre pares son medibles. En la acción 5, el Foro sobre Prácticas Fiscales Perniciosas ha revisado $319$ regímenes preferenciales, se han producido casi $50.000$ intercambios de información sobre $23.000$ resoluciones fiscales identificadas, y $73$ jurisdicciones cumplen plenamente el estándar, mientras las $58$ restantes han recibido $61$ recomendaciones de mejora. La crítica de fondo, sin embargo, persiste. Los países en desarrollo entraron en un marco cuyas reglas ya estaban escritas, y se comprometieron a aplicar estándares en cuyo diseño no habían participado.


\subsubsection{Implementación: Instrumento Multilateral}
Ahora bien, una vez alcanzado el acuerdo relativo al Plan, $3.000$ convenios debían renegociarse uno a uno, pues muchas de las medidas de BEPS, en particular las de las acciones $2, 6, 7$ y $14$, exigían modificar los convenios de doble imposición. Como cabe esperar, esta fue una gran limitación práctica. 

Por tanto, se adoptó una solución un tanto heterodoxo: un \textbf{convenio sobre convenios}. La acción $15$ previó la creación de un instrumento multilateral que permitiera a las jurisdicciones interesadas introducir las medidas de BEPS en sus convenios bilaterales sin tener que renegociarlos. El Convenio Multilateral para aplicar las medidas relacionadas con los tratados fiscales para prevenir la erosión de bases y el traslado de beneficios (MLI) se adoptó el 24 de noviembre de 2016. Se firmó en París el 7 de junio de 2017 y entró en vigor el 1 de julio de 2018, tras alcanzar el número mínimo de ratificaciones. Para España entró en vigor el 1 de enero de 2022.

Su funcionamiento merece explicarse porque es genuinamente ingenioso. En primer lugar, cada país firmante \textbf{elabora una lista de los convenios que quiere que el MLI modifique}, los llamados convenios cubiertos. En segundo lugar, para cada disposición del MLI, \textbf{cada país indica si la acepta o formula una reserva}. Por último, un convenio bilateral concreto solo se modifica si ambos países lo han incluido en su lista, y cada disposición solo se aplica si ambos la han aceptado.

El resultado fue un sistema de \textbf{casación de preferencias}. Un único texto multilateral que produce efectos distintos en cada relación bilateral. El único elemento que no admitía reservas eran los estándares mínimos, en particular la cláusula de propósito principal de la acción $6$, que prácticamente todos los firmantes han adoptado.\footnote{%
    La Unión Europea no se limitó a esperar a que cada Estado aplicara BEPS voluntariamente. Aprobó en 2016 la Directiva antielusión conocida como ATAD 1, que incorporó al Derecho de la Unión, entre otras, las normas sobre deducibilidad de gastos financieros (acción 4), transparencia fiscal internacional (acción 3) y asimetrías híbridas. La respuesta a estas últimas era solo parcial, por lo que al año siguiente se aprobó la ATAD 2, que incorpora de forma más completa las recomendaciones de la acción 2. La consecuencia fue, por tanto, muy importante, puesto que lo que para el resto del mundo eran recomendaciones, para los Estados miembros de la Unión eran obligaciones, con un plazo de trasposición y la posibilidad de un procedimiento de infracción si no se cumplen. España tuvo una aplicación asimétrica, temprana en unas medidas, tardía en otras:
\begin{la}
    \item Limitación de intereses. España se adelantó a BEPS: ya en 2012 introdujo un límite a la deducción de gastos financieros netos del 30\% del beneficio operativo, que la Ley 27/2014 del Impuesto sobre Sociedades mantuvo en su artículo 16. Coincide con el máximo de la horquilla recomendada por la acción 4;
    \item Informe país por país. Se aplica en España desde los periodos impositivos iniciados en 2016, mediante el modelo 231, con el mismo umbral de 750 millones de euros;
    \item Patent box. España adaptó su reducción de rentas procedentes de intangibles al enfoque del nexo de la acción 5; y,
    \item Trasposición de la ATAD. Aquí España llegó tarde. La Ley 11/2021, de medidas de prevención y lucha contra el fraude fiscal, traspuso la Directiva ATAD e introdujo cambios relevantes en el "impuesto de salida" (artículo 19 de la Ley del Impuesto sobre Sociedades) y en el texto refundido del IRNR. Las normas sobre híbridos llegaron antes, mediante el Real Decreto-ley 4/2021, de 9 de marzo. España, como el resto de países de la Unión, debía haber traspuesto la ATAD 2 antes de que acabara 2019, y también había incumplido el plazo de trasposición de las disposiciones antiabuso de la ATAD 1, que vencía el 31 de diciembre de 2018, por lo que la Comisión Europea abrió procedimientos de infracción. La trasposición se completó después mediante la Ley 5/2022, con un nuevo artículo 15 bis en la Ley del Impuesto sobre Sociedades, y el Real Decreto-ley 18/2022 para las asimetrías híbridas invertidas.
\end{la}

La doctrina señaló que el uso del Real Decreto-ley, en lugar de una ley ordinaria, podía plantear dudas de constitucionalidad a la luz del artículo $86$ de la Constitución y de la doctrina del Tribunal Constitucional. Fue la propia inacción legislativa la que generó la situación de "extraordinaria y urgente necesidad" invocada para evitar los controles del procedimiento legislativo ordinario.
}


\paragraph{Valoración: ¿cumplió sus objetivos?}
Los resultados fueron ambiguos. El avance más claro se produjo en \textbf{transparencia}. Antes de 2015, las Administraciones tributarias apenas sabían dónde declaraban sus beneficios los grandes grupos; hoy reciben informes país por país de todos los grupos con más de $750$ millones de facturación e intercambian miles de resoluciones fiscales. Se cerraron además los canales más burdos: los híbridos, el abuso de convenios más evidente y los regímenes preferenciales sin actividad real. La estructura Double Irish, por ejemplo, desapareció en este periodo.

Sin embargo, tenía tres limitaciones que explican por qué BEPS $1.0$ no fue el final del camino. En primer lugar, \textbf{no cambió la lógica del reparto}. Al mantener el principio de plena competencia y la exigencia de presencia física, dejó sin respuesta el problema de la economía digital, como la propia acción $1$ reconocía. En segundo lugar, \textbf{no fijó ningún suelo}. Un país seguía siendo libre de aplicar un tipo del $0\%$ siempre que hubiera actividad real. La competencia fiscal pasó de las estructuras de papel a los tipos nominales y a los incentivos, pero no desapareció. Por último, \textbf{dependía de la voluntad} de cada país. Fuera de los cuatro estándares mínimos, la aplicación fue desigual. Y el mayor país exportador de capital, Estados Unidos, se mantuvo al margen del MLI.\footnote{%
    Además, al reforzar un sistema basado en el análisis caso por caso de funciones, activos y riesgos, BEPS $1.0$ aumentó considerablemente la complejidad del cumplimiento y de la inspección. Esa complejidad beneficia relativamente a los grandes grupos, que pueden permitirse los mejores asesores, y a las Administraciones con más recursos, en detrimento de los países en desarrollo. Esta crítica, junto con la frustración de los países que querían gravar la economía digital, preparó el terreno para una reforma que sí cambiara la lógica del sistema.
}


\subsection{BEPS 2.0 (2019-hoy): del reparto por presencia al reparto por mercado y al suelo mínimo}
El gran fracaso fue el de la acción $1$, relativa a la economía digital. Se limitaba a reconocer que la economía digital planteaba un problema de fondo, pero no logró ninguna solución: no sabía \textbf{cómo gravar a una empresa digital sin presencia física}, y no ponía \textbf{ningún límite a la competencia a la baja en los tipos}. Entre 2016 y 2021 la falta de acuerdo covirtió esos dos vacíos en un problema político urgente. 

Por un lado, varios países decidieron actuar por su cuenta con impuestos sobre servicios digitales: gravámenes indirectos sobre los ingresos (no sobre el beneficio) de las grandes plataformas por publicidad en línea, intermediación o venta de datos de usuarios.\footnote{%
    EEUU ya había introducido en 2017, con la reforma fiscal de la Administración Trump, un impuesto mínimo propio sobre las rentas extranjeras de sus multinacionales (conocido como GILTI). Era un precedente directo de lo que sería el Pilar 2: EEUU ya gravaba a sus propias empresas cuando tributaban poco en el exterior, y le interesaba que el resto del mundo hiciera lo mismo para que sus empresas no quedaran en desventaja.
} Francia fue la pionera en 2019. España aprobó el suyo mediante la Ley 4/2020: un impuesto indirecto del $3\%$ sobre publicidad en línea, servicios de intermediación y transmisión de datos de usuarios, aplicable solo a grandes grupos que superen umbrales de ingresos mundiales y en España. Entró en vigor el 16 de enero de 2021 y sigue vigente. Para los países que los aprobaron, estos impuestos eran una respuesta legítima a la parálisis de la negociación. Para las empresas afectadas y para EEUU, eran una forma de discriminación. Por otro lado, la Oficina del Representante Comercial de EEUU abrió en 2019 y 2020 investigaciones que concluyeron en que los impuestos digitales de Austria, Francia, India, Italia, España, Turquía y Reino Unido eran discriminatorios y dirigidos sobre todo contra los gigantes tecnológicos estadounidenses. Llegó a amenazar con aranceles del $25\%$ sobre más de $2.000$ millones de dólares de importaciones de esos países. 

En definitiva, el \textbf{conflicto fiscal amenazaba con convertirse en conflicto comercial}, lo que dio a todas las partes un incentivo real para buscar un acuerdo. Y, en 2021, la Administración Biden, a través de la secretaria del Tesoro Janet Yellen, impulsó activamente un impuesto mínimo global, lo que desbloqueó la negociación. El 8 de octubre de 2021, $136$ de los $140$ países del Marco Inclusivo alcanzaron un compromiso político sobre una solución de dos pilares, la lógica de cada uno responde a uno de los dos vacíos de BEPS $1.0$.

\subsubsection{Pilar 1: Gravar donde está el mercado}
El Pilar $1$ reasigna una parte del beneficio de los mayores grupos a los países donde están sus clientes, aunque no tengan presencia física allí. Por tanto, rompe por primera vez con la exigencia de presencia física y con el principio de plena competencia para una parte del beneficio.

¿Cómo se hace esto? Lo primero que se define es la \textbf{Cantidad A}. La mecánica es la siguiente. Del beneficio total del grupo, se considera normal el que corresponde a un margen del $10\%$ sobre sus ventas mundiales. El resto es beneficio residual. El $25\%$ de ese beneficio residual se reparte entre los países donde el grupo tiene ventas, en proporción a esas ventas, aunque no tenga allí ningún establecimiento.

Por ejemplo, imaginemos un grupo digital que factura $100.000$ millones de euros en todo el mundo y obtiene un beneficio de $30.000$ millones: un margen del $30\%$. Sobre este beneficio, se considera normal el $10\%$ de las ventas: $10.000$ millones. Sin embargo, el beneficio residual queda como la diferencia, $30.000 - 10.000 = 20.000$ millones. Por tanto, la Cantidad A, que es el $25\%$ del residual, son $5.000$ millones. Entonces, si el $5\%$ de las ventas del grupo se produce en España, a España le corresponde gravar el $5\%$ de esos $5.000$ millones: $250$ millones de beneficio, aunque el grupo no tenga aquí ni una oficina.

La intuición económica es que una parte del beneficio extraordinario de estos grupos se debe al mercado donde venden, y no solo a los países donde están sus ingenieros o su propiedad intelectual. Es, por primera vez en un siglo, una concesión parcial al reparto por fórmula que BEPS $1.0$ había rechazado. Esta cantidad \textbf{solo a aplica que grupos empresariales} que facturan más de $20.000$ millones de euros y tienen un margen de beneficio antes de impuestos superior al $10\%$ sobre sus ventas. Se excluyen las industrias extractivas y los servicios financieros regulados. En la práctica afectaría a unos $100$ grupos, en su mayoría grandes tecnológicas.

Por otro lado, se define la \textbf{Cantidad B}, que es un enfoque simplificado para fijar los precios de transferencia de las actividades básicas de comercialización y distribución. En lugar de un análisis caso por caso, que para muchos países en desarrollo resulta muy costoso, fija rentabilidades estandarizadas para estas actividades. La OCDE publicó su informe en febrero de 2024 y cada país decide si la aplica.

\paragraph{Resultados: Bloqueado}
Al firmarse el acuerdo de 2021, los países con impuestos digitales (Austria, Francia, Italia, España y Reino Unido) pactaron con EEUU un compromiso. Mantendrían sus impuestos de forma transitoria hasta la entrada en vigor del Pilar 1, pero lo recaudado por ellos que superase lo que las empresas habrían pagado con el Pilar 1 se descontaría después de lo que esos países recibieran por la Cantidad A. A cambio, EEUU retiró los aranceles propuestos contra los cinco países; India y Turquía no se sumaron inicialmente al acuerdo. El compromiso vencía en principio al implantarse el Pilar 1 o el 31 de diciembre de 2023; ante los retrasos, se prorrogó hasta el 30 de junio de 2024.

Sin embargo, el convenio multilateral para aplicar la Cantidad A solo puede entrar en vigor si lo ratifica una masa crítica de países que incluya a los de las matrices de la mayoría de los grupos afectados, como EEUU, China, Japón, Reino Unido, Alemania o Francia. El mayor obstáculo político era la aprobación de EEUU. Obstáculo que se volvió insalvable el 20 de enero de 2025, primer día de la segunda Administración Trump, cuando un memorándum presidencial declaró que el acuerdo fiscal global no tiene fuerza ni efecto en EEUU salvo que el Congreso lo apruebe, desautorizó los compromisos de la Administración anterior y encargó al Tesoro estudiar medidas de protección frente a países con normas fiscales que afecten desproporcionadamente a empresas estadounidenses.

El resultado es que la Cantidad A no se ha aplicado en España ni hay normas internas que asignen derechos de gravamen conforme a ella, y España sigue aplicando su impuesto unilateral sobre servicios digitales a la espera de una solución internacional. La paradoja es completa: el Pilar 1 nació para acabar con los impuestos digitales unilaterales, y su bloqueo los ha consolidado.

\subsubsection{Pilar 2: un suelo mínimo a la competencia fiscal}
El Pilar 2, también conocido como reglas GloBE (Global anti-Base Erosion), establece una tributación mínima efectiva global del $15\%$ para los grupos multinacionales con ingresos superiores a $750$ millones de euros. De esta forma, responde a la competencia fiscal fijando un tipo efectivo mínimo, de modo que bajar el tipo por debajo de ese umbral deja de ser rentable.\footnote{%
    ¿Es el $15\%$ demasiado bajo? ZUCMAN y el Observatorio Fiscal de la Unión han criticado que el tipo sea bajo, inferior al tipo nominal de la mayoría de países desarrollados, y que la exclusión por sustancia reduzca considerablemente la recaudación. Desde este punto de vista, el Pilar 2 legitima una tributación del $15\%$ que, antes, muchos habrían considerado propia de un paraíso fiscal. Además, al limitarse la competencia en tipos, los países tienden a competir con otros instrumentos que el Pilar 2 trata con más benevolencia, como ciertos créditos fiscales reembolsables o las subvenciones directas. La competencia no desaparece: cambia de forma.
}

Tres precisiones son esenciales para entenderlo. En primer lugar, es un \textbf{tipo efectivo}, no nominal. No importa el tipo que fije la ley de cada país, sino lo que el grupo paga realmente en relación con su beneficio, calculado con reglas propias y comunes. En segundo lugar, se \textbf{calcula país por país}. Si un grupo tributa al $25\%$ en Alemania y al $5\%$ en un paraíso, no se hace la media: se comprueba cada jurisdicción por separado, y la del $5\%$ no alcanza el mínimo. Por último, no obliga a ningún país a subir su tipo: un país puede seguir aplicando un $5\%$, lo que cambia es que la diferencia hasta el $15\%$ la cobrará otro país.

Evidentemente, los críticos sostienen que funciona como un cártel de gobiernos que elimina la disciplina de la competencia fiscal. No obstante, sus defensores responden que solo fija un suelo, y un suelo relativamente bajo, dejando todo el margen de competencia por encima de él. Por otro lado, muchos países en desarrollo usan incentivos fiscales para atraer inversión que el Pilar 2 los neutraliza. Si un país concede una exención, la diferencia hasta el $15\%$ la cobra el país de la matriz, que suele ser un país rico. Además, las reglas de prioridad benefician sobre todo a los países de la matriz. 

\paragraph{Implementación: Cuatro reglas (subsidiarias)}
La pregunta que queda por resolver es, ¿cómo se puede implementar este objetivo en la práctica? El mecanismo consta de cuatro piezas con un orden de prioridad muy concreto:
\begin{lnum}
    \item \textbf{Impuesto complementario nacional mínimo cualificado} (QDMTT, por sus siglas en inglés). El país donde la renta tributa poco puede cobrar él mismo la diferencia hasta el $15\%$. Tiene prioridad sobre las demás reglas;
    \item \textbf{Regla de inclusión de rentas} (IIR). Si el país de baja tributación no cobra la diferencia, la cobra el país de la matriz, que impone un impuesto complementario a la entidad matriz por las rentas de sus filiales sujetas a baja tributación;
    \item \textbf{Regla de beneficios insuficientemente gravados} (UTPR). En defecto de las dos anteriores, la tercera es una regla de apoyo. Si ni el país de baja tributación ni el de la matriz cobran la diferencia, la cobran los demás países donde opera el grupo, negando deducciones o exigiendo un ajuste equivalente; y,
    \item \textbf{Cláusula de sujeción a impuestos} (STTR). Por último, la cuarta es una regla de convenio, que permite al país de origen gravar determinados pagos entre partes vinculadas (intereses, cánones) cuando en el país de destino tributan por debajo de un tipo mínimo. Ese mínimo es del $9\%$, y su principal beneficiario son los países en desarrollo, que suelen ser el origen de estos pagos.
\end{lnum}

No obstante, cabe destacar que no toda la renta de baja tributación paga el complemento. En particular, se excluye un rendimiento normal ligado a la actividad real: un porcentaje de los salarios y de los activos materiales del grupo en cada país. ¿Por qué? En esencia, porque el Pilar $2$ no persigue a quien tiene fábricas y empleados en un país de tipo bajo, sino el beneficio de papel sin actividad real detrás.

Veamos como funcionaría su implementación con un ejemplo. Supongamos un grupo con matriz en España tiene una filial en el país $X$, cuyo impuesto sobre sociedades es del $5\%$. La filial obtiene un beneficio de $100$ millones y tiene salarios y activos que dan derecho a una exclusión por sustancia de $20$ millones.

El impuesto pagado en $X$ es del $5\%$, $5$ millones, un tipo efectivo del $5\%$. La diferencia hasta el mínimo queda entonces en $10$ puntos. Por tanto, la base del complemente, que se calcula como el beneficio menos exclusión por sustancia, $= 100 - 20$, son $80$ millones; y el impuesto complementario del $10\%$ sobre $80 = 8$ millones.

Ahora bien, ¿quién cobra esos $8$ millones? Si $X$ ha aprobado un QDMTT, los cobra $X$. Si no, los cobra España como país de la matriz, a través de la IIR. Y, por último, si ninguno de los dos lo hace, los cobran los demás países donde opere el grupo, a través de la UTPR.

En definitiva, en los tres casos el grupo paga lo mismo: $13$ millones en total. Lo único que cambia es qué país se queda con la diferencia.

¿Por qué el Pilar 2 cambia el juego? Si recordamos el dilema del prisionero, cada país tenía incentivo a bajar su tipo porque así atraía base ajena. En este mismo contexto, el Pilar 2 altera ese cálculo de forma radical. Si el país $X$ baja su tipo al $5\%$ para atraer beneficios, esos beneficios siguen pagando el $15\%$ en total. La única diferencia es que los $10$ puntos restantes se los queda otro país. Bajar del $15\%$ deja de atraer base y solo sirve para regalar recaudación a otros. Por tanto, la respuesta racional de los países de baja tributación es aprobar su propio QDMTT y cobrar ellos mismos la diferencia. Y eso es exactamente lo que ha ocurrido. 

Irlanda, Suiza y varios otros países de baja tributación han introducido impuestos complementarios nacionales para los grandes grupos. Como es normal, preferían cobrar ellos la diferencia a que la cobraran otros. Más de $56$ jurisdicciones, incluidos todos los Estados miembros de la Unión, las habían incorporado a su derecho interno las reglas GloBE, pero el Congreso estadounidense no aprobó una legislación similar.

Tras el memorándum de enero de 2025, la Administración estadounidense presionó para que sus multinacionales no quedaran sujetas a la IIR y la UTPR aplicadas por otros países. En junio de 2025, el G7 alcanzó un acuerdo político para excluir a los grupos con matriz estadounidense de esas dos reglas, con el argumento de que la legislación estadounidense ya grava suficientemente sus beneficios domésticos y extranjeros; y se convirtió en norma el 5 de enero de 2026, cuando el Marco Inclusivo aprobó un paquete de orientaciones administrativas que, en la práctica, exime a las multinacionales estadounidenses de la mayor parte de las reglas del Pilar 2.\footnote{
El paquete fue respaldado por más de $145$ jurisdicciones y sus elementos principales son:
\begin{la}
    \item Un puerto seguro lado a lado. Los grupos cuya matriz esté en un país con un régimen fiscal suficientemente sólido pueden optar por no aplicar la IIR ni la UTPR en los ejercicios que comiencen a partir del 1 de enero de 2026. A esa fecha, EEUU era el único país reconocido con un régimen cualificado para ello, aunque otros países pueden solicitar que se evalúe el suyo;
    \item Por otro lado, los Impuestos complementarios nacionales (QDMTT) no se ven afectados y seguirán aplicándose a los grupos estadounidenses donde estén en vigor. Es decir, el país de baja tributación donde opera una filial estadounidense sigue pudiendo cobrar la diferencia hasta el $15\%$.
\end{la}
} Este acuerdo, como es normal, divide a la doctrina. Para los defensores del marco, era el precio necesario para salvar el Pilar 2 del rechazo frontal de su principal economía y evitar una guerra comercial. Por otro lado, para sus críticos, rompe el principio de igualdad de trato entre grupos según la nacionalidad de su matriz, y confirma que las reglas fiscales internacionales se doblan ante la presión del país más poderoso. En todo caso, el QDMTT sale reforzado como la pieza central del sistema.



\begin{specialblock}{\textbf{España}: Aplicación de BEPS $2.0$}
    La Ley de Presupuestos Generales del Estado para 2022 ya se adelantó en cierta medida a la Directiva europea al incluir una tributación mínima del $15\%$ para determinados contribuyentes del Impuesto sobre Sociedades. Conviene no confundirla con el Pilar 2. Aquella norma fija una cuota líquida mínima del $15\%$ de la base imponible del propio impuesto sobre sociedades español, para empresas con una cifra de negocios de al menos $20$ millones de euros o que tributen en régimen de consolidación fiscal. Mide el tipo sobre la base imponible española, mientras que el Pilar 2 mide un tipo efectivo con reglas propias y país por país.

    El 21 de diciembre de 2024 se publicó en el BOE la Ley 7/2024, de 20 de diciembre, que traspone la Directiva (UE) 2022/2523. La ley crea un Impuesto Complementario aplicable a los grupos con una facturación consolidada superior a $750$ millones de euros en al menos dos de los cuatro ejercicios anteriores, para garantizar en todos los casos un umbral mínimo de tributación del $15\%$. Tiene efectos para los periodos impositivos iniciados desde el 1 de enero de 2024, con algunas reglas de aplicación diferida. Su desarrollo reglamentario se aprobó mediante el Real Decreto 252/2025, de 1 de abril.
    
    El impuesto tiene tres modalidades, que reproducen las tres reglas del Pilar 2. Una de ellas es el impuesto complementario nacional, cuya finalidad es garantizar que, si las entidades de un gran grupo situadas en España no alcanzan una tributación mínima del $15\%$ en territorio español, sea la Administración española quien cobre la diferencia. Es, en esencia, el QDMTT español: \textbf{España se asegura de que nadie más cobre el complemento correspondiente a rentas obtenidas en su territorio}.
    
    Por último, el bloqueo del Pilar 1 ha hecho que España mantenga el Impuesto sobre Determinados Servicios Digitales de la Ley 4/2020. Su futuro depende de dos variables que España no controla: si algún día se desbloquea el Pilar 1, y si la Administración estadounidense decide reactivar la vía de las represalias comerciales frente a los países que mantienen estos impuestos.
\end{specialblock}

\begin{specialblock}{\textbf{Unión Europea}: el laboratorio que va por delante, y su límite}
    La Unión comparte con la OCDE el objetivo de combatir la erosión de bases, pero tiene además un problema propio que la OCDE no tiene: un mercado interior con libertad de circulación de capitales y libertad de establecimiento. Dentro de ese mercado, una empresa puede trasladar su sede, sus filiales o sus beneficios de un Estado miembro a otro sin ninguna barrera. Esto convierte la competencia fiscal entre Estados miembros en un problema más intenso que la competencia global: el dilema del prisionero se juega aquí entre socios que comparten moneda, mercado y reglas de competencia, pero no fiscalidad directa. 
    
    Ante esta situación, la Unión dispone de tres instrumentos:
    \begin{la}
        \item Directivas de armonización fiscal, basadas en el art. $115$ del TFUE. Obligan a todos los Estados miembros, pero exigen unanimidad en el Consejo;
        \item El control de ayudas de Estado del art. $107$ del TFUE. La Comisión puede actuar sola, sin necesidad de acuerdo de los Estados; y,
        \item Tribunales de la Unión, que son los que deciden en último término.
    \end{la}

    Cualquier directiva de fiscalidad directa necesita el voto favorable de los $27$ Estados miembros, lo que da a cada país un derecho de veto. Esto introduce, dentro de la propia Unión, el mismo problema de acción colectiva que vimos a escala global, pues basta con que un solo Estado que se beneficia de la situación actual se niegue para bloquear cualquier reforma. La Comisión propuso en 2019 pasar gradualmente a la mayoría cualificada en materia fiscal, utilizando las cláusulas del Tratado que lo permiten, pero la propuesta no prosperó precisamente porque su adopción también exigía unanimidad. Ante la dificultad de armonizar, la Comisión encontró en la última década un camino alternativo: tratar ciertas resoluciones fiscales concedidas por un Estado a una empresa concreta como una ayuda de Estado ilegal. La lógica es que, si una Administración tributaria concede a una multinacional un tratamiento más favorable que el que aplicaría a cualquier otra empresa en su misma situación, le está otorgando una ventaja selectiva equivalente a una subvención.
    
    No obstante, los resultados judiciales de aplicar esta segunda vía han sido desiguales. Por ejemplo, la Comisión ganó definitivamente en el Caso de Apple (irlanda), cuando el TJUE confirmó en septiembre de 2024 que Irlanda había concedido una ayuda ilegal de unos $13.000$ millones de euros más intereses. Sin embargo, en otros casos (Starbucks, Países Bajos; Fiat, Luxemburgo; y, Amazon, Luxemburgo) la Comisión perdió. En el caso Fiat, el TJUE anuló la decisión de la Comisión, al considerar que esta no podía evaluar la ventaja con arreglo a un principio de plena competencia distinto del que aplica el propio Derecho luxemburgués. 
    
    Por tanto, la lección de fondo es que el control de ayudas de Estado no puede sustituir a la armonización. El TJUE exige que la comparación se haga con las propias normas del Estado miembro, no con un estándar europeo que todavía no existe. Si el Derecho interno de un país es laxo para todos, no hay ventaja selectiva que combatir.
    
    \textbf{Transparencia: la saga de las directivas de cooperación administrativa}
    
    La transparencia es el terreno donde la Unión ha avanzado de forma más sostenida, porque intercambiar información no obliga a ningún Estado a renunciar a su soberanía sobre los tipos o las bases, y es mucho más fácil de aprobar. La Directiva 2011/16/UE de cooperación administrativa (conocida como DAC) se ha modificado sucesivamente hasta nueve veces para ampliar el tipo de información que las Administraciones se intercambian. Conviene agruparlas por la información que cubren:
    \begin{la}
        \item Rentas financieras y de capital (DAC1 y DAC2). Intercambio automático de información sobre rentas de los contribuyentes y cuentas financieras, siguiendo el estándar común de la OCDE;
        \item Resoluciones fiscales (DAC3). Intercambio automático de las resoluciones fiscales transfronterizas concedidas a empresas: es la respuesta europea directa a LuxLeaks, y la traducción europea de la acción 5 de BEPS;
        \item Informe país por país (DAC4). Intercambio de los informes de la acción $13$ de BEPS entre Administraciones;
        \item Titularidad real (DAC5). Acceso de las Administraciones tributarias a la información sobre los titulares reales de sociedades, recogida por las normas contra el blanqueo de capitales;
        \item Esquemas de planificación agresiva (DAC6). Obliga a los intermediarios, como asesores y bancos, a comunicar a la Administración los esquemas transfronterizos que presenten determinados indicios de planificación agresiva. Es la traducción europea de la acción 12 de BEPS;
        \item Plataformas digitales (DAC7). Impone a los operadores de plataformas digitales la obligación de informar sobre los ingresos obtenidos por quienes venden bienes o prestan servicios a través de ellas;
        \item Criptoactivos (DAC8). Establece el intercambio automático de información sobre criptoactivos entre los Estados miembros. Se adoptó el 17 de octubre de 2023, se basa en el estándar internacional de la OCDE sobre criptoactivos y se aplica desde el 1 de enero de 2026. Los proveedores de servicios de criptoactivos deben recopilar datos de las operaciones de sus usuarios residentes en la Unión desde esa fecha, y presentar la primera información entre el 1 de enero y el 30 de septiembre de 2027;
        \item Pilar 2 (DAC9). La Directiva (UE) 2025/872 establece un marco para el intercambio de información del Pilar 2 entre Estados miembros, con plazo de trasposición hasta el 31 de diciembre de 2025. Su utilidad práctica es considerable, pues permite a un gran grupo presentar una única declaración del impuesto complementario en un Estado miembro, que después se comparte con el resto.
    \end{la}
    
    De hecho, en el ámbito de la transparencia, la Unión fue más lejos que BEPS con la Directiva (UE) 2021/2101, que obliga a los grandes grupos a publicar su informe país por país, no solo a entregarlo a las Administraciones. Estas normas están ya en vigor en todos los Estados miembros, y la primera ronda de publicaciones de la mayoría de grupos afectados vence a finales de 2026, referida al ejercicio 2025. Como podemos apreciar, se trata de un cambio cualitativo, pues la información deja de estar reservada a las Administraciones y pasa a estar sometida al escrutinio de la opinión pública, los medios y los investigadores, lo que convierte la reputación en un instrumento más de disciplina fiscal.\footnote{%
        La Comisión ha anunciado para 2026 un paquete de simplificación fiscal (ómnibus fiscal) centrado en ordenar las interacciones entre las distintas normas de la Unión, en particular la directiva de cooperación administrativa, previsto para el segundo trimestre de 2026. Conviene verificar su estado exacto antes del examen, porque podría refundir buena parte de esta saga.
    }

    \textbf{Antiabuso: la ATAD, la propuesta Unshell y la lista de jurisdicciones no cooperativas}
    
    Conviene añadir aquí dos elementos en los que la Unión fue más allá de BEPS; (i) una \textbf{norma general antiabuso obligatoria} para todos los Estados miembros, que permite a la Administración ignorar cualquier mecanismo cuyo propósito principal sea obtener una ventaja fiscal contraria al objeto de la norma; y, (ii) un \textbf{impuesto de salida que grava las plusvalías latentes de los activos} cuando una empresa traslada su residencia o sus activos fuera del país. 
    
    Por otro lado, la propuesta de Directiva, presentada en 2021, pretendía combatir la planificación fiscal vinculada al uso de sociedades pantalla, obligando a las empresas de la Unión a comunicar a la Administración la información necesaria para evaluar si tenían presencia sustancial y actividad económica real, y negando beneficios fiscales a las que no la tuvieran. El proyecto preveía una serie de pruebas para identificar a las sociedades sin sustancia suficiente, a las que se negaría el acceso a los beneficios de los convenios y de las directivas europeas. Tras más de tres años bloqueada, sin que los Estados miembros alcanzaran un consenso sobre sus aspectos clave, el informe del ECOFIN de 18 de junio de 2025 la dio oficialmente por abandonada. La Comisión la retiró formalmente en su Programa de Trabajo para 2026, publicado el 21 de octubre de 2025, aunque los indicadores de alto riesgo que proponía podrían incorporarse a las señales de alerta de la DAC6 en el marco del ómnibus fiscal. 
    
    Además, desde 2017, el Consejo mantiene una lista de países terceros que no cumplen los estándares europeos de buena gobernanza fiscal: transparencia, competencia fiscal leal y aplicación de los estándares mínimos de BEPS. La lista se revisa dos veces al año. Su eficacia depende de las consecuencias defensivas que los Estados miembros asocien a ella, y ha sido criticada por dos motivos: (i) excluir sistemáticamente a los propios Estados miembros, algunos de los cuales figuran entre los principales destinos del traslado de beneficios; y, (ii) su carácter político, pues un país puede salir de ella con compromisos formales antes de haber cambiado realmente su régimen.
    
    \textbf{Base común: del proyecto de BICCIS a BEFIT, y el debate entre fórmula y plena competencia}
    
    Desde 2011, la Comisión intenta que los grupos que operan en varios Estados miembros calculen su beneficio con unas reglas comunes, consoliden el resultado de todo el grupo en la Unión y lo repartan después entre los Estados miembros según una fórmula. Es exactamente el segundo camino que BEPS $1.0$ rechazó. La propuesta de Base Imponible Consolidada Común del Impuesto sobre Sociedades (BICCIS) de 2011 no prosperó; se relanzó en 2016 dividida en dos fases, la base común y la consolidación, y tampoco salió adelante.
    
    Su atractivo es que elimina de raíz el problema de los precios de transferencia dentro de la Unión. Si el beneficio del grupo se reparte por fórmula, da igual a qué precio se vendan bienes o servicios entre sus filiales europeas. 
    
    Por ejemplo, supongamos que un grupo obtiene en la Unión un beneficio consolidado de $100$ millones. Si la fórmula reparte a partes iguales según ventas, empleo y activos, y en España el grupo tiene el $30\%$ de sus ventas, el $20\%$ de su empleo y el $10\%$ de sus activos europeos, a España le corresponde gravar un $20\%$ del beneficio: $20$ millones, con independencia de dónde haya declarado el grupo sus beneficios mediante precios intragrupo. La dificultad es que toda fórmula genera ganadores y perdedores. Los países con grandes mercados o mucho empleo ganan; los que concentran activos intangibles, sedes o beneficios contables pierden. Y los perdedores tienen derecho de veto.
    
    Por eso, en septiembre de 2023, la Comisión presentó la propuesta BEFIT (Business in Europe: Framework for Income Taxation). Consiste en una serie de normas comunes para calcular la base imponible de las entidades de la Unión que formen parte de grupos con ingresos consolidados mundiales superiores a un umbral determinado. Es decir, se armoniza el cálculo de la base, pero se congela el reparto existente, precisamente para no crear perdedores que bloqueen la propuesta. BEFIT (y la propuesta paralela de tributación según las normas del Estado de la sede para pymes, HOT) siguen pendientes en el Consejo, sin fecha de adopción.
    
    Por otro lado, la Directiva relativa a los precios de transferencia, presentada junto a BEFIT en 2023. Esta pretende incorporar al Derecho de la Unión normas comunes sobre precios de transferencia. No obstante, fue retirada en el Programa de Trabajo para 2026 y podría sustituirse por una plataforma europea que busque soluciones consensuadas y no vinculantes a problemas prácticos de precios de transferencia.
    
    DEBRA fue otra propuesta, de 2022, que pretendía corregir el sesgo del impuesto sobre sociedades a favor de la deuda mediante dos medidas: una reducción de la base ligada al aumento de fondos propios y una limitación adicional a la deducibilidad de los gastos financieros. El Consejo suspendió las negociaciones en diciembre de 2022; y la Comisión la retiró formalmente en el mismo Programa de Trabajo para 2026, junto con Unshell, la Directiva de precios de transferencia y el impuesto sobre transacciones financieras.

    
    \textbf{Directiva del Pilar 2: la unanimidad en acción}
    
    [Seguro] La Directiva (UE) 2022/2523, de 15 de diciembre de 2022, trasladó el Pilar 2 al Derecho de la Unión. Añadió además normas para evitar cualquier discriminación entre situaciones transfronterizas y nacionales: el tipo mínimo efectivo del 15\% se aplica tanto a los grupos multinacionales como a los grandes grupos que operan solo dentro del mercado interior. [Probable] Esta extensión a los grupos puramente nacionales no es un capricho: sin ella, la Directiva podría haber vulnerado la libertad de establecimiento, al tratar peor a un grupo con filiales en otros Estados miembros que a uno que opere en un solo país.
    
    Su aprobación ilustra perfectamente lo que vimos. Polonia bloqueó primero la Directiva durante varios meses en 2022, y después fue Hungría la que ejerció su derecho de veto. El bloqueo húngaro se resolvió en diciembre de 2022 dentro de un acuerdo político más amplio, que afectaba también a la aprobación del plan de recuperación húngaro y a los fondos europeos que tenía congelados. Es un ejemplo claro de cómo la regla de la unanimidad permite a un Estado miembro utilizar su voto en materia fiscal como moneda de cambio en negociaciones sobre asuntos completamente distintos.
    
    En definitiva, todos los Estados miembros de la Unión han incorporado las reglas del Pilar 2 a su derecho interno. Sin embargo, el acuerdo lado a lado de enero de 2026 obliga ahora a los Estados miembros a adaptar su aplicación para eximir a los grupos estadounidenses de la IIR y la UTPR. La Unión, que había actuado como primer motor del acuerdo global, termina adaptando su propia normativa a una excepción negociada por un país que nunca lo aplicó.
\end{specialblock}



\subsection{Frente abierto: ¿qué foro debe gobernar la fiscalidad internacional?}
Una crítica que BEPS nunca resolvió es el de los países en desarrollo, que entraron en el Marco Inclusivo cuando las reglas de BEPS ya estaban escritas, y la recaudación de forma que favorece sobre todo a los países donde están las matrices, que suelen ser ricos. 

Para una parte importante de los países en desarrollo, estas dos críticas llevan a una misma conclusión: la OCDE, una organización de $38$ países mayoritariamente desarrollados, no es un foro legítimo para fijar las reglas fiscales de todo el mundo, por muy abierto que sea su Marco Inclusivo. Participar en igualdad de condiciones en un foro cuya agenda, secretaría técnica y calendario controlan otros no es, a su juicio, igualdad real.

El debate enfrenta dos argumentos, ambos serios. En primer lugar, está el argumento de la \textbf{legitimidad}, defendido por el Grupo Africano, el G77 y buena parte de la sociedad civil. Naciones Unidas es el único foro verdaderamente universal, donde cada país tiene un voto, y por tanto el único con legitimidad para fijar reglas que afectan a todos. Además, sus prioridades serían distintas: reforzar la tributación en la fuente, que es donde los países en desarrollo generan sus rentas, frente a la preferencia de la OCDE por la tributación en la residencia. Por otro lado, estaría el argumento de la \textbf{capacidad técnica}, defendido por la mayoría de países desarrollados. La OCDE lleva sesenta años elaborando el modelo de convenio que ya vimos en el bloque 4, tiene la experiencia técnica y funciona por consenso, lo que garantiza que las reglas se apliquen de verdad. Un proceso en Naciones Unidas por mayoría podría aprobar normas que los principales países exportadores de capital nunca aplicarían, duplicando foros y generando más incertidumbre.\footnote{%
    Ya vimos que existe un Modelo de convenio de la ONU, más favorable a la tributación en la fuente que el de la OCDE, y que España lo tiene en cuenta en algunos convenios con países en desarrollo. En 2021, el Comité de Expertos de la ONU incorporó a ese modelo un nuevo artículo 12B, que permite al país de la fuente gravar los ingresos por servicios digitales automatizados. Es la respuesta de Naciones Unidas al mismo problema que el Pilar 1 pretendía resolver, con una diferencia importante: no exige umbrales tan altos como los $20.000$ millones de euros de la Cantidad A, de modo que alcanza a muchas más empresas, algo especialmente relevante para países cuyos mercados, individualmente, son pequeños.
}

\subsubsection{Convención Marco de Naciones Unidas: calendario, contenido y actores}
Por impulso del Grupo Africano, con Nigeria como portavoz en nombre del G77 y China, la Asamblea General de Naciones Unidas adoptó a finales de 2023 una resolución para iniciar negociaciones hacia una Convención Marco sobre Cooperación Fiscal Internacional. Es el primer intento de construir un marco fiscal multilateral vinculante bajo los auspicios de Naciones Unidas, rompiendo con décadas de gobernanza dominada por la OCDE y sus miembros mayoritariamente del Norte.

Al votarse los términos de referencia en agosto de 2024, $110$ países votaron a favor, 8 en contra y 43 se abstuvieron, con el apoyo de casi todo el continente africano. La mayoría de los Estados miembros de la Unión, incluida España, figuró entre las abstenciones, mientras EEUU y otros países desarrollados votaron en contra. Es un reparto que refleja con bastante precisión la línea que separa a los exportadores de capital de los importadores, la misma que explicaba la preferencia de unos por el método de imputación y de otros por el de exención. Mediante su resolución 79/235, la Asamblea General creó un comité intergubernamental de negociación, dirigido por los Estados miembros, para redactar la Convención Marco junto con dos protocolos tempranos, y decidió que el comité se reuniría en 2025, 2026 y 2027 con al menos tres sesiones sustantivas al año, en Nueva York y en otras sedes, incluida Nairobi. El objetivo es concluir las negociaciones en 2027.

El primer protocolo temprano tratará la tributación de las rentas procedentes de la prestación de servicios transfronterizos en una economía cada vez más digitalizada y globalizada. En la sesión organizativa de febrero de 2025 se eligió como tema del segundo protocolo la prevención y resolución de controversias fiscales. Otros temas prioritarios, como la tributación de la economía digital, las medidas contra los flujos financieros ilícitos relacionados con los impuestos, y la lucha contra la evasión y elusión de las grandes fortunas, se abordarán en protocolos posteriores a 2027. [Probable] Este último punto conecta con el debate sobre el impuesto mínimo a los milmillonarios que ya trabajamos en el tema de imposición patrimonial.

Además, en esa misma sesión se acordó que la Convención se aprobará por mayoría simple y los protocolos por mayoría de dos tercios. Los países desarrollados preferían el consenso y los países en desarrollo y emergentes, la mayoría simple. La diferencia no es menor: en la OCDE, cualquier gran país puede bloquear un acuerdo, como hizo EEUU con el Pilar 1; en Naciones Unidas, los países en desarrollo, que son mayoría, pueden aprobar normas sin el apoyo de los países ricos. Pero una convención aprobada sin los principales países exportadores de capital corre el riesgo de quedarse en papel, porque son precisamente sus multinacionales las que habría que gravar.

El 3 de febrero de 2025, primer día de la sesión organizativa, EEUU anunció que se retiraba de las negociaciones, afirmando que eran incompatibles con sus prioridades y que suponían una extralimitación indeseada, y que se opondría a cualquier resultado. EEUU llamó a otros países a abandonar las negociaciones con él, pero ninguno lo hizo. [Probable] La decisión es coherente con su posición ante el acuerdo de la OCDE que vimos en el 5.3: la Administración estadounidense rechaza cualquier marco fiscal multilateral que limite su capacidad de decisión, sea en la OCDE o en Naciones Unidas.

Las negociaciones sustantivas comenzaron en agosto de 2025 y continúan hasta 2027. La quinta sesión del comité se celebró en Nueva York del 3 al 13 de agosto de 2026, y la siguiente está prevista en Nairobi del 30 de noviembre al 11 de diciembre de 2026. [Suposición] No tengo información fiable sobre los resultados de la sesión de agosto de 2026; conviene verificarlos antes del examen, porque es el punto de todo el bloque con mayor probabilidad de haber cambiado.

\subsubsection{Valoración: ¿competencia o complementariedad entre foros?}
Existen tres escenarios posibles, y la doctrina está dividida sobre cuál es más probable.

El primero de ellos sería el de la \textbf{complementariedad}. Bajo este escenario, la Convención de la ONU se centra en los principios generales, la resolución de controversias y los temas que la OCDE no ha resuelto, como la tributación de los servicios transfronterizos o de las grandes fortunas, mientras la OCDE conserva el desarrollo técnico de los estándares ya existentes. Por otro lado, es posible que los conflictos deriven en un escenario de \textbf{fragmentación}. Dos sistemas de reglas conviven sin coordinarse, donde los países en desarrollo aplican el marco de Naciones Unidas entre ellos y los desarrollados siguen con el de la OCDE. Sería el peor escenario para la seguridad jurídica de las empresas y el más cercano, paradójicamente, al problema de descoordinación con el que empezó este bloque. Por último, seria el de la irrelevancia. El más plausible, puesto que sin los principales países exportadores de capital, la Convención se aprobará pero apenas se aplica.

No obstante, es importante entender que la mera existencia del proceso de Naciones Unidas ha cambiado el equilibrio de poder dentro de la propia OCDE. El paquete lado a lado de enero de 2026, por ejemplo, reforzó expresamente el impuesto complementario nacional como mecanismo principal para proteger las bases de los países en desarrollo. 

Con independencia de su resultado final, la presión de un foro alternativo obliga al foro establecido a tomarse más en serio a quienes antes podía ignorar.

























































\clearpage
\section{Conclusión} 


\subsection{Recapitulación}


\subsection{Extensiones y relación con otras partes del Temario}



\subsection{Opinión}

\subsubsection{Idea final (Salida o cierra)}

\clearpage
\pagenumbering{gobble}
\MyBibliography

Texto refundido de la Ley del IRNR (Real Decreto Legislativo 5/2004), arts. 9, 18, 19.2 (gravamen complementario), 24, 25 y 46.
OCDE: Model Tax Convention, art. 24.3 (no discriminación de los establecimientos permanentes).
Sentencia del Tribunal de Justicia de 14 de febrero de 1995, Schumacker, asunto C-279/93.
Sentencia del Tribunal de Justicia de 12 de junio de 2003, Gerritse, asunto C-234/01.
Directiva 2003/48/CE del Consejo, sobre fiscalidad de los rendimientos del ahorro, derogada por la Directiva (UE) 2015/2060 del Consejo, de 10 de noviembre de 2015.
Convenio entre España y Estados Unidos para evitar la doble imposición, de 22 de febrero de 1990, y protocolo de 14 de enero de 2013.
Constitución Española, arts. 94.1, 96 y 149.1.3 y 14.
Convenios entre España y Alemania (Bonn, 1966, y Madrid, 3 de febrero de 2011), Reino Unido (Londres, 1975, y 2013) y Estados Unidos (1990, y protocolo de 14 de enero de 2013).
Ministerio de Hacienda: notas de prensa de 14 de enero de 2013 y 13 de diciembre de 2021 (ratificación del Instrumento Multilateral).
Directiva 2011/96/UE del Consejo (matriz-filial); Directiva 2003/49/CE del Consejo (intereses y cánones); Convenio 90/436/CEE (arbitraje en precios de transferencia); Directiva (UE) 2017/1852 del Consejo (resolución de litigios fiscales).
Ley 27/2014, del Impuesto sobre Sociedades, arts. 21, 22, 31, 32 y 107-108 (ETVE); Ley 35/2006, del IRPF, art. 80.
Ley 11/2020, de 30 de diciembre, de Presupuestos Generales del Estado para 2021 (exención del 95% del art. 21 LIS).
Real Decreto 1080/1991, de 5 de julio (lista original de paraísos fiscales).
Ley 36/2006, de 29 de noviembre, de medidas para la prevención del fraude fiscal, disposición adicional primera (en la redacción de la Ley 11/2021).
Orden HFP/115/2023, de 9 de febrero (jurisdicciones no cooperativas), y proyecto de orden de modificación sometido a información pública el 22 de mayo de 2026.
OCDE: Model Tax Convention on Income and on Capital, arts. 23 A (método de exención, incluido el apartado 4), 23 B (método de imputación) y 29 (derecho a los beneficios), y sus Comentarios.
OCDE (1998): Tax Sparing: A Reconsideration, OECD Publishing, París.
Tema de imposición directa (teoría y comparaciones internacionales), bloque sobre el sujeto: neutralidad en la exportación de capital, en la importación de capital y en la propiedad del capital (Desai y Hines, 2003).
Convenio entre el Imperio austrohúngaro y Prusia para evitar la doble imposición, de 21 de junio de 1899.
BRUINS, G. W. J., EINAUDI, L., SELIGMAN, E. R. A. y STAMP, J. (1923): Report on Double Taxation Submitted to the Financial Committee, Sociedad de Naciones, Doc. E.F.S.73.F.19.
Sociedad de Naciones (1927): Double Taxation and Tax Evasion: Report of the Committee of Technical Experts, Doc. C.216.M.85.1927.II.
Modelos de convenio de la Sociedad de Naciones de México (1943) y Londres (1946).
JOGARAJAN, S. (2018): Double Taxation and the League of Nations, Cambridge University Press.
BAKER, P. (2023): "The history of international tax treaties from the UK perspective", Milano University Press.
Convención de Viena sobre el Derecho de los Tratados, de 23 de mayo de 1969.
Constitución Española, art. 96; texto refundido de la Ley del IRNR, art. 4.
OCDE: Model Tax Convention on Income and on Capital, versiones de 1963, 1977, 1992, 2017 y actualización de 2025 (publicada el 19 de noviembre de 2025).
Naciones Unidas: Model Double Taxation Convention between Developed and Developing Countries (1980 y actualizaciones posteriores; art. 12A de 2017 y art. 12B de 2021).
Naciones Unidas, Asamblea General: Resolución 78/230, de 22 de diciembre de 2023, "Promoción de una cooperación internacional en cuestiones de tributación inclusiva y eficaz en las Naciones Unidas".
Naciones Unidas, Comité Especial: términos de referencia de la Convención Marco de las Naciones Unidas sobre Cooperación Fiscal Internacional, agosto de 2024.
Naciones Unidas, Asamblea General: Resolución 79/235, de diciembre de 2024 (creación del comité intergubernamental de negociación).
Comité intergubernamental de negociación: sesión organizativa (Nueva York, 3-6 de febrero de 2025), tercera sesión (Nairobi, noviembre de 2025) y quinta sesión (Nueva York, 3-13 de agosto de 2026).
Naciones Unidas: Model Double Taxation Convention between Developed and Developing Countries, actualización de 2021 (artículo 12B, servicios digitales automatizados).
Compromiso de Sevilla, Cuarta Conferencia Internacional sobre la Financiación para el Desarrollo, 2025.
Tratado de Funcionamiento de la Unión Europea, arts. 107 (ayudas de Estado) y 115 (aproximación de legislaciones, unanimidad).
Sentencia del Tribunal de Justicia de 8 de noviembre de 2022, Fiat Chrysler Finance Europe/Comisión, asuntos acumulados C-885/19 P y C-898/19 P.
Sentencia del Tribunal de Justicia de 14 de diciembre de 2023, Comisión/Amazon.com y otros, asunto C-457/21 P.
Directiva 2011/16/UE del Consejo, de 15 de febrero de 2011, de cooperación administrativa, y sus sucesivas modificaciones (DAC2 a DAC9).
Directiva (UE) 2018/822 del Consejo (DAC6), Directiva (UE) 2021/514 del Consejo (DAC7), Directiva (UE) 2023/2226 del Consejo (DAC8) y Directiva (UE) 2025/872 del Consejo (DAC9).
Directiva (UE) 2021/2101 del Parlamento Europeo y del Consejo, de 24 de noviembre de 2021 (informe país por país público).
Comisión Europea: propuestas de Directiva COM(2021) 565 (Unshell), COM(2022) 216 (DEBRA), COM(2023) 529 (precios de transferencia), COM(2023) 532 (BEFIT), y COM(2011) 121 y COM(2016) 683/685 (BICCIS).
Comisión Europea: Programa de Trabajo para 2026, 21 de octubre de 2025 (retirada de Unshell, DEBRA, Directiva de precios de transferencia e impuesto sobre transacciones financieras).
Consejo de la UE: informe ECOFIN 9960/25, de 18 de junio de 2025.
Directiva (UE) 2022/2523 del Consejo, de 15 de diciembre de 2022 (Pilar 2).
Consejo de la UE: lista de jurisdicciones no cooperativas a efectos fiscales (desde diciembre de 2017, revisiones semestrales).
OCDE/G20 Marco Inclusivo (2021): Statement on a Two-Pillar Solution to Address the Tax Challenges Arising from the Digitalisation of the Economy, 8 de octubre de 2021.
OCDE (2021): Tax Challenges Arising from the Digitalisation of the Economy — Global Anti-Base Erosion Model Rules (Pillar Two), diciembre de 2021.
OCDE (2023): Multilateral Convention to Implement Amount A of Pillar One (texto de octubre de 2023, no firmado).
OCDE (2024): Pillar One — Amount B, 19 de febrero de 2024.
OCDE/G20 Marco Inclusivo (2026): Side-by-Side Package, administrative guidance, 5 de enero de 2026.
Declaración del G7 sobre el Pilar 2, junio de 2025.
Memorándum presidencial de EEUU sobre el acuerdo fiscal global de la OCDE, 20 de enero de 2025.
USTR: investigaciones de la Sección 301 sobre impuestos a los servicios digitales (2019-2020) y declaración conjunta de 21 de octubre de 2021 con Austria, Francia, Italia, España y Reino Unido.
Ley 4/2020, de 15 de octubre, del Impuesto sobre Determinados Servicios Digitales.
Ley 22/2021, de 28 de diciembre, de Presupuestos Generales del Estado para 2022 (tributación mínima en el IS).
Directiva (UE) 2022/2523 del Consejo, de 15 de diciembre de 2022.
Ley 7/2024, de 20 de diciembre, por la que se establecen un Impuesto Complementario, un Impuesto sobre el margen de intereses y comisiones de determinadas entidades financieras y un Impuesto sobre los líquidos para cigarrillos electrónicos.
Real Decreto 252/2025, de 1 de abril, Reglamento del Impuesto Complementario.
ZODROW, G. R. y MIESZKOWSKI, P. (1986): "Pigou, Tiebout, Property Taxation, and the Underprovision of Local Public Goods", Journal of Urban Economics, vol. 19, núm. 3.
WILSON, J. D. (1986): "A Theory of Interregional Tax Competition", Journal of Urban Economics, vol. 19, núm. 3.
KANBUR, R. y KEEN, M. (1993): "Jeux Sans Frontières: Tax Competition and Tax Coordination When Countries Differ in Size", American Economic Review, vol. 83, núm. 4.
TIEBOUT, C. M. (1956): "A Pure Theory of Local Expenditures", Journal of Political Economy, vol. 64, núm. 5.
BRENNAN, G. y BUCHANAN, J. M. (1980): The Power to Tax, Cambridge University Press (ya citado en el Tema 12).
SLEMROD, J.: obra sobre costes sociales del cumplimiento y la elusión tributaria.
OCDE (2015): Measuring and Monitoring BEPS, Action 11 — 2015 Final Report.
TØRSLØV, T., WIER, L. y ZUCMAN, G. (2023): "The Missing Profits of Nations", Review of Economic Studies, vol. 90, núm. 3, pp. 1499-1534.
BLOUIN, J. y ROBINSON, L. (2023): "Accounting for the Profits of Multinational Enterprises: Double Counting and Misattribution of Foreign Affiliate Income", SSRN 3491451 (versión inicial de 2019-2020 como "Double Counting Accounting").
CLAUSING, K. (2016 y 2020): estimaciones de la pérdida recaudatoria de EEUU por traslado de beneficios.
Sentencia del Tribunal de Justicia de 10 de septiembre de 2024, Comisión/Irlanda y otros, asunto C-465/20 P (caso Apple).
- OCDE (2013): Action Plan on Base Erosion and Profit Shifting, OECD Publishing, París.
- OCDE (2015): BEPS 2015 Final Reports (informes finales de las 15 acciones), 5 de octubre de 2015.
- OCDE/G20: Inclusive Framework on BEPS — Progress Report September 2022-September 2023.
- OCDE (2025): Revised BEPS Action 5 Transparency Framework on Tax Rulings, 8 de septiembre de 2025.
- Convenio Multilateral para aplicar las medidas relacionadas con los tratados fiscales para prevenir la erosión de las bases imponibles y el traslado de beneficios (MLI), adoptado el 24 de noviembre de 2016.
- Directiva (UE) 2016/1164 del Consejo, de 12 de julio de 2016 (ATAD 1).
- Directiva (UE) 2017/952 del Consejo, de 29 de mayo de 2017 (ATAD 2).
- Ley 27/2014, de 27 de noviembre, del Impuesto sobre Sociedades, arts. 15 bis, 16, 19 y 23.
- Real Decreto-ley 4/2021, de 9 de marzo (asimetrías híbridas), y Ley 5/2022, de 9 de marzo.
- Ley 11/2021, de 9 de julio, de medidas de prevención y lucha contra el fraude fiscal.
- Real Decreto-ley 18/2022 (asimetrías híbridas invertidas).
- SANZ GADEA, E. (2021): "Asimetrías híbridas: Real Decreto-Ley 4/2021", Revista de Contabilidad y Tributación. CEF.


\MyBibItem{DAWID y GATTI}{Chapter 2 - Agent-Based Macroeconomics}{2018}{https://doi.org/10.1016/bs.hescom.2018.02.006}


% ANEXOS
\clearpage










\end{document}